\documentclass[11pt]{amsart}

\usepackage[a4paper,margin=2.5cm]{geometry}
\usepackage{amsmath,amssymb,mathtools}
\usepackage{enumitem}
\usepackage[hidelinks]{hyperref}

\newtheorem{theorem}{Theorem}[section]
\newtheorem{proposition}[theorem]{Proposition}
\newtheorem{lemma}[theorem]{Lemma}
\newtheorem{corollary}[theorem]{Corollary}
\newtheorem{remark}[theorem]{Remark}
\newtheorem{definition}[theorem]{Definition}

\newcommand{\R}{\mathbb R}
\newcommand{\N}{\mathbb N}
\newcommand{\Om}{\Omega}
\newcommand{\Et}{E_{\hat t}}
\newcommand{\dT}{\delta_T}
\newcommand{\AF}{\mathcal A_F}
\newcommand{\HH}{\mathcal H}
\newcommand{\eps}{\varepsilon}
\newcommand{\Etil}{\widetilde E}
\newcommand{\Gtil}{\widetilde G}
\newcommand{\Sclass}{\mathcal S}
\newcommand{\wt}{\widetilde}
\DeclareMathOperator{\dist}{dist}
\DeclareMathOperator{\Lip}{Lip}
\newcommand{\dHnm}{d\mathcal H^{n-1}\!\restriction_{S^{n-1}}}

\title{Part 4: The Selection--Level--Brandolini Hybrid Route
       (rigorous version, dimension-free main chain)}
\author{}
\date{May 2026}

\begin{document}
\maketitle

\section{Purpose and status}

This note records the hybrid Selection--Level--Brandolini route in
one place, with every step proved.  The main chain is
\emph{dimension-free}: each building block holds for arbitrary
$n\ge 2$, and the route delivers
$\AF(U)^{2}\le C\,\dT(U)$ with the sharp exponent in every
dimension.  The explicit constants are written out for $n=2$ (where
they are most transparent); the general-$n$ specialisation is
catalogued in \S\ref{ss:general-n}.  The proof combines:

\begin{enumerate}[label=\arabic*.]
\item the quantitative selection reduction and its $C^{1,\gamma}$
regularity output (Alt--Caffarelli, no Schauder bootstrap);
\item a level extraction at the $\sqrt{\dT}$ boundary-layer scale;
\item a quantitative $D_H\to L^\infty$-Serrin upgrade on the stopped
level (using $\gamma$-Hölder interpolation);
\item inheritance of the selected $C^{1,\gamma}$ regularity by the
inset;
\item the Brandolini--Nitsch--Salani--Trombetti Serrin stability
theorem in its $C^{1,\gamma}$ form (proved here as
Theorem~\ref{thm:BNST-Cgamma}), giving a concentric annular core;
\item a dimension-free \emph{radial-graph promotion}
(Lemma~\ref{lem:radialgraph}): the $C^{1,\gamma}$ regularity and the
Brandolini inclusion together force $\partial\Etil$ to be a
$C^{1,\gamma}$ radial graph from a barycentre, with no pocket-filling
needed;
\item a barycentre-recentring step;
\item the direct small-amplitude radial-graph stability theorem.
\end{enumerate}

Step~6 is the load-bearing dimension-free part of the route: in any
dimension $n\ge 2$, the chain delivers the sharp exponent
$\AF(U)^2\lesssim\dT(U)$ once a dimension-free radial-graph theorem
$(G_n)$ is supplied (the planar case is proved in
\S\ref{sec:graph-proof}; the general-$n$ dimension-free derivation is
cited in \S\ref{ss:G12-conclude}).  For $n=2$ we additionally record,
in \S\ref{sec:hole-fill}, the elementary polar-slicing
``hole-filling'' alternative to Step~6: a planar radial-hull pocket
lemma that gives an $O(\dT)$-volume radialisation by the sharp
planar isoperimetric inequality $|Q|\le P(Q)^2/(4\pi)$ followed by
super-additivity.  The planar alternative is cute and self-contained;
it does not extend to $n\ge 3$ at the sharp scaling because in higher
dimensions the analogous quadratic isoperimetric bound is replaced
by a $P(Q)^{n/(n-1)}$ bound, which yields only
$v=O(\dT^{n/(2(n-1))})\not=O(\dT)$.  The radial-graph route of
Step~6 sidesteps this obstruction entirely.

The route does \emph{not} use Brasco--De Philippis--Velichkov's
quantitative stability theorem as a black box.  It does use, with
explicit statements:

\begin{enumerate}[label=(I\arabic*)]
\item the planar volume-variable level identity (Plan~2/Plan~3);
\item the $C^{1,\gamma}$ Brandolini--Nitsch--Salani--Trombetti
Serrin stability theorem;
\item the direct radial-graph theorem proved in the Plan~3 graph note.
\end{enumerate}

The selected-minimizer regularity package $\Sclass$ and the
quantitative selection principle (P) --- previously imported as
Plan~1 black boxes --- are now \emph{proved in full} in
\S\ref{sec:selection-package} below.  The only deep ingredient
required for that proof is the $C^{1,\gamma}$ Alt--Caffarelli
free-boundary regularity for selected quasi-minimisers; everything
in $(S1)$--$(S6)$ beyond $(S3)$ is then derived from Lieberman
boundary-Schauder estimates for the torsion equation on
$C^{1,\gamma}$ domains and the Hopf lemma, in a quantitative form.
\emph{No further Schauder bootstrap to $C^{2,\alpha}$ is required},
because the BNST theorem (B) is itself proved here at the
$C^{1,\gamma}$ regularity level (\S\ref{sec:BNST-proof}).

Everything else is proved below in full.  In particular, all asymmetry
quantities are kept in their Fraenkel form \(\AF\), so that the final
estimate is sharp in the exponent without any auxiliary smooth-asymmetry
detour.

\section{Conventions, the selection class, and imported inputs}

\subsection{Torsion convention}

Throughout this note, \emph{the torsion function} \(u_D\) of a bounded
open set \(D\subset\R^2\) is the unique \(H^1_0(D)\) solution of
\begin{equation}\label{eq:torsion-pde}
   -\Delta u_D=1\quad\hbox{in }D,
   \qquad
   u_D=0\quad\hbox{on }\partial D.
\end{equation}
Equivalently, \(\Delta u_D=-1\) and \(u_D\ge 0\) in \(D\).  We set
\[
   T(D):=\int_D u_D\,dx.
\]
The (unnormalised) Saint--Venant deficit of \(D\) at fixed area is
\[
   \dT(D):=T(B_D)-T(D),
\]
where \(B_D\) is any disk with \(|B_D|=|D|\).  Under the convention
\eqref{eq:torsion-pde}, the explicit torsion of a planar disk is
\begin{equation}\label{eq:disk-torsion}
   T(B_r)=\frac{\pi r^4}{8},
   \qquad
   T(B_A)=\frac{A^2}{8\pi}\quad(A=|B_A|).
\end{equation}
On the disk \(B_1\) of radius \(1\) (area \(\pi\)) the gradient satisfies
\(|\nabla u_{B_1}|(x)=|x|/2\), and in particular
\(|\nabla u_{B_1}|\equiv 1/2\) on \(\partial B_1\).

We normalize the comparison ball to area \(\pi\), so the area-\(\pi\)
disk has radius \(1\).  Constants change harmlessly under any other
normalisation by an explicit power of the dilation factor.

The \emph{Fraenkel asymmetry} of a bounded measurable set
\(E\subset\R^2\) with \(|E|>0\) is
\[
   \AF(E):=\inf_{z\in\R^2}\frac{|E\Delta B(z)|}{|E|},
\]
where \(B(z)\) is the disk of area \(|E|\) centred at \(z\).  This is
the asymmetry we control: the corollary in \S10 is stated for \(\AF\).

\subsection{The selection class}

Selected minimisers of the BDV variational problem live in a class
with explicit uniform regularity.  Fix once and for all the
structural constants
\[
   R_0,\ K_{1,\gamma},\ \rho_0,\ \Lambda_\gamma,\ q_-,\ q_+,\ r_0,\
   c_0,\ C_0>0,
   \qquad
   0<\gamma<1.
\]

\begin{definition}[Selection class \(\Sclass\)]\label{def:Sclass}
The \emph{selection class} \(\Sclass=\Sclass(R_0,K_{1,\gamma},\rho_0,
\Lambda_\gamma,q_\pm,r_0,c_0,C_0,\gamma)\) is the family of all open
sets \(U\subset\R^2\) satisfying:
\begin{enumerate}[label=(S\arabic*)]
\item \(U\) is bounded, connected, and simply connected;\label{S1}
\item \(|U|=\pi\) and \(U\subset B_{R_0}\);\label{S2}
\item \(\partial U\) is a closed \(C^{1,\gamma}\) Jordan curve with
$C^{1,\gamma}$ atlas pair $(K_{1,\gamma},\rho_0)$ \textup{(}i.e.\
every \(x_0\in\partial U\) admits a chart
\(\psi:B_{\rho_0}(x_0)\to\R^2\) with
\(\|\psi\|_{C^{1,\gamma}},\|\psi^{-1}\|_{C^{1,\gamma}}\le K_{1,\gamma}\)
straightening \(\partial U\) to \(\partial\R^2_+\)\textup{)};\label{S3}
\item the torsion function $u=u_U$ satisfies
$\|u\|_{C^{1,\gamma}(\overline U)}\le M_{1,\gamma}$ and the
gradient Hölder bound $\|Du\|_{C^{0,\gamma}(\overline U)}\le
\Lambda_\gamma$;\label{S4}
\item the Hopf bound \(q_-\le|\nabla u|\le q_+\) holds on
\(\partial U\);\label{S5}
\item the linear-growth pair holds in the collar:
\[
   c_0\,\mathrm{dist}(x,\partial U)\le u(x)\le C_0\,
   \mathrm{dist}(x,\partial U)
   \quad\hbox{for }x\in\mathcal N_{r_0}(\partial U)\cap U.
\]\label{S6}
\end{enumerate}
Here $M_{1,\gamma}=M_{1,\gamma}(R_0,K_{1,\gamma},\rho_0)$ is the
universal constant from boundary Lieberman--Schauder estimates for
$-\Delta u=1$ on $C^{1,\gamma}$ domains, depending only on
$(R_0,K_{1,\gamma},\rho_0)$; we absorb it into the structural
constants.
\end{definition}

The constants in \(\Sclass\) are fixed throughout; every constant
generated below is allowed to depend on them.  We write
\[
   C\ \hbox{or}\ C(\Sclass)
\]
for a constant depending only on \(\Sclass\).

\begin{remark}\label{rem:S-source}
Properties (S1)--(S6) constitute the regularity package that the
Alt--Caffarelli step of the BDV selection attaches to a selected
minimiser; crucially, only the $C^{1,\gamma}$ free-boundary regularity
(R4) of Theorem~\ref{thm:BDV-reg} is required to verify (S3), and (S4)
follows from boundary Lieberman--Schauder estimates for $-\Delta u=1$
on $C^{1,\gamma}$ domains.  No further $C^{2,\alpha}$ bootstrap is
needed, since the BNST input (B) of \S\ref{sec:BNST-import} is the
$C^{1,\gamma}$ version (proved in \S\ref{sec:BNST-proof}).  We use
$\Sclass$ as a black box throughout \S\S3--8.  \S9 imports the
quantitative selection principle (P) that delivers $U\in\Sclass$ from
a small-asymmetry near-counterexample.  The two BDV-conditional
inputs of this note are thus (i) the $C^{1,\gamma}$ regularity class
$\Sclass$ (Definition~\ref{def:Sclass}), derived in \S11 from the
Alt--Caffarelli regularity package (R1)--(R4), and (ii) the selection
principle (P) of \S\ref{sec:transfer}.
\end{remark}

\subsection{The Brandolini--Nitsch--Salani--Trombetti theorem
under \texorpdfstring{$C^{1,\gamma}$}{C(1,gamma)} regularity}
\label{sec:BNST-import}

We use the following $C^{1,\gamma}$ version of the planar
Brandolini--Nitsch--Salani--Trombetti theorem~\cite{BNST}.  The
original BNST result is stated for $C^{2,\alpha}$ domains; we
strengthen it to $C^{1,\gamma}$ in \S\ref{sec:BNST-proof}, where it
is proved in full as Theorem~\ref{thm:BNST-Cgamma}.  This
strengthening allows the manuscript to bypass any Schauder bootstrap
of the boundary: the $C^{1,\gamma}$ regularity required below is
delivered directly by the Alt--Caffarelli step (R4) of the BDV
regularity package (Theorem~\ref{thm:BDV-reg}).

\begin{quote}
\textbf{(B)} Let $\gamma\in(0,1]$ and let the \emph{atlas pair}
$(K_{1,\gamma},\rho_0)$ consist of the $C^{1,\gamma}$ chart norm
$K_{1,\gamma}$ and a uniform chart radius $\rho_0>0$.  There exist
$\eta_B\in(0,1/4)$ and constants $C_{\rm Br}>0$ and $\mu=1/34$
\textup{(the planar value of $\mu=1/(2(4n+9)(n-1))$ at $n=2$)},
depending polynomially on $(K_{1,\gamma},1/\rho_0,R_0,\Lambda_\gamma)$
and on $1/\gamma$, such that the following holds.  If $D\subset\R^2$
is bounded, connected, with diameter $\le R_0$ and $C^{1,\gamma}$
atlas pair bounded by $(K_{1,\gamma},\rho_0)$, and if
$w\in C^{1,\gamma}(\overline D)\cap C^\infty(D)$ solves
\[
   \Delta w=2\ \hbox{in }D,
   \qquad
   w=0\ \hbox{on }\partial D,
   \qquad
   \bigl\||Dw|-1\bigr\|_{L^\infty(\partial D)}\le \eta\le\eta_B,
\]
with $\|Dw\|_{C^{0,\gamma}(\overline D)}\le\Lambda_\gamma$, then
there exists $z\in\R^2$ and $r_{\rm in},r_{\rm out}>0$ with
\begin{equation}\label{eq:brandolini-gap}
   B_{r_{\rm in}}(z)\subset D\subset B_{r_{\rm out}}(z),
   \qquad
   r_{\rm out}-r_{\rm in}\le C_{\rm Br}\,\eta^\mu,
   \qquad
   \bigl|r_{\rm in}-\rho_D\bigr|+\bigl|r_{\rm out}-\rho_D\bigr|
   \le C_{\rm Br}\,\eta^\mu,
\end{equation}
where $\rho_D:=\sqrt{|D|/\pi}$ is the radius of the disk of area
$|D|$.
\end{quote}

Two observations: \textbf{(i)} No \emph{smallness} of the
$C^{1,\gamma}$ atlas norm is required, only finiteness;
$K_{1,\gamma}$ and $1/\rho_0$ enter polynomially in $C_{\rm Br}$,
and the only smallness parameter is $\eta$.  \textbf{(ii)} The
exponent $\mu=1/34$ is unchanged from the $C^{2,\alpha}$ version:
the rate comes from the $L^1$ multi-ball argument of BNST (their
Theorem~5; see \S\ref{ss:BNST-anatomy} step \ref{B1}), which uses
no boundary regularity above Lipschitz.  Only the constant
$C_{\rm Br}$ is affected by the $C^{2,\alpha}\to C^{1,\gamma}$
relaxation.

\subsection{Imported planar level identity}

We use the volume-variable form of the level-set deficit identity
(Plan~2/Plan~3, planar specialisation):

\begin{quote}
\textbf{(L)} For every bounded open \(U\subset\R^2\) of finite
perimeter, with torsion function \(u=u_U\) and inverse distribution
function \(v(s):=\inf\{t>0:|\{u>t\}|\le s\}\), and with
\[
   E_t:=\{u>t\},\quad
   f_t:=|\nabla u|\restriction_{\partial E_t},\quad
   P_t:=\HH^1(\partial E_t),\quad
   A_t:=|E_t|,
\]
the level deficits
\begin{equation}\label{eq:DH-DI-def}
   D_H(t):=
      \Bigl(\int_{\partial E_t}f_t\,d\HH^1\Bigr)
      \Bigl(\int_{\partial E_t}\frac1{f_t}\,d\HH^1\Bigr)
      -P_t^2,
   \qquad
   D_I(t):=P_t^2-4\pi A_t,
\end{equation}
are non-negative and satisfy
\begin{equation}\label{eq:level-identity}
   \int_0^{|U|}\bigl(D_H(v(s))+D_I(v(s))\bigr)\,ds
      =8\pi\,\dT(U).
\end{equation}
\end{quote}

The constant \(8\pi\) is the planar specialisation of the general
identity, which has the dimensional factor
\(n^2\omega_n^{2/n}\,s^{1-2/n}\); in \(n=2\) this equals \(4\pi\cdot 1\)
and the equality with the deficit is obtained after multiplication by
\(2\) and by \(|U|\) (cf.\ Plan~2, agent1).

\subsection{Imported radial-graph theorem}\label{ss:G-import}

We state the radial-graph theorem in dimension-free form
$(G_n)$, valid for every $n\ge 2$; the planar case ($n=2$) is the
form used in the main chain of \S\ref{sec:selected-domain}, and
is proved in full as Theorem~\ref{thm:G12-explicit} in
\S\ref{sec:graph-proof}.  The dimension-$n$ version is proved as
Theorem~\ref{thm:Gn} in \S\ref{ss:Gn}.

\begin{quote}
\textbf{($G_n$)} Let $n\ge 2$.  There exist
$\eta_G^{(n)},c_G^{(n)},C_G^{(n)}>0$ such that the following holds.
If $\varphi\in C^0(S^{n-1})$ and
\[
   G=\{z+r\omega:0\le r<1+\varphi(\omega),\ \omega\in S^{n-1}\},
\]
satisfies
\[
   \|\varphi\|_{L^\infty(S^{n-1})}\le\eta_G^{(n)},
   \qquad
   |G|=\omega_n,
   \qquad
   \int_G(x-z)\,dx=0
\]
\textup{(i.e.\ the barycentre of $G$ is at the radial centre $z$)},
then
\begin{equation}\label{eq:graph-theorem}
   \AF(G)^{2}\le C_G^{(n)}\,\dT(G).
\end{equation}
At $n=2$ this is the planar (G) used in
\S\ref{sec:selected-domain}; we then drop the superscript and write
$\eta_G,c_G,C_G$.  Admissible explicit constants are
$\eta_G=1/13,c_G=15/1024,C_G<48$ (planar), or
$c_G^{(n)}\ge 1/(16 n^{2})$ (general).
\end{quote}

Both the smallness of $\|\varphi\|_{L^\infty}$ and the barycentre
condition are needed in the proof of $(G_n)$, the latter to suppress
the $P_1$-mode in the spherical-harmonic expansion of $\varphi$.  We
carry out the barycentre recentring in
Lemma~\ref{lem:recenter} below.

\subsection{General-\texorpdfstring{$n$}{n} specialisation}
\label{ss:general-n}

The detailed write-up of \S\S\ref{sec:level-id}--\ref{sec:transfer}
is carried out at $n=2$, where every constant is explicit.  The same
chain extends verbatim to general $n\ge 2$ via the following
substitutions:

\begin{itemize}[leftmargin=2em]
\item Ambient: $\R^2\to\R^n$, $S^1\to S^{n-1}$,
$\HH^1\to\HH^{n-1}$, Jordan curve\,$\to$\,closed embedded
$C^{1,\gamma}$ hypersurface.
\item Normalisations: $|U|=\pi\to|U|=\omega_n$ (volume of the unit
ball); area of optimal disk $\pi\to\omega_n$; torsion of the unit
ball $T(B_1)=\pi/8\to T(B_1)=\omega_n/(n(n+2))$; gradient of the
torsion function on $\partial B_1$ is $1/n$ (so $|Dw|=1$ on
$\partial B_1$ corresponds to $w=-2 u_0$ with $u_0=(1-|x|^2)/(2n)$).
\item Level identity \eqref{eq:level-identity}: the constant
$8\pi=4\pi\cdot 2$ becomes
$C_{\mathrm{lev}}(n):=n^{2}\omega_n^{2/n}\,\bigl(|U|^{2-2/n}\,
n/(2(n-1))\bigr)$ — concretely, the integral identity
$\int_0^{|U|}(D_H(v)+D_I(v))ds=C_{\mathrm{lev}}(n)\,\dT(U)$ with
$C_{\mathrm{lev}}(n)$ as above and $D_I(t)=P_t^2-n^{2}
\omega_n^{2/n} A_t^{2-2/n}$ (the $n$-D isoperimetric deficit);
cf.\ Plan~2/Plan~3.  At $n=2$, $C_{\mathrm{lev}}(2)=8\pi$ and
$n^{2}\omega_n^{2/n}=4\pi$ recover the planar form.
\item BNST exponent: $\mu=1/34\to\mu=1/(2(4n+9)(n-1))$ (see
Theorem~\ref{thm:BNST-Cgamma} for the $C^{1,\gamma}$ form, which is
stated for $\R^2$ but whose proof — \S\ref{sec:BNST-proof} — uses
only Lipschitz boundary in steps (B1) and (B3), and the
Hölder-modulus replacement of $D^2 w$ in (B2) is dimension-free).
\item $D_H\to L^\infty$ Serrin interpolation (Lemma~\ref{lem:interp}):
exponent $\gamma/(1+2\gamma)\to\gamma/(2\gamma+n-1)$ (the
denominator counts the surface dimension, with $n-1$ in place of
$1$), with constant depending on the lower density of $\partial\Et$
and the chart radius (both controlled by $\Sclass$).
\item Radial-graph promotion (Lemma~\ref{lem:radialgraph}): already
stated and proved for arbitrary $n\ge 2$.
\item Radial-graph theorem: planar (G) (Theorem~\ref{thm:G12-explicit})
is replaced by $(G_n)$ (Theorem~\ref{thm:Gn} of \S\ref{ss:Gn}),
with $c_G^{(n)}=1/(16 n^{2})$ and $C_G^{(n)}=8 n^{2}\bigl(C_n/
\omega_n\bigr)^{2}$.
\item Hole-filling (\S\ref{sec:hole-fill}): \emph{not} available in
general $n$, since the planar isoperimetric bound $|Q|\le P(Q)^{2}/
(4\pi)$ is replaced by $|Q|\le C_n P(Q)^{n/(n-1)}$, which yields
only $v=O(\delta^{n/(2(n-1))})$, strictly weaker than $O(\delta)$
for $n\ge 3$.  The radial-graph route of \S\ref{sec:radialgraph}
is the dimension-free replacement.
\end{itemize}

With these substitutions, the selected-domain bound
$\AF(U)^{2}\le C_\star(\Sclass,n)\,\dT(U)$ of
Theorem~\ref{thm:selected} holds for every $n\ge 2$ with the same
sharp exponent.  In the rest of this article we write all explicit
constants for $n=2$ (where they are most transparent); the reader
substitutes the entries above for general $n$.

\begin{remark}\label{rem:Plan1}
The regularity class $\Sclass$ (Definition~\ref{def:Sclass}) and the
quantitative selection principle (P) of \S\ref{sec:transfer} are
\emph{proved in \S\ref{sec:selection-package}} (see in particular
Theorem~\ref{thm:S-package} for (S) and Theorem~\ref{thm:P-package}
for (P)); the rest of \S\S3--10 uses them as already established
hypotheses.  The only ``deep'' input required for the proof of (S)
is $C^{1,\gamma}$ Alt--Caffarelli free-boundary regularity of a
selected quasi-minimiser (Brasco--De~Philippis--Velichkov 2015,
Theorem 4.18), invoked as a single black box and recorded as (R4) in
Theorem~\ref{thm:BDV-reg}.  $(S1),(S2),(S4),(S5),(S6)$ are then
derived from $(S3)=$(R4) by boundary Lieberman--Schauder estimates
for the torsion equation on $C^{1,\gamma}$ domains and the Hopf
lemma, in a quantitative form, with no smallness hypothesis on the
$C^{1,\gamma}$ atlas norm.
\end{remark}

\section{Level-set identities and the stopped level}\label{sec:level-id}

\begin{lemma}[Good level in a prescribed layer]\label{lem:good-level}
Let \(U\in\Sclass\), set \(\delta:=\dT(U)\), and fix
\(0<\lambda\le|U|/2=\pi/2\).  There is a regular value \(\hat t>0\)
of \(u\) such that
\begin{equation}\label{eq:layer-size}
   \frac{\lambda}2\le |U\setminus E_{\hat t}|\le \lambda,
\end{equation}
and
\begin{equation}\label{eq:good-D}
   D_H(\hat t)+D_I(\hat t)
      \le \frac{16\pi}{\lambda}\,\delta.
\end{equation}
\end{lemma}

\begin{proof}
In the volume variable \(s=|E_t|\), the layer
\(|U|-\lambda\le s\le|U|-\lambda/2\) has length \(\lambda/2\).  By
\eqref{eq:level-identity} and non-negativity of \(D_H,D_I\),
\[
   \int_{|U|-\lambda}^{|U|-\lambda/2}
      \bigl(D_H(v(s))+D_I(v(s))\bigr)\,ds
   \le 8\pi\delta,
\]
so for some \(s\) in this layer
\(D_H(v(s))+D_I(v(s))\le (8\pi\delta)/(\lambda/2)=16\pi\delta/\lambda\).
Sard's theorem applied to \(u\) gives that regular values of \(u\) are
full-measure in the volume variable (the singular set of values is
mapped under \(t\mapsto |E_t|\) to a Lebesgue-null subset of
\([0,|U|]\)).  We may therefore choose \(s\) so that \(\hat t:=v(s)\) is
a regular value, which gives \eqref{eq:layer-size} and
\eqref{eq:good-D}.
\end{proof}

\begin{lemma}[Boundary-layer height]\label{lem:height}
Let \(U\in\Sclass\).  There are constants \(c_h,C_h>0\) depending on
\(\Sclass\) such that, for every regular value \(t\) with
\(0<t<r_0\,c_0/2\),
\begin{equation}\label{eq:t-layer}
   c_h\,t\le |U\setminus E_t|\le C_h\,t.
\end{equation}
Consequently, when \(\hat t\) is the level from
Lemma~\ref{lem:good-level} with \(\lambda<\lambda_0:=
\min(\pi/2,r_0 c_0 C_h/2)\),
\begin{equation}\label{eq:t-lambda}
   \frac{\lambda}{2C_h}\le \hat t\le \frac{\lambda}{c_h}.
\end{equation}
\end{lemma}

\begin{proof}
By (S6), in the collar of width $r_0$,
\[
   \{\mathrm{dist}(\cdot,\partial U)\le t/C_0\}\cap U
      \subset
   \{0<u\le t\}
      \subset
   \{\mathrm{dist}(\cdot,\partial U)\le t/c_0\}\cap U.
\]
The chosen range $t<r_0c_0/2$ makes $t/c_0<r_0/2$, so both
inclusions lie strictly inside the collar.  Since $\partial U\in
C^{1,\gamma}$ with atlas pair $(K_{1,\gamma},\rho_0)$ (S3), the
normal vector $\nu_U:\partial U\to S^1$ is $\gamma$-Hölder.  By the
standard tubular-neighbourhood formula for $C^{1,\gamma}$
hypersurfaces \textup{(}cf.\ \cite[Theorem~3.1]{HP}\textup{)}, normal
coordinates $\Phi:(y,s)\mapsto y-s\nu_U(y)$ define a bi-Lipschitz
parametrisation of the collar
$\mathcal N_{r_0'}(\partial U)\cap U$ for some
$r_0'=r_0'(K_{1,\gamma},\rho_0)>0$, with Jacobian satisfying
$|J(y,s)-1|\le c(K_{1,\gamma})\,s^\gamma$ (a $\gamma$-Hölder
modulus of $\nu_U$).  In particular
\[
   1-c(K_{1,\gamma})\,r^\gamma\le J(y,s)\le 1+c(K_{1,\gamma})\,r^\gamma
   \qquad\text{for }s\in[0,r],
\]
giving
\[
   |\{\mathrm{dist}(\cdot,\partial U)\le r\}\cap U|
      =\int_{\partial U}\int_0^r J(y,s)\,ds\,d\HH^1(y)
      =r\,\HH^1(\partial U)+O(r^{1+\gamma}),
\]
with constants in $O(\cdot)$ controlled by $K_{1,\gamma}$.  Since
$U\in\Sclass$ is bounded, $|U|=\pi$ and the boundary has uniform
$C^{1,\gamma}$ atlas, the perimeter satisfies $P(U)\in[P_-,P_+]$ with
$P_\pm$ depending only on $\Sclass$ (lower bound from the planar
isoperimetric inequality $P\ge 2\pi$; upper bound from the
$C^{1,\gamma}$ chart covering of $\partial U\subset B_{R_0}$).
Substituting $r=t/c_0$ and $r=t/C_0$ yields \eqref{eq:t-layer} with
$c_h=P_-/(2C_0)$ and $C_h=2P_+/c_0$ (the higher-order term
$O(r^{1+\gamma})=O(t^{1+\gamma})$ being absorbed into the leading
$r$-linear term for $t\le t_*(K_{1,\gamma})$ small).  The statement
\eqref{eq:t-lambda} is immediate from \eqref{eq:layer-size}.
\end{proof}

\begin{lemma}[Monotone deficit transfer to the stopped level]
\label{lem:def-transfer}
Let \(U\in\Sclass\), set \(\delta:=\dT(U)\), and let \(E:=\Et\) be a
regular level with \(\lambda:=|U\setminus E|\le\pi/2\).  Then
\begin{equation}\label{eq:def-transfer}
   \dT(E)\le\delta,
\end{equation}
where \(\dT(E)=T(B_{|E|})-T(E)\) is computed without area
renormalisation.  After rescaling \(E\) by
\(s:=\sqrt{\pi/|E|}\) to obtain \(\Etil:=sE\) with \(|\Etil|=\pi\),
\begin{equation}\label{eq:def-transfer-norm}
   \dT(\Etil)=s^4\,\dT(E)
      \le\Bigl(\frac{\pi}{\pi-\lambda}\Bigr)^{2}\delta
      \le 4\delta.
\end{equation}
\end{lemma}

\begin{proof}
Set \(\hat t=v(|E|)\) so \(E=\{u>\hat t\}\) and the function
\(u_E:=u-\hat t\) is the torsion function of \(E\)
(\(-\Delta u_E=1\) in \(E\), \(u_E=0\) on \(\partial E\)).  For
\(\tau>0\) and \(x\in E\), \(u_E(x)>\tau\iff u(x)>\hat t+\tau\), so
\[
   \{u_E>\tau\}=\{u>\hat t+\tau\}\quad(\tau>0),
\]
and on the level \(\{u_E=\tau\}\) we have
\(|\nabla u_E|=|\nabla u|\) and the same perimeter and area.
Consequently, with the obvious notation
\(D_H^E(\tau)=D_H^U(\hat t+\tau)\) and
\(D_I^E(\tau)=D_I^U(\hat t+\tau)\).

Now apply the identity \eqref{eq:level-identity} to \(E\):
\[
   8\pi\,\dT(E)
      =\int_0^{|E|}
      \bigl(D_H^E(v_E(\sigma))+D_I^E(v_E(\sigma))\bigr)
      \,d\sigma.
\]
In the volume variable \(\sigma=|\{u_E>\tau\}|=|\{u>\hat t+\tau\}|\),
the substitution \(\sigma\mapsto s=\sigma\) (over the range
\([0,|E|]\), with the same area variable) makes the integrand on the
right equal to \(D_H^U(v(s))+D_I^U(v(s))\) for \(s\in[0,|E|]\).  Hence
\[
   8\pi\,\dT(E)
      =\int_0^{|E|}\bigl(D_H^U(v(s))+D_I^U(v(s))\bigr)\,ds
      \le\int_0^{|U|}\bigl(D_H^U(v(s))+D_I^U(v(s))\bigr)\,ds
      =8\pi\,\delta,
\]
using non-negativity of the integrand on \((|E|,|U|)\).  This is
\eqref{eq:def-transfer}.

For \eqref{eq:def-transfer-norm}, recall that torsion satisfies
\(T(sE)=s^4 T(E)\) (because \(u_{sE}(y)=s^2 u_E(y/s)\) solves
\(-\Delta=1\) in \(sE\)), and \(B_{|\Etil|}\) is the scaled
\(B_{|E|}\), so \(\dT(\Etil)=s^4\dT(E)\).  With
\(s^2=\pi/(\pi-\lambda)\) and \(\lambda\le\pi/2\), the factor
\(s^4\le 4\).
\end{proof}

\section{From \(D_H\) to pointwise Serrin}

\begin{lemma}[Exact variance identity]\label{lem:variance}
Let \(U\in\Sclass\) and let \(t\) be a regular value of \(u\).  Set
\(\bar f_t:=A_t/P_t=|E_t|/P_t\).  Then
\begin{equation}\label{eq:variance}
   \int_{\partial E_t}\frac{(|\nabla u|-\bar f_t)^2}{|\nabla u|}\,
      d\HH^1
      =
   \frac{A_t}{P_t^{2}}\,D_H(t).
\end{equation}
\end{lemma}

\begin{proof}
By the divergence theorem applied to \(\nabla u\) on \(E_t\),
\[
   \int_{\partial E_t}|\nabla u|\,d\HH^1
      =-\int_{E_t}\Delta u\,dx=|E_t|=A_t.
\]
Hence, expanding the square and using \(\int_{\partial E_t}
(\bar f_t)^2/|\nabla u|=\bar f_t^{\,2}\int 1/|\nabla u|\),
\[
\begin{aligned}
\int_{\partial E_t}\frac{(|\nabla u|-\bar f_t)^2}{|\nabla u|}
   &=
   \int_{\partial E_t}|\nabla u|\,d\HH^1
   -2\bar f_t P_t
   +\bar f_t^{\,2}\int_{\partial E_t}\frac{d\HH^1}{|\nabla u|}\\
   &=A_t-2\bar f_t P_t
     +\frac{A_t^2}{P_t^2}\int_{\partial E_t}\frac{d\HH^1}{|\nabla u|}\\
   &=\frac{A_t}{P_t^2}
     \Bigl[A_t\int_{\partial E_t}\frac{d\HH^1}{|\nabla u|}-P_t^2\Bigr]
   =\frac{A_t}{P_t^2}\,D_H(t),
\end{aligned}
\]
where we used \(2\bar f_t P_t=2A_t\) in the third line.
\end{proof}

\begin{lemma}[\(L^2\) Serrin on the good level]\label{lem:L2}
Let \(\hat t\) be the level produced by Lemma~\ref{lem:good-level} with
\(\lambda\le\lambda_0\).  Then
\begin{equation}\label{eq:L2-serrin}
   \int_{\partial \Et}
      \bigl(|\nabla u|-\bar f_{\hat t}\bigr)^2\,d\HH^1
      \le
   C_{L^2}\,\frac{\delta}{\lambda},
\end{equation}
with \(C_{L^2}\) depending only on \(\Sclass\).
\end{lemma}

\begin{proof}
By \eqref{eq:t-lambda} we have \(\hat t\le\lambda/c_h\le\lambda_0/c_h<
r_0/2\), so the good level lies in the selected collar.  By (S5) and
(S4), $|\nabla u|\le q_++\Lambda_\gamma\,r_0^\gamma$ in the collar
(from the Hölder modulus of $Du$).  By Lemma~\ref{lem:variance} and
$A_{\hat t}\le|U|=\pi$,
\[
   \int_{\partial\Et}
      \bigl(|\nabla u|-\bar f_{\hat t}\bigr)^2\,d\HH^1
      \le
   (q_++\Lambda_\gamma r_0^\gamma)\,
   \int_{\partial\Et}
      \frac{(|\nabla u|-\bar f_{\hat t})^2}{|\nabla u|}\,d\HH^1
      \le
   \frac{(q_++\Lambda_\gamma r_0^\gamma)\,\pi}{P_-^2}\,D_H(\hat t).
\]
Using \(D_H(\hat t)\le 16\pi\delta/\lambda\) from \eqref{eq:good-D}
finishes the proof with
$C_{L^2}=16\pi^2(q_++\Lambda_\gamma r_0^\gamma)/P_-^2$.
\end{proof}

\begin{lemma}[One-dimensional interpolation on a level curve, Hölder
version]\label{lem:interp}
Let \(\Sigma\subset\R^2\) be a compact \(C^1\) curve satisfying the
uniform lower length density
\begin{equation}\label{eq:density}
   \HH^1(\Sigma\cap B_r(x))\ge c_\Sigma\,r
   \qquad(x\in\Sigma,\ 0<r\le r_\Sigma).
\end{equation}
Let \(\gamma\in(0,1]\).  If \(g:\Sigma\to\R\) is $\gamma$-Hölder
continuous with seminorm
\([g]_{C^{0,\gamma}(\Sigma)}\le L_\gamma\) and
\(\|g\|_{L^2(\Sigma)}^2\le B\), then
\begin{equation}\label{eq:interp}
   \|g\|_{L^\infty(\Sigma)}
      \le
   C_{\rm int}\Bigl[B^{\gamma/(1+2\gamma)}\,L_\gamma^{1/(1+2\gamma)}
      +B^{1/2}\Bigr],
   \qquad
   C_{\rm int}=C_{\rm int}(c_\Sigma,r_\Sigma,\gamma)>0.
\end{equation}
\end{lemma}

\begin{proof}
Set \(A:=\|g\|_{L^\infty(\Sigma)}\) and pick \(x_0\in\Sigma\) with
\(|g(x_0)|\ge A/2\).  Assume \(A>0\) (otherwise the bound is trivial).

\emph{Case 1: $L_\gamma>0$ and $r:=(A/(4L_\gamma))^{1/\gamma}\le
r_\Sigma$.}
For $y\in\Sigma\cap B_r(x_0)$,
$|g(y)|\ge|g(x_0)|-L_\gamma|y-x_0|^\gamma\ge A/2-L_\gamma r^\gamma
=A/2-A/4=A/4$, so
\[
   B\ge\int_{\Sigma\cap B_r(x_0)}g^2\,d\HH^1
      \ge\frac{A^2}{16}\,c_\Sigma r
      =\frac{c_\Sigma A^2}{16}\Bigl(\frac{A}{4L_\gamma}\Bigr)^{1/\gamma}
      =\frac{c_\Sigma\,A^{2+1/\gamma}}{16\cdot 4^{1/\gamma}\,L_\gamma^{1/\gamma}},
\]
giving $A^{2+1/\gamma}\le(16\cdot 4^{1/\gamma}/c_\Sigma)\,B\,
L_\gamma^{1/\gamma}$, i.e.,
$A\le C_1(\gamma,c_\Sigma)\,B^{\gamma/(1+2\gamma)}\,
L_\gamma^{1/(1+2\gamma)}$, with
$C_1(\gamma,c_\Sigma):=(16\cdot 4^{1/\gamma}/c_\Sigma)^{\gamma/(1+2\gamma)}$.

\emph{Case 2: $L_\gamma=0$, or $L_\gamma>0$ with $(A/(4L_\gamma))^{1/
\gamma}>r_\Sigma$.}
In either case $|g|\ge A/2-L_\gamma r_\Sigma^\gamma\ge A/4$ on
$\Sigma\cap B_{r_\Sigma}(x_0)$ (the case $L_\gamma=0$ being trivial,
the second case using $L_\gamma r_\Sigma^\gamma<A/4$ by hypothesis),
so
\[
   B\ge\frac{A^2}{16}\,c_\Sigma r_\Sigma,
   \qquad
   A\le\Bigl(\frac{16}{c_\Sigma r_\Sigma}\Bigr)^{1/2}B^{1/2}.
\]
Taking the maximum of the two cases yields \eqref{eq:interp} with
\(C_{\rm int}=\max\bigl(C_1(\gamma,c_\Sigma),(16/(c_\Sigma r_\Sigma))^{1/2}
\bigr)\).  For $\gamma=1$ the exponent $\gamma/(1+2\gamma)=1/3$ and
$1/(1+2\gamma)=1/3$ recover the Lipschitz case
$A\le C(BL)^{1/3}+CB^{1/2}$.
\end{proof}

\begin{lemma}[Pointwise Serrin on the good level]\label{lem:Linf}
Under the hypotheses of Lemma~\ref{lem:L2}, with the inset
regularity constants from Lemma~\ref{lem:inset-regularity} below,
\begin{equation}\label{eq:Linf-serrin}
   \bigl\||\nabla u|-\bar f_{\hat t}\bigr\|_{L^\infty(\partial \Et)}
      \le
   C_{L^\infty}\Bigl[
      \Bigl(\frac{\delta}{\lambda}\Bigr)^{\!\gamma/(1+2\gamma)}
      +\Bigl(\frac{\delta}{\lambda}\Bigr)^{\!1/2}
   \Bigr],
\end{equation}
with $C_{L^\infty}$ depending only on $\Sclass$.
\end{lemma}

\begin{proof}
By (S4), the function $g:=|\nabla u|-\bar f_{\hat t}$ is
$\gamma$-Hölder on $\partial\Et$ with seminorm
$[g]_{C^{0,\gamma}(\partial\Et)}\le\Lambda_\gamma$ (since
$\bar f_{\hat t}$ is constant); by
Lemma~\ref{lem:inset-regularity}, $\partial\Et$ satisfies the
lower length density \eqref{eq:density} with constants
$c_\Sigma,r_\Sigma$ depending only on $\Sclass$.  Apply
Lemma~\ref{lem:interp} with $B=C_{L^2}\delta/\lambda$ from
Lemma~\ref{lem:L2} and $L_\gamma=\Lambda_\gamma$.  The first term
in \eqref{eq:interp} becomes
$C_{\rm int}\cdot C_{L^2}^{\gamma/(1+2\gamma)}\,
\Lambda_\gamma^{1/(1+2\gamma)}\cdot
(\delta/\lambda)^{\gamma/(1+2\gamma)}$;
the second is $C_{\rm int}\cdot C_{L^2}^{1/2}\,(\delta/\lambda)^{1/2}$.
Both fit into \eqref{eq:Linf-serrin} with
$C_{L^\infty}:=C_{\rm int}\,\max\bigl(C_{L^2}^{\gamma/(1+2\gamma)}\,
\Lambda_\gamma^{1/(1+2\gamma)},C_{L^2}^{1/2}\bigr)$.  For
$\gamma=1$ the exponent $\gamma/(1+2\gamma)=1/3$ recovers the
Lipschitz case.
\end{proof}

\section{Inset regularity, recentring, and Brandolini}

\begin{lemma}[Inset regularity]\label{lem:inset-regularity}
Let $U\in\Sclass$.  There are $t_0,\rho_0^*,K_0^*,d_0>0$ depending
only on $\Sclass$ such that every regular value $t\in(0,t_0)$
yields an inset $E_t$ which is open, connected, simply connected,
of diameter $\le d_0$, and bounded by a closed $C^{1,\gamma}$ Jordan
curve with atlas pair $(K_0^*,\rho_0^*)$.
\end{lemma}

\begin{proof}
By (S5) and the Hölder continuity of $\nabla u$ on the collar
(from (S4): $\|Du\|_{C^{0,\gamma}}\le\Lambda_\gamma$), after shrinking
the collar radius once and for all to
$r_0'\in(0,r_0]$ depending only on $\Sclass$, we have
$|\nabla u|\ge q_-/2$ throughout $\mathcal N_{r_0'}(\partial U)\cap U$.

Introduce the normal coordinate map
\[
   \Phi(y,s):=y-s\,\nu_U(y),
   \qquad
   y\in\partial U,\ s\in[0,r_0').
\]
By (S3) ($C^{1,\gamma}$ atlas of $\partial U$), $\nu_U\in
C^{0,\gamma}(\partial U)$, so $\Phi$ is a $C^{1,\gamma}$ map onto the
collar.  Its differential at $s=0$ is the identity-plus-perturbation
on $\partial U\times\R$, so for $r_0'$ small enough, $\Phi$ is a
$C^{1,\gamma}$ diffeomorphism with uniformly bounded norm and inverse
norm.  The composition $F(y,s):=u(\Phi(y,s))$ is $C^{1,\gamma}$ with
norm bounded by $M_{1,\gamma}$ from (S4) and $\partial_s F(y,0)=
|\nabla u(y)|\ge q_-$.  After further shrinking the collar to
$r_0''\le r_0'$ (depending only on $M_{1,\gamma}, q_-,\gamma$) we have
$|\partial_s F|\ge q_-/2$ throughout.

Set $t_0:=\min(q_-r_0''/2,r_0 c_0/2)$.  For $0<t<t_0$, the
quantitative implicit function theorem for $C^{1,\gamma}$ maps gives
a unique $C^{1,\gamma}$ function $s_t:\partial U\to(0,r_0'')$ with
$\|s_t\|_{C^{1,\gamma}}\le K_0^*$ and
\[
   \partial E_t=\{y-s_t(y)\,\nu_U(y):y\in\partial U\}.
\]
This makes $\partial E_t$ a closed $C^{1,\gamma}$ Jordan curve with
atlas pair $(K_0^*,\rho_0^*)$ depending only on $\Sclass$, and makes
$E_t$ a Jordan domain (open, connected, simply connected) by the
normal-graph structure.  The diameter satisfies
$\mathrm{diam}(E_t)\le\mathrm{diam}(U)\le 2R_0=:d_0$.
\end{proof}

\begin{lemma}[Lower length density of $\partial\Et$]
\label{lem:density-quant}
Under the hypotheses of Lemma~\ref{lem:inset-regularity}, the curve
$\Sigma:=\partial \Et$ satisfies \eqref{eq:density} with explicit
constants $c_\Sigma:=1/(2(1+K_0^*))$ and $r_\Sigma:=\rho_0^*/2$.
\end{lemma}

\begin{proof}
By Lemma~\ref{lem:inset-regularity}, $\partial\Et$ admits at every
point a $C^{1,\gamma}$ chart of size $\rho_0^*$ with chart norm
bounded by $K_0^*$.  Since $C^{1,\gamma}\subset
C^{0,1}$ (Lipschitz), the chart is bi-Lipschitz with constant
$\le 1+K_0^*$.  For any disk $B_\rho(x_0)$ centred on $\Sigma$ with
$\rho\le\rho_0^*/2$, the chart pulls $B_\rho(x_0)\cap\Sigma$ back to
a subset of $\partial\R^2_+$ containing an interval of length
$\ge\rho/(1+K_0^*)$; pushing forward (Lipschitz with constant
$\le 1+K_0^*$ in the opposite direction) gives the curve segment of
$\HH^1$-measure $\ge\rho/((1+K_0^*)^2)$.  Hence
$\HH^1(B_\rho(x_0)\cap\Sigma)\ge\rho/(2(1+K_0^*))$ for the slightly
more conservative constant $c_\Sigma=1/(2(1+K_0^*))$.
\end{proof}

\begin{lemma}[Brandolini annular core for the rescaled inset]
\label{lem:Brandolini-core}
Let \(U\in\Sclass\), set \(\delta:=\dT(U)\), and let \(K\ge1\) be a
constant to be chosen later.  Apply
Lemma~\ref{lem:good-level} with \(\lambda=K\sqrt\delta\), assuming
\(\delta\le\delta_1(K,\Sclass)\) so that
\(\lambda<\lambda_0\).  Let \(E:=\Et\) and let \(\Etil:=sE\) with
\(s=\sqrt{\pi/|E|}\) so \(|\Etil|=\pi\).  There exists
\(\delta_2(K,\Sclass)\in(0,\delta_1]\) such that, for
\(\delta\le\delta_2\),
\begin{equation}\label{eq:eta-explicit}
   \eta(K,\delta):=
   2C_{L^\infty}\Bigl(\frac{\sqrt\delta}{K}\Bigr)^{\!\gamma/(1+2\gamma)}
   +C_h^*\sqrt\delta\le\eta_B,
\end{equation}
and there exist \(z\in\R^2\), \(r_{\rm in},r_{\rm out}>0\) such that
\begin{equation}\label{eq:annular-core}
   B_{r_{\rm in}}(z)\subset \Etil\subset B_{r_{\rm out}}(z),
   \qquad
   r_{\rm out}-r_{\rm in}+|r_{\rm in}-1|+|r_{\rm out}-1|
   \le 3\,C_{\rm Br}\,\eta(K,\delta)^\mu.
\end{equation}
Here \(C_h^*=C_h^*(\Sclass)\) is a constant depending only on
\(\Sclass\).
\end{lemma}

\begin{proof}
\emph{Step 1: Pointwise Serrin on $\Etil$.}
By Lemma~\ref{lem:Linf} applied with $\lambda=K\sqrt\delta$,
\[
   \bigl\||\nabla u|-\bar f_{\hat t}\bigr\|_{L^\infty(\partial E)}
      \le
   C_{L^\infty}\Bigl[
      \Bigl(\frac{\sqrt\delta}{K}\Bigr)^{\!\gamma/(1+2\gamma)}
      +\Bigl(\frac{\sqrt\delta}{K}\Bigr)^{\!1/2}
   \Bigr]
   \le
   2C_{L^\infty}\Bigl(\frac{\sqrt\delta}{K}\Bigr)^{\!\gamma/(1+2\gamma)}
\]
for $\delta\le\delta_1$ small (the $1/2$-power being dominated by the
$\gamma/(1+2\gamma)$-power for small $\sqrt\delta/K$, since
$\gamma/(1+2\gamma)\le 1/3<1/2$).

Now consider \(\tilde u(y):=s^2\,u(y/s)\) on \(\Etil\); then
\(-\Delta\tilde u=1\) in \(\Etil\), \(\tilde u=0\) on \(\partial\Etil\),
and \(|\nabla\tilde u|(y)=s\,|\nabla u|(y/s)\).  Set
\(\bar{\tilde f}:=|\Etil|/P(\Etil)=\pi/(sP(E))=s\bar f_{\hat t}\).
Then
\[
   \bigl\||\nabla\tilde u|-\bar{\tilde f}\bigr\|_{L^\infty
     (\partial\Etil)}
   =s\,\bigl\||\nabla u|-\bar f_{\hat t}\bigr\|
     _{L^\infty(\partial E)}
   \le 2s\,C_{L^\infty}\Bigl(\frac{\sqrt\delta}{K}\Bigr)^{\!\gamma/(1+2\gamma)}.
\]
With \(s=\sqrt{\pi/|E|}=1+O(\sqrt\delta)\), the factor \(s\le 2\) for
\(\delta\) small.

\emph{Step 2: \(\bar{\tilde f}\) close to \(1/2\).}  From
\(D_I(\hat t)\le 16\pi\delta/\lambda=16\pi\sqrt\delta/K\) we get, by
the planar isoperimetric inequality applied to \(\Etil\),
\[
   P(\Etil)^2=s^2P(E)^2\le s^2\bigl(4\pi|E|+16\pi\sqrt\delta/K\bigr)
      =4\pi^2+\frac{16\pi^2\sqrt\delta}{K|E|}\cdot 1
      \le 4\pi^2+\frac{32\pi\sqrt\delta}{K},
\]
so \(P(\Etil)\le 2\pi\sqrt{1+(8/(\pi K))\sqrt\delta}
\le 2\pi+(8/K)\sqrt\delta\) (using \(\sqrt{1+x}\le 1+x/2\)).
Consequently
\[
   \Bigl|\bar{\tilde f}-\tfrac12\Bigr|
   =\Bigl|\frac{\pi}{P(\Etil)}-\frac12\Bigr|
   =\frac{P(\Etil)-2\pi}{2P(\Etil)}
   \le\frac{1}{2\cdot 2\pi}\cdot(8/K)\sqrt\delta
   =\frac{2}{\pi K}\sqrt\delta.
\]
Set \(C_h^*:=4/\pi+1\) (the additional \(1\) absorbs the \(s\ne 1\)
inflation).

\emph{Step 3: Apply (B).}  Set \(w:=-2\tilde u\), which satisfies
\(\Delta w=2\) in \(\Etil\) and \(w=0\) on \(\partial\Etil\).  Then
\[
   \bigl\||Dw|-1\bigr\|_{L^\infty(\partial\Etil)}
      =\bigl\|2|\nabla\tilde u|-1\bigr\|_{L^\infty(\partial\Etil)}
      \le 2\bigl\||\nabla\tilde u|-\bar{\tilde f}\bigr\|_{L^\infty}
        +2\bigl|\bar{\tilde f}-\tfrac12\bigr|
      \le\eta(K,\delta),
\]
with \(\eta(K,\delta)\) given by \eqref{eq:eta-explicit}.  For
\(\delta\le\delta_2\) small enough we have
$\eta(K,\delta)\le\eta_B$.  By Lemma~\ref{lem:inset-regularity},
$\Etil$ is bounded, connected, with $C^{1,\gamma}$ atlas pair
$(K_0^*,\rho_0^*)$ depending only on $\Sclass$ (rescaling by $s\sim
1$ preserves these constants up to a universal factor), and
$\|D\tilde u\|_{C^{0,\gamma}(\overline{\Etil})}\le 2\Lambda_\gamma$
\textup{(}by the chain-rule scaling
$\nabla\tilde u(y)=s\nabla u(y/s)$ and $s\le 2$ for $\delta$ small,
with $\Lambda_\gamma$ from (S4)\textup{)}; and diameter $\le 2R_0$.
Apply the $C^{1,\gamma}$ Brandolini theorem (B) of
\S\ref{sec:BNST-import} \textup{(}proved as
Theorem~\ref{thm:BNST-Cgamma}\textup{)}: there exist $z$ and
$r_{\rm in}\le\rho_{\Etil}=1\le r_{\rm out}$ with
$B_{r_{\rm in}}(z)\subset\Etil\subset B_{r_{\rm out}}(z)$ and the
three-term bound \eqref{eq:annular-core} (the third side coming
from BNST's $|r_{\rm in}-1|+|r_{\rm out}-1|\le 2C_{\rm Br}\eta^\mu$).
\end{proof}

\begin{lemma}[Barycentre recentring]\label{lem:recenter}
Let \(\Etil\) and \(z\) be as in Lemma~\ref{lem:Brandolini-core}, and
set \(\eta:=3C_{\rm Br}\eta(K,\delta)^\mu\).  Let
\(y_\star:=\frac1\pi\int_{\Etil}(x-z)\,dx\) be the barycentre of
\(\Etil-z\).  Then \(|y_\star|\le 6\eta\), and
\begin{equation}\label{eq:recenter-bound}
   B_{1-7\eta}(z+y_\star)\subset \Etil\subset
      B_{1+7\eta}(z+y_\star).
\end{equation}
\end{lemma}

\begin{proof}
By Lemma~\ref{lem:Brandolini-core}, \(B_{1-\eta}(z)\subset\Etil\subset
B_{1+\eta}(z)\).  Since the disk \(B_{1-\eta}(z)\) is symmetric about
\(z\), its barycentre contribution to \(\int_{\Etil-z}\!x\,dx\)
vanishes.  Using \(|\Etil\setminus B_{1-\eta}(z)|\le|B_{1+\eta}(z)
\setminus B_{1-\eta}(z)|=\pi\bigl((1+\eta)^2-(1-\eta)^2\bigr)=4\pi\eta\),
\[
   |y_\star|=\frac1\pi\Bigl|\int_{\Etil\setminus B_{1-\eta}(z)}
      (x-z)\,dx\Bigr|
      \le\frac{1+\eta}{\pi}\cdot|\Etil\setminus B_{1-\eta}(z)|
      \le(1+\eta)\cdot 4\eta\le 6\eta,
\]
where the last step uses \(\eta\le\eta_B<1/2\).  Then
\(z':=z+y_\star\) satisfies \(|z'-z|\le 6\eta\), and any \(x\in
B_{1-7\eta}(z')\) lies in \(B_{1-7\eta+6\eta}(z)=B_{1-\eta}(z)
\subset\Etil\); conversely any \(x\in\Etil\subset B_{1+\eta}(z)\)
lies in \(B_{1+\eta+6\eta}(z')=B_{1+7\eta}(z')\).
\end{proof}

\section{Radial-graph promotion from $C^{1,\gamma}$ regularity}
\label{sec:radialgraph}

The Brandolini core (Lemma~\ref{lem:Brandolini-core}) yields a
near-spherical inclusion
\(B_{1-\eta}(z)\subset\Etil\subset B_{1+\eta}(z)\) for the rescaled
inset, but does not by itself say that \(\Etil\) is a radial graph
from \(z\).  The next lemma supplies that final step: combined with
the global \(C^{1,\gamma}\) atlas of \(\partial\Etil\) inherited from
inset regularity (Lemma~\ref{lem:inset-regularity}), the
near-spherical inclusion automatically promotes to a \(C^{1,\gamma}\)
radial-graph parametrisation, provided the Hausdorff gap \(\eta\) is
small relative to the atlas constants.  In that regime there are no
pockets to fill: \(\Etil\) is itself a radial graph, and may be
handed directly to the small-amplitude radial-graph theorem (G).

The criterion below is purely geometric and dimension-free.  We
state and prove it for arbitrary \(n\ge 2\).  In the main chain of
\S\ref{sec:selected-domain} it is invoked at \(n=2\); the
dimension-free statement is what would extend the chain to
arbitrary \(n\) once (G) is replaced by its dimension-free version
(Theorem~\ref{thm:G12-explicit} carries over verbatim to
\(S^{n-1}\), with \(c_G=1/(16 n^2)\); cf.\ the H12 derivation cited
in \S\ref{ss:G12-conclude}).

\begin{lemma}[Radial-graph promotion]\label{lem:radialgraph}
Let \(n\ge 2\), \(\gamma\in(0,1]\), \(K\ge 1\), \(\rho\in(0,1/2)\),
and let \(E\subset\R^n\) be a bounded open set with \(\partial E\) a
closed connected \(C^{1,\gamma}\)-embedded hypersurface admitting a
uniform atlas of \emph{normal-graph charts} of size \(\rho\) and
norm \(K\): around every \(p\in\partial E\),
\[
   \partial E\cap B_\rho(p)
      =\bigl\{p+\xi+h_p(\xi)\,\nu(p):\xi\in T_p\partial E,\
         |\xi|\le\rho\bigr\},
\]
with \(h_p(0)=0\), \(\nabla h_p(0)=0\), and
\(\|h_p\|_{C^{1,\gamma}(B_\rho(0))}\le K\).  There exist
\(\eta_{\rm RG}=\eta_{\rm RG}(K,\rho,\gamma,n)\in(0,\rho/4)\) and
\(K_{\rm RG}=K_{\rm RG}(K,\rho,\gamma,n)>0\) such that, for any
\(z\in\R^n\) and \(\eta\in(0,\eta_{\rm RG}]\) satisfying
\begin{equation}\label{eq:RG-shell}
   B_{1-\eta}(z)\subset E\subset B_{1+\eta}(z),
\end{equation}
there exists a unique \(\varphi\in C^{1,\gamma}(S^{n-1})\) with
\begin{equation}\label{eq:RG-graph}
   \partial E
      =\bigl\{z+(1+\varphi(\omega))\omega:\omega\in S^{n-1}\bigr\},
   \qquad
   E=\bigl\{z+r\omega:0\le r<1+\varphi(\omega),\,\omega\in S^{n-1}
   \bigr\},
\end{equation}
and
\begin{equation}\label{eq:RG-bounds}
   \|\varphi\|_{L^\infty(S^{n-1})}\le\eta,
   \quad
   \|\nabla_\omega\varphi\|_{L^\infty(S^{n-1})}\le\tfrac12,
   \quad
   \|\varphi\|_{C^{1,\gamma}(S^{n-1})}\le K_{\rm RG}.
\end{equation}
\end{lemma}

\begin{proof}
Translate so that \(z=0\) and write \(\widehat p:=p/|p|\) (defined
on \(\partial E\) since \(|p|\ge 1-\eta>1/2\)).

\emph{Step 1: $C^{0,\gamma}$-Hölder bound on the outward unit
normal $\nu$.}  In the chart at \(p\), the outward unit normal at
\(\xi+h_p(\xi)\nu(p)\) is
\(\nu(\xi)=\bigl(\nu(p)-\nabla h_p(\xi)\bigr)/\sqrt{1+|\nabla h_p
(\xi)|^2}\), where we identify \(\nabla h_p\) with its lift to
\(\nu(p)^\perp\subset\R^n\).  Since \(\nabla h_p(0)=0\) and
\([\nabla h_p]_{C^{0,\gamma}}\le K\), \(|\nabla h_p(\xi)|\le K|\xi
|^\gamma\) on \(B_\rho(0)\); differentiating and bounding,
\([\nu]_{C^{0,\gamma}(\partial E\cap B_\rho(p))}\le 2K\) (the factor
\(2\) absorbs the chain-rule contribution of the unit-norm
denominator, which has \(C^{0,\gamma}\) seminorm \(\le K^2\) on
\(B_\rho\)).  Globally, two points \(p_1,p_2\in\partial E\) at
Euclidean distance \(\le\rho/2\) lie in either chart, giving
\(|\nu(p_1)-\nu(p_2)|\le 2K|p_1-p_2|^\gamma\).  For \(p_1,p_2\) at
larger distance, one chains \(\le 3/\rho\) such hops along
\(\partial E\) (whose Euclidean diameter is \(\le 3\) because
\(\partial E\subset B_{1+\eta}\subset B_{3/2}\)), giving a global
seminorm
\begin{equation}\label{eq:nu-Holder}
   [\nu]_{C^{0,\gamma}(\partial E)}
      \le K_\nu,
   \qquad
   K_\nu:=6K(1+K/\rho).
\end{equation}

\emph{Step 2: Radial distance and its tangential gradient.}
Set \(\psi(p):=|p|-1\) for \(p\in\partial E\).  Then \(|\psi|\le
\eta\) on \(\partial E\).  The tangential gradient is
\(\nabla^{\partial E}\psi(p)=\widehat p-(\widehat p\cdot\nu(p))
\nu(p)\), so
\begin{equation}\label{eq:DEpsi-formula}
   |\nabla^{\partial E}\psi(p)|^{2}
      =1-(\widehat p\cdot\nu(p))^{2}
      =\sin^{2}\theta(p),
\end{equation}
where \(\theta(p)\in[0,\pi]\) is the angle between \(\nu(p)\) and
\(\widehat p\).

\emph{Step 3: Hölder seminorm of $\nabla^{\partial E}\psi$.}
The map \(p\mapsto\widehat p\) on the shell \(\{1-\eta\le|p|\le 1+
\eta\}\) is smooth with Lipschitz constant \(\le 1/(1-\eta)\le 2\)
for \(\eta\le 1/2\), so its \(C^{0,\gamma}\) seminorm on
\(\partial E\) is bounded by \(2\cdot\mathrm{diam}(\partial E)^{1
-\gamma}\le 2\cdot 3^{1-\gamma}\le 6\).  Combining with
\eqref{eq:nu-Holder},
\begin{equation}\label{eq:DEpsi-Holder}
   [\nabla^{\partial E}\psi]_{C^{0,\gamma}(\partial E)}
      \le L_\nu,
   \qquad
   L_\nu:=6+5 K_\nu=6+30 K(1+K/\rho).
\end{equation}

\emph{Step 4: Hölder interpolation.}
By the Gagliardo--Nirenberg--Hölder inequality on a
\(C^{1,\gamma}\) surface of bounded geometry, applied to
\(\psi\):
\begin{equation}\label{eq:GN}
   \|\nabla^{\partial E}\psi\|_{L^\infty(\partial E)}
   \le C_{\rm GN}\Bigl(
      \|\psi\|_{L^\infty}^{\gamma/(1+\gamma)}\,
      [\nabla^{\partial E}\psi]_{C^{0,\gamma}}^{1/(1+\gamma)}
      +\|\psi\|_{L^\infty}
   \Bigr),
\end{equation}
with \(C_{\rm GN}=C_{\rm GN}(K,\rho,\gamma,n)\).  \textup{(Direct
proof: at the point \(p_0\) realising \(\|\nabla^{\partial E}\psi\|
_\infty=:A\), localise to the chart of size \(\rho\) at \(p_0\); in
flat coordinates the inequality reduces to the standard Euclidean
\(\R^{n-1}\) interpolation, with constants depending on \(K\) and
\(\rho\).)}  Substituting Step~2 (\(\|\psi\|_\infty\le\eta\)) and
\eqref{eq:DEpsi-Holder},
\[
   \sup_{\partial E}\sin\theta
   =\|\nabla^{\partial E}\psi\|_\infty
   \le C_{\rm GN}\bigl(\eta^{\gamma/(1+\gamma)}\,L_\nu^{1/(1+\gamma)}
      +\eta\bigr).
\]

\emph{Step 5: Choice of $\eta_{\rm RG}$.}
Set
\begin{equation}\label{eq:etaRG-def}
   \eta_{\rm RG}:=\min\Bigl\{\tfrac{\rho}{4},\
      \bigl(16\,C_{\rm GN}\,(L_\nu^{1/(1+\gamma)}+1)\bigr)^{-(1+\gamma)
      /\gamma}\Bigr\}.
\end{equation}
For \(\eta\le\eta_{\rm RG}\) the right-hand side of Step~4 is
\(\le 1/8\), so
\begin{equation}\label{eq:cos-bound}
   \sin\theta(p)\le 1/8,
   \quad
   |\widehat p\cdot\nu(p)|=|\cos\theta(p)|\ge\sqrt{63}/8>15/16
   \qquad(p\in\partial E).
\end{equation}

\emph{Step 6: $\widehat p\cdot\nu(p)>0$ throughout $\partial E$.}
The function \(p\mapsto\widehat p\cdot\nu(p)\) is continuous on the
connected set \(\partial E\) and never vanishes by
\eqref{eq:cos-bound}, so its sign is constant on \(\partial E\).
Pick \(p_\star\in\partial E\) with \(|p_\star|=\max_{\partial
E}|\cdot|=:r_\star\in[1,1+\eta]\).  Then \(p_\star\) is an interior
maximum of \(|q|^{2}\) restricted to \(\partial E\), so
\(\nabla^{\partial E}|q|^2(p_\star)=0\), i.e.,
\(p_\star=(p_\star\cdot\nu(p_\star))\nu(p_\star)\), and
\(\nu(p_\star)=\pm\widehat p_\star\).  Outward orientation (taking
\(\nu(p_\star)\) to point out of \(E\), where \(0\in E\) is on the
inward side) selects the \(+\) sign, so \(\widehat p_\star\cdot
\nu(p_\star)=+1\).  By constant sign, \(\widehat p\cdot\nu(p)>0\)
everywhere on \(\partial E\).

\emph{Step 7: The radial map $F$ is a $C^{1,\gamma}$
diffeomorphism.}  Define \(F:\partial E\to S^{n-1}\) by
\(F(p):=\widehat p\).  \(F\) is \(C^{1,\gamma}\) (composition of the
inclusion \(\partial E\hookrightarrow\R^n\setminus\{0\}\) and the
smooth normalisation).  The differential at \(p\) is
\[
   dF_p:T_p\partial E\to T_{\widehat p}S^{n-1},
   \qquad
   dF_p(v)=|p|^{-1}\bigl(v-(v\cdot\widehat p)\widehat p\bigr),
\]
whose kernel is \(\{v\in T_p\partial E:v\parallel\widehat p\}\) —
trivial precisely when \(\widehat p\notin T_p\partial E\), i.e., when
\(\widehat p\cdot\nu(p)\ne 0\).  By Step~6 this holds at every
\(p\), so \(F\) is a local \(C^{1,\gamma}\) diffeomorphism.

Compactness of \(\partial E\) makes \(F\) proper.  A proper local
diffeomorphism between connected boundaryless manifolds of equal
dimension is a covering map.  For \(n\ge 3\), \(S^{n-1}\) is simply
connected, so the covering is trivial and \(F\) is a global
\(C^{1,\gamma}\) diffeomorphism onto \(S^{n-1}\).  For \(n=2\), the
degree of \(F\) equals the winding number of \(\partial E\) around
\(0\in B_{1-\eta}\subset E\), which is \(\pm 1\) (\(\partial E\) is
a Jordan curve enclosing \(0\)); the orientation
\(\widehat p\cdot\nu(p)>0\) forces \(+1\), so \(F\) is again a
\(C^{1,\gamma}\) diffeomorphism.

\emph{Step 8: Construction of $\varphi$.}
Set \(\varphi(\omega):=|F^{-1}(\omega)|-1\) for \(\omega\in S^{n-1}\).
Then \(\varphi\in C^{1,\gamma}(S^{n-1})\) as a composition of
\(C^{1,\gamma}\) maps; the shell inclusion gives
\(\|\varphi\|_\infty\le\eta\).  The radial-graph identity
\eqref{eq:RG-graph} follows from \(F^{-1}(\omega)=|F^{-1}(\omega)|
\widehat F^{-1}(\omega)=(1+\varphi(\omega))\omega\), since \(F(F^{-1}
(\omega))=\omega\).  A chart-by-chart computation (differentiating
\(|F^{-1}(\omega)|\) via \(d(F^{-1})=(dF)^{-1}\), bounded by
Step~6) yields the \(C^{1,\gamma}\)-norm bound
\(\|\varphi\|_{C^{1,\gamma}(S^{n-1})}\le K_{\rm RG}(K,\rho,\gamma,n)\)
and, more sharply, the slope bound:
\[
   |\nabla_\omega\varphi(\omega)|
      =\bigl|\,|p|\cdot(\widehat p\cdot\nu(p))^{-1}\,\sin\theta(p)
      \,\bigr|
      \le(1+\eta)\cdot(15/16)^{-1}\cdot 1/8\le 1/2,
\]
evaluated at \(p=F^{-1}(\omega)\), giving \eqref{eq:RG-bounds}.
\end{proof}

\begin{remark}\label{rem:RG-dimfree}
Lemma~\ref{lem:radialgraph} is dimension-free and uses no smallness
of the atlas norm \(K\) or chart radius \(\rho\): the threshold
\(\eta_{\rm RG}\) absorbs both polynomially.  In the application of
\S\ref{sec:selected-domain} below, \(K=K_0^\ast\) and
\(\rho=\rho_0^\ast\) of Lemma~\ref{lem:inset-regularity} depend only
on \(\Sclass\), so \(\eta_{\rm RG}\) is a universal constant
\(\eta_{\rm RG}(\Sclass,\gamma)\), independent of \(\delta\).  Since
the Brandolini gap \(\eta\) of \eqref{eq:annular-core} is
\(O(\delta^{\mu\gamma/(2(1+2\gamma))})\to 0\) as \(\delta\to 0\), the
hypothesis \(\eta\le\eta_{\rm RG}\) is automatic for
\(\delta\le\delta_\star(\Sclass,\gamma)\).
\end{remark}

\section{Hole-filling: an alternative route for $n=2$}
\label{sec:hole-fill}

This section records an alternative, polar-slicing argument that, in
dimension two only, replaces the radial-graph promotion of
\S\ref{sec:radialgraph} by a hole-filling step.  The route uses no
smallness condition on \(\eta\) relative to the atlas norm — only the
Brandolini inclusion and the perimeter deficit
\(P(\Etil)\le 2\pi+\eps\).  The resulting bound
\(|G\setminus\Etil|\le\eps^{2}/(4\pi)\) is \emph{sharp} in
\(\delta\): \(O(\delta)\), same scaling as the radial-graph route
of \S\ref{sec:radialgraph}.  This sharp \(O(\delta)\) scaling
relies on the planar isoperimetric inequality
\(|Q|\le P(Q)^{2}/(4\pi)\) and the super-additivity inequality
\(\sum a_j^{2}\le(\sum a_j)^{2}\); the corresponding \(n\)-dimensional
inequality \(|Q|\le C_n P(Q)^{n/(n-1)}\) gives only
\(v\le C\eps^{n/(n-1)}=O(\delta^{n/(2(n-1))})\), strictly weaker than
\(O(\delta)\) for \(n\ge 3\).  Hence the polar-slicing route is
\emph{planar} in essence; the radial-graph promotion of
\S\ref{sec:radialgraph} provides the dimension-free counterpart.
\S\ref{sec:selected-domain}, Remark~\ref{rem:planar-alt} records the
interchange.

\subsection{The polar sheet lemma and quadratic pocket bound}

\begin{lemma}[Polar sheet decomposition]\label{lem:hull-decomp}
Let \(E\subset\R^2\) be a bounded open set with \(C^{1,\gamma}\) boundary,
satisfying
\begin{equation}\label{eq:core}
   B_{r_-}(z_0)\subset E\subset B_{r_+}(z_0),
   \qquad
   0<r_-<r_+<\infty.
\end{equation}
Define
\[
   R(\theta):=\sup\{r>0:z_0+r e^{i\theta}\in E\}\quad(\theta\in S^1),
   \qquad
   G:=\bigl\{z_0+re^{i\theta}:0\le r<R(\theta),\theta\in S^1\bigr\}.
\]
Let \(\{Q_j\}_{j\in J}\) denote the connected components of the open
set \(G\setminus\overline E\).  Then \(G\) is open, of finite perimeter,
each \(Q_j\) is open and of finite perimeter, and
\begin{equation}\label{eq:hull-decomp}
   P(G)+\sum_{j\in J}P(Q_j)\le P(E).
\end{equation}
\end{lemma}

\begin{proof}
Translate so that \(z_0=0\).  Use the bi-Lipschitz polar map
\[
   \Phi:S^1\times(0,r_+)\to B_{r_+}\setminus\{0\},
   \qquad
   \Phi(\theta,r)=re^{i\theta},
   \qquad
   \det D\Phi=r.
\]
Lengths in the cylinder \(S^1\times(0,r_+)\) with metric
\(g=d\theta^2+r^{-2}dr^2\) match Euclidean arclengths under \(\Phi\),
since the Jacobian factor cancels.  For our perimeter computation
below it is sufficient that \(\Phi\) is bi-Lipschitz on
\(S^1\times(r_-/2,r_+)\), which forces both Euclidean and cylindrical
perimeters to agree up to bi-Lipschitz constants on every subset of
that annulus.

\emph{Step 1: Polar slice structure.}  Because \(\partial E\) is
\(C^{1,\gamma}\) and \(E\) is open with \(B_{r_-}\subset E\), for every
direction \(\theta\) the radial slice \(E_\theta:=
\{r>0:re^{i\theta}\in E\}\) is open in \((0,\infty)\) and contains
\((0,r_-]\).  Sard's theorem applied to the restriction of \(\partial
E\) to a polar chart shows that, for \(\HH^1\)-a.e.\
\(\theta\in S^1\), the radial line \(\{re^{i\theta}:r>0\}\) is
transverse to \(\partial E\) (it crosses \(\partial E\) in finitely
many points, none of which is a tangent point).  Call such \(\theta\)
\emph{regular}.

For regular \(\theta\) the slice \(E_\theta\) is a finite union of
disjoint open intervals starting at \(0\):
\[
   E_\theta=(0,r_1(\theta))\sqcup(r_2(\theta),r_3(\theta))\sqcup
       \cdots\sqcup(r_{2m(\theta)}(\theta),r_{2m(\theta)+1}(\theta)),
\]
with crossings \(0<r_1<r_2<\cdots<r_{2m+1}\le r_+\), and \(R(\theta)=
r_{2m(\theta)+1}(\theta)\) is the outermost crossing.  All
\(r_k(\theta)e^{i\theta}\in\partial E\) for regular \(\theta\).

\emph{Step 2: Sheets are graphs.}  For each \(k\in\N\), set
\[
   \Sigma_k:=\{r_k(\theta)e^{i\theta}:
      \theta\hbox{ regular},\ m(\theta)\ge\lceil k/2\rceil\}.
\]
By the implicit function theorem applied to the \(C^{1,\gamma}\) defining
function of \(\partial E\) (which is non-degenerate radially at
regular \(\theta\)), \(\Sigma_k\) is locally a \(C^{1,\gamma}\) polar graph
\(r=r_k(\theta)\), and
\[
   \HH^1(\Sigma_k)=\int_{\{\theta:\ \Sigma_k\hbox{ defined}\}}
      \sqrt{r_k(\theta)^2+r_k'(\theta)^2}\,d\theta.
\]
The sheets \(\{\Sigma_k\}_{k\ge 1}\) are pairwise disjoint, and the
projection \(\pi_{\rm rad}:\partial E\to S^1\),
\(\pi_{\rm rad}(re^{i\theta})=\theta\), is Lipschitz on
\(\partial E\cap\{r\ge r_-\}\) (with constant \(1/r_-\)), so by the
area formula
\(\HH^1\bigl(\pi_{\rm rad}^{-1}(\hbox{regular }\theta)\bigr)=
\HH^1(\partial E)-\HH^1\bigl(\pi_{\rm rad}^{-1}(\hbox{singular }
\theta)\bigr)=\HH^1(\partial E)\), because the singular angle set is
\(S^1\)-null and the radial projection has finite multiplicity
(bounded by \(2m_{\max}+1\)).  Hence
\begin{equation}\label{eq:sheet-PE}
   P(E)=\HH^1(\partial E)=\sum_{k\ge 1}\HH^1(\Sigma_k).
\end{equation}

\emph{Step 3: \(\partial G\) is the outermost sheet.}  By definition
\(R(\theta)=r_{2m(\theta)+1}(\theta)\), so for each regular \(\theta\)
the boundary \(\partial G\) above direction \(\theta\) is the single
point \(R(\theta)e^{i\theta}=r_{2m(\theta)+1}(\theta)e^{i\theta}\),
which lies in \(\Sigma_{2m(\theta)+1}\).  Indexing sheets so that the
outermost piece always belongs to a single distinguished sheet
\(\Sigma_{\rm out}:=\bigcup_\theta\{R(\theta)e^{i\theta}\}\),
\[
   P(G)=\HH^1(\Sigma_{\rm out})
      \le\sum_{k:\ k\hbox{ is outermost at some }\theta}\HH^1(\Sigma_k).
\]
Concretely \(\Sigma_{\rm out}\) is obtained from
\(\{(\theta,R(\theta))\}\) under \(\Phi\), and its \(\HH^1\)-measure is
the polar-graph perimeter of \(G\).

\emph{Step 4: Pockets are bounded by sheet pairs.}  Each \(Q_j\) is a
connected component of \(G\setminus\overline E\), hence \(Q_j\subset
G\setminus E\).  Since \(B_{r_-}\subset E\), \(Q_j\) lies in the open
annulus \(\{r_-<r<r_+\}\) where \(\Phi\) is bi-Lipschitz, so \(\Phi
^{-1}(Q_j)\) is a connected open subset of the cylinder.  For each
regular \(\theta\) in the angular range of \(Q_j\), the slice
\(\Phi^{-1}(Q_j)\cap\{\theta\}\times(0,r_+)\) is a single interval
\((r_{2i-1}(\theta),r_{2i}(\theta))\) with \(i\le m(\theta)\) (it is
non-empty by definition of the angular range; uniqueness of \(i\)
follows from connectedness of \(\Phi^{-1}(Q_j)\)), and we set
\(i=i_j(\theta)\).  Because \(\Phi^{-1}(Q_j)\) is connected and the
pair indices \(i_j(\theta)\) are integer-valued and locally constant
in \(\theta\) on the angular range (the radial coordinates
\(r_{2i-1},r_{2i}\) vary continuously, so a jump in \(i\) would
disconnect the slice), \(i_j(\theta)\) is in fact constant on the
angular range of \(Q_j\); call its common value \(i_j\).  By Sard's
theorem applied to \(\theta\mapsto r_{2i_j-1}(\theta)-r_{2i_j}(\theta)\)
on the angular range, the endpoint angles where
\(r_{2i_j-1}=r_{2i_j}\) form a finite set (for the \(C^{1,\gamma}\)
boundary) at which the two sheet branches meet transversally as
\(C^{1,\gamma}\) arcs of \(\partial E\).  At such an endpoint, the pocket
closes off via a single \(C^{1,\gamma}\) tangent meeting (no radial wall
piece is created, because \(Q_j\) is connected and the meeting is
internal to \(\partial E\), which is a single Jordan curve).  Hence
\[
   P(Q_j)=\HH^1(\partial Q_j)
      =\HH^1\bigl(\Sigma_{2i_j-1}\cap{\rm angle\ range\ of\ }Q_j\bigr)
      +\HH^1\bigl(\Sigma_{2i_j}\cap{\rm angle\ range\ of\ }Q_j\bigr).
\]
Distinct pockets have disjoint angular ranges (because each pocket
contains an open interior, so different components are separated by
\(\partial E\) intersecting some interior ray), and the union of all
pocket-sheet pieces is contained in the union of all inner sheets
(those with \(k<2m(\theta)+1\)):
\[
   \sum_j P(Q_j)\le \sum_{k\ne {\rm out}}\HH^1(\Sigma_k)
      =P(E)-P(G),
\]
using \eqref{eq:sheet-PE} and Step~3.  Rearranging gives
\eqref{eq:hull-decomp}.
\end{proof}

\begin{remark}\label{rem:lipschitz}
The hypothesis $C^{1,\gamma}$ on $\partial E$ (for any $\gamma>0$)
is used in Step~1 and Step~2 for the transverse polar slicing.  In
the application below, $E=\Etil$ inherits the $C^{1,\gamma}$ atlas
pair $(K_0^*,\rho_0^*)$ from the inset regularity
Lemma~\ref{lem:inset-regularity}, so the hypothesis is satisfied.
\end{remark}

\begin{lemma}[Quadratic pocket area]\label{lem:pocket}
Let \(E\subset\R^2\) be bounded, open, with \(C^{1,\gamma}\) boundary,
\(|E|=\pi\), satisfying \eqref{eq:core} and
\begin{equation}\label{eq:perim-eps}
   P(E)\le 2\pi+\eps.
\end{equation}
Let \(G\) be the radial hull of \(E\) about \(z_0\).  Then
\begin{equation}\label{eq:pocket-area}
   |G\setminus E|\le\frac{\eps^2}{4\pi}.
\end{equation}
\end{lemma}

\begin{proof}
Set \(v:=|G\setminus E|\) and \(h:=\sum_j P(Q_j)\).  By
Lemma~\ref{lem:hull-decomp},
\[
   2\pi+\eps\ge P(E)\ge P(G)+h\ge 2\sqrt{\pi(\pi+v)}+h\ge 2\pi+h,
\]
where the third inequality is the planar isoperimetric inequality
applied to \(G\) (which has area \(\pi+v\)), and the last uses
\(\sqrt{\pi(\pi+v)}\ge\pi\) for \(v\ge 0\).  Hence \(h\le\eps\).
Apply the isoperimetric inequality to each pocket
\(|Q_j|\le P(Q_j)^2/(4\pi)\) and sum, using \(\sum a_j^2\le
(\sum a_j)^2\) for non-negative \(a_j\):
\[
   v=\sum_j|Q_j|\le\frac1{4\pi}\sum_j P(Q_j)^2
      \le\frac1{4\pi}\Bigl(\sum_j P(Q_j)\Bigr)^{\!2}
      \le\frac{\eps^2}{4\pi}.
\]
\end{proof}

\subsection{Filling, rescaling, and asymmetry transfer}

\begin{lemma}[Filling and rescaling]\label{lem:fill-rescale}
Let \(E\subset H\subset B_{R_*}\) be open planar sets with
\(|E|=\pi\) and \(|H|=\pi+v\), \(0\le v\le \pi/2\).  Let
\(s:=(\pi/|H|)^{1/2}\) and \(\Gtil:=sH\), so \(|\Gtil|=\pi\).  Then
\begin{equation}\label{eq:fill-rescale}
   \dT(\Gtil)\le 2\bigl(\dT(E)+v\bigr).
\end{equation}
\end{lemma}

\begin{proof}
Torsional rigidity is monotone under inclusion (replace \(u_H\) by
\(u_H|_E\) in the variational characterisation; we obtain
\(T(H)\ge T(E)\)).  Using \(T(B_A)=A^2/(8\pi)\),
\[
   T(B_H)-T(B_E)=\frac{(\pi+v)^2-\pi^2}{8\pi}
      =\frac{2\pi v+v^2}{8\pi}\le\frac{v}{4}+\frac{v^2}{8\pi}
      \le\frac{v}{4}+\frac{v}{16}\le\frac{v}{2}.
\]
Therefore
\[
   \dT(H)=T(B_H)-T(H)\le T(B_H)-T(E)
      =\bigl(T(B_H)-T(B_E)\bigr)+\dT(E)\le\dT(E)+v/2.
\]
Scaling by \(s\) multiplies torsion by \(s^4\) and sends \(B_H\) to
\(B_\pi\); so \(\dT(\Gtil)=s^4\dT(H)\le(\pi/|H|)^2(\dT(E)+v/2)\le
4(\dT(E)+v/2)\le 2(\dT(E)+v)\) (for the chosen range of \(v\)).
\end{proof}

\begin{lemma}[Fraenkel asymmetry under symmetric difference and
rescaling]\label{lem:AF-lip}
For bounded measurable \(E_1,E_2\subset\R^2\) with
\(|E_1|=|E_2|=:V>0\),
\begin{equation}\label{eq:AF-lip}
   |\AF(E_1)-\AF(E_2)|\le\frac{|E_1\Delta E_2|}{V}.
\end{equation}
The Fraenkel asymmetry is dilation-invariant: \(\AF(sE)=\AF(E)\) for
\(s>0\).  It is also translation-invariant.
\end{lemma}

\begin{proof}
For any disk \(B\) with \(|B|=V\), \(|E_1\Delta B|\le|E_1\Delta E_2|
+|E_2\Delta B|\), so taking the infimum over \(B\) gives
\(V\AF(E_1)\le |E_1\Delta E_2|+V\AF(E_2)\), and similarly with the
roles exchanged.  Dilation and translation invariance follow
because \(\Delta\) commutes with dilation/translation and the disk
class is preserved.
\end{proof}

\section{The selected-domain theorem}\label{sec:selected-domain}

\begin{theorem}[Hybrid estimate for selected planar domains]
\label{thm:selected}
There exist \(C_\star=C_\star(\Sclass)>0\) and
\(\delta_\star=\delta_\star(\Sclass)>0\) such that, for every
\(U\in\Sclass\) with \(\delta:=\dT(U)\le\delta_\star\),
\begin{equation}\label{eq:selected-final}
   \AF(U)^2\le C_\star\,\delta.
\end{equation}
\end{theorem}

\begin{proof}
We argue along the dimension-free chain of \S\ref{sec:radialgraph},
applying the radial-graph promotion Lemma~\ref{lem:radialgraph} to
the rescaled inset $\Etil$ once Brandolini delivers the
near-spherical inclusion.  The hole-filling alternative of
\S\ref{sec:hole-fill} is recorded in Remark~\ref{rem:planar-alt}.

Fix \(K=1\) (any other choice of \(K\ge 1\) works and yields a
\(K\)-dependent constant; \(K=1\) is taken for definiteness).  Let
\(\delta\le\delta_2(1,\Sclass)\) be small enough that
Lemma~\ref{lem:Brandolini-core} applies, and so that all subsequent
smallness conditions hold; we record the final smallness threshold
as \(\delta_\star\).

\emph{Step~1: Good level and rescaled inset.}
By Lemma~\ref{lem:good-level} with \(\lambda=\sqrt\delta\), let
\(E=\Et\) be the resulting good level, satisfying
\(\sqrt\delta/2\le|U\setminus E|\le\sqrt\delta\) and
\(D_H(\hat t)+D_I(\hat t)\le 16\pi\sqrt\delta\).  By
Lemma~\ref{lem:def-transfer}, \(\dT(\Etil)\le 4\delta\), where
\(\Etil:=sE\) with \(s=\sqrt{\pi/|E|}\in[1,(1-\sqrt\delta/\pi)^{-1/2}]\).

\emph{Step~2: Brandolini core and first barycentre recentring.}
Set \(\eta:=3C_{\rm Br}\eta(1,\delta)^\mu\) (Lemma~\ref{lem:Brandolini-core}
notation).  Then $\Etil$ admits a centre $z\in\R^2$ with
\(B_{1-\eta}(z)\subset\Etil\subset B_{1+\eta}(z)\).  By
Lemma~\ref{lem:recenter}, with \(z':=z+y_\star\) and
\(|y_\star|\le 6\eta\),
\begin{equation}\label{eq:rec-annulus}
   B_{1-7\eta}(z')\subset\Etil\subset B_{1+7\eta}(z').
\end{equation}

\emph{Step~3: Radial-graph promotion at $z'$.}
By Lemma~\ref{lem:inset-regularity}, $\partial\Etil$ admits a uniform
$C^{1,\gamma}$ atlas of normal-graph charts of size $\rho_0^\ast$ and
norm $K_0^\ast$, with both constants depending only on $\Sclass$
(the rescaling by $s\sim 1$ inflating them by a universal factor).
Let
\(\eta_{\rm RG}:=\eta_{\rm RG}(K_0^\ast,\rho_0^\ast,\gamma,2)\) be
the threshold of Lemma~\ref{lem:radialgraph}.  Since $\eta\to 0$ as
$\delta\to 0$ (with rate $O(\delta^{\mu\gamma/(2(1+2\gamma))})$;
cf.\ Lemma~\ref{lem:Brandolini-core}), we tighten $\delta_\star$ so
that
\begin{equation}\label{eq:smallness-RG-1}
   7\eta\le\eta_{\rm RG}.
\end{equation}
Apply Lemma~\ref{lem:radialgraph} with $E=\Etil$, $z=z'$, and the
shell gap $7\eta$: there is a unique
$\varphi_1\in C^{1,\gamma}(S^{1})$ with
\(\partial\Etil=\{z'+(1+\varphi_1(\omega))\omega:\omega\in S^{1}\}\)
and $\|\varphi_1\|_\infty\le 7\eta$.  In particular $\Etil$ is a
radial graph from $z'$; there are no pockets.

\emph{Step~4: Second barycentre recentring on the radial graph.}
Let \(y_\star':=\pi^{-1}\int_{\Etil}(x-z')\,dx\) be the barycentre of
$\Etil-z'$.  By Lemma~\ref{lem:recenter} applied to
\eqref{eq:rec-annulus} (with $\eta$ replaced by $7\eta$),
\(|y_\star'|\le 6\cdot 7\eta=42\eta\), and
\begin{equation}\label{eq:rec-annulus-2}
   B_{1-49\eta}(z'')\subset\Etil\subset B_{1+49\eta}(z''),
   \qquad
   z'':=z'+y_\star'.
\end{equation}
Tighten $\delta_\star$ further so that
\begin{equation}\label{eq:smallness-RG-2}
   49\eta\le\eta_{\rm RG},
   \qquad
   49\eta\le\eta_G,
\end{equation}
where $\eta_G$ is the (G)-threshold of \S\ref{sec:BNST-import}.
Apply Lemma~\ref{lem:radialgraph} a second time with $z=z''$ and
shell gap $49\eta$: there is
$\varphi_2\in C^{1,\gamma}(S^{1})$ with
\(\partial\Etil=\{z''+(1+\varphi_2(\omega))\omega:\omega\in S^{1}\}\)
and $\|\varphi_2\|_\infty\le 49\eta$.  By construction $z''$ is the
barycentre of $\Etil$: \(\int_{\Etil}(x-z'')\,dx=0\).

\emph{Step~5: Apply (G).}  The translated $\Etil$ at centre $z''$
satisfies all hypotheses of (G) of \S\ref{sec:BNST-import}:
$|\Etil|=\pi$, radial centre at the barycentre,
$\|\varphi_2\|_\infty\le\eta_G$ by \eqref{eq:smallness-RG-2}.
Together with translation invariance of $\dT$,
\begin{equation}\label{eq:G-on-Etil}
   \AF(\Etil)^{2}\le C_G\,\dT(\Etil)\le 4 C_G\,\delta.
\end{equation}

\emph{Step~6: Transfer back to $U$.}  By
Lemma~\ref{lem:AF-lip} (dilation invariance), $\AF(\Etil)=\AF(E)$.
We compare $E$ and $U$ via a concentric disk pair: let $B'$ be the
optimal disk for $E$ (area $|E|=\pi-\lambda$, $\lambda\le\sqrt\delta$)
and $B$ its concentric dilation to area $\pi$.  Then
\(|U\Delta B|\le|U\Delta E|+|E\Delta B'|+|B\Delta B'|
   =|U\setminus E|+|E|\AF(E)+\lambda\le 2\sqrt\delta+\pi\AF(E)\),
giving
\begin{equation}\label{eq:transfer-U}
   \AF(U)
   \le\frac{|U\Delta B|}{\pi}
   \le\frac{2\sqrt\delta}{\pi}+\AF(E)
   \le\frac{2\sqrt\delta}{\pi}+\sqrt{4 C_G\delta}
   =\Bigl(\frac{2}{\pi}+2\sqrt{C_G}\Bigr)\sqrt\delta
   =:C_2\sqrt\delta.
\end{equation}
Squaring, \(\AF(U)^{2}\le C_2^{2}\,\delta\), giving
\eqref{eq:selected-final} with $C_\star:=C_2^{2}=(2/\pi+2\sqrt{C_G})
^{2}$.
\end{proof}

\begin{remark}[Alternative planar route via hole-filling]
\label{rem:planar-alt}
For $n=2$, Steps~3--5 above admit a polar-slicing alternative that
uses the planar lemmas of \S\ref{sec:hole-fill} in place of
Lemma~\ref{lem:radialgraph}, at the cost of slightly larger constants
but without any smallness condition on $\eta$ relative to the atlas
norm.  Concretely, replace Steps~3--5 by:
\begin{enumerate}[label=(\arabic*)]
\item Form the radial hull $G$ of $\Etil$ about $z'$
\textup{(Lemma~\ref{lem:hull-decomp})}.  The same rescaled
isoperimetric estimate as in Step~2 of
Lemma~\ref{lem:Brandolini-core} gives $P(\Etil)\le 2\pi+4\sqrt\delta$.
By Lemma~\ref{lem:pocket} with $\eps=4\sqrt\delta$,
\(v:=|G\setminus\Etil|\le 16\delta/(4\pi)=4\delta/\pi\).
\item Set $\Gtil:=s' G$, $s'=(\pi/|G|)^{1/2}\in(0,1]$;
Lemma~\ref{lem:fill-rescale} gives $\dT(\Gtil)\le 2(\dT(\Etil)+v)
\le 12\delta$.  The same containment-plus-rescaling computation as in
the old Step~4 shows $B_{s'(1-7\eta)}(s' z')\subset\Gtil\subset
B_{s'(1+7\eta)}(s' z')$; reading off the centred amplitude,
$\|\widetilde\varphi\|_\infty\le 8\eta$.
\item Recentre $\Gtil$ to its barycentre, with the same calculation
as in Lemma~\ref{lem:recenter}: $|y_\star^{(\Gtil)}|\le 48\eta$ and
the centred shell becomes $B_{1-56\eta}(z''')\subset\Gtil\subset
B_{1+56\eta}(z''')$.  Tighten $\delta_\star$ to ensure
$56\eta\le\eta_G$; then (G) gives
$\AF(\Gtil)^{2}\le 12 C_G\,\delta$.
\item Transfer from $\Gtil$ back to $\Etil$ via
Lemma~\ref{lem:AF-lip} and the symmetric-difference bound
$|\Etil\Delta\Gtil|\le|\Etil\Delta G|+|G\Delta\Gtil|\le 4\delta/\pi+
8\delta/\pi=12\delta/\pi$, giving
$\AF(\Etil)\le 2\sqrt{3 C_G\delta}+12\delta/\pi^{2}$.  Then
\eqref{eq:transfer-U} runs unchanged, with the constant
$C_2$ replaced by $C_2^{\rm alt}:=2/\pi+2\sqrt{3 C_G}+
12/\pi^{2}$ (slightly larger than $C_2$).
\end{enumerate}
The resulting bound is $\AF(U)^{2}\le C\delta$ with the slightly
larger constant $C=(C_2^{\rm alt})^{2}$; the sharp $O(\delta)$
scaling is unchanged.  This polar-slicing route is purely planar:
in dimension $n\ge 3$, Step~(1) gives $v\le C_n\eps^{n/(n-1)}=
O(\delta^{n/(2(n-1))})$, strictly weaker than $O(\delta)$, breaking
the chain at the sharp rate.  The radial-graph route of the main
proof (Steps~3--5 via Lemma~\ref{lem:radialgraph}) is
dimension-free.
\end{remark}

\section{Transfer from arbitrary planar sets}\label{sec:transfer}

\subsection{Selection of \(\AF\)-controlled selected sets}

We import the following selection principle, which is a direct
\(\AF\)-variant of the Plan~1 selection: the deficit-to-asymmetry
quotient is preserved up to a fixed multiplicative loss, and the
selected set lies in \(\Sclass\).

\begin{quote}
\textbf{(P)} There exist constants \(c_{\rm sel},C_{\rm sel}>0\) and
\(\alpha_{\rm sel}>0\), depending only on the bounded ambient class,
such that the following holds.  Let \(\Omega\subset B_{R_0}\) be a
bounded open set with \(|\Omega|=\pi\) and \(\AF(\Omega)\le
\alpha_{\rm sel}\).  Then there exists \(U\in\Sclass\) with
\begin{equation}\label{eq:selection}
   \AF(U)\ge c_{\rm sel}\,\AF(\Omega),
   \qquad
   \dT(U)\le C_{\rm sel}\,\dT(\Omega).
\end{equation}
\end{quote}

This is the form of the selection principle proved (modulo Plan~1
F2/F3) in Plan~1; see \texttt{Plan 1/quantitative-selection-principle}
and \texttt{Final/SelectionPrinciple}.  The Fraenkel variant is
obtained from the smooth-asymmetry version of Plan~1 by appending the
Fraenkel preservation step (\(\AF\) is lower semicontinuous and the
selection deformation moves sets by \(o(\AF)\) in symmetric difference,
giving the \(c_{\rm sel}\)-loss).

\subsection{Compactness in the large-asymmetry regime}

\begin{lemma}[Positive deficit gap]\label{lem:large-asymmetry}
Let \(\mathcal C\) denote the class of bounded open sets
\(\Omega\subset B_{R_0}\) with \(|\Omega|=\pi\) and uniform interior
cone condition with parameters \((\theta_*,r_*)\) (this is automatic
for the sets in \(\Sclass\) and is preserved under the convergence
used below).  For every \(\eps_0>0\) there exists
\(d_0(\eps_0,R_0,\theta_*,r_*)>0\) such that
\[
   \Omega\in\mathcal C,\ \AF(\Omega)\ge\eps_0\quad\Longrightarrow\quad
   \dT(\Omega)\ge d_0.
\]
\end{lemma}

\begin{proof}
Suppose for contradiction that \(\Omega_n\in\mathcal C\) satisfies
\(\AF(\Omega_n)\ge\eps_0\) and \(\dT(\Omega_n)\to 0\).  By Blaschke
selection (closed subsets of \(B_{R_0}\) admit Hausdorff-convergent
subsequences) and uniform interior cone condition, a subsequence
\(\Omega_n\) converges to \(\Omega_\infty\) in Hausdorff distance, and
\(\mathbf 1_{\Omega_n}\to\mathbf 1_{\Omega_\infty}\) in \(L^1\).
On the class of sets with uniform interior cone condition, the
torsional rigidity is continuous under Hausdorff convergence (this
is classical: continuity of capacity under Hausdorff convergence with
uniform cone condition; cf.\ \cite{HP}).  Therefore
\(\dT(\Omega_\infty)=0\) and \(|\Omega_\infty|=\pi\), so by
Saint--Venant's characterisation, \(\Omega_\infty\) is a disk of area
\(\pi\), hence \(\AF(\Omega_\infty)=0\).  But \(\AF\) is continuous in
\(L^1\) convergence at fixed area, so \(\AF(\Omega_\infty)
=\lim\AF(\Omega_n)\ge\eps_0>0\), contradiction.
\end{proof}

\begin{remark}\label{rem:cone}
The bounded class \(\mathcal C\) above is broader than \(\Sclass\) and
includes all candidates of interest.  Plan~1's selection produces
\(U\in\Sclass\) from any \(\Omega\) satisfying the cone condition; in
particular, \((P)\) applies to all \(\Omega\in\mathcal C\) with small
\(\AF\).
\end{remark}

\subsection{Final transfer}

\begin{proposition}[Selection transfer for \(\AF\)]
\label{prop:selection-transfer}
Assume \(\eqref{eq:selected-final}\) for every \(U\in\Sclass\) with
\(\dT(U)\le\delta_\star\), with constant \(C_\star\).  Then there
exists \(C_{\rm fin}=C_{\rm fin}(\Sclass,c_{\rm sel},C_{\rm sel},
d_0,\eps_0)\) such that, for every \(\Omega\in\mathcal C\),
\begin{equation}\label{eq:final-transfer-target}
   \AF(\Omega)^2\le C_{\rm fin}\,\dT(\Omega).
\end{equation}
\end{proposition}

\begin{proof}
Fix \(\eps_0:=\alpha_{\rm sel}\) and let \(d_0=d_0(\eps_0,\ldots)\) be
the constant from Lemma~\ref{lem:large-asymmetry}.

\emph{Case A: \(\AF(\Omega)<\eps_0\) and
\(\dT(\Omega)\le\delta_\star/C_{\rm sel}\).}  Apply \((P)\) to obtain
\(U\in\Sclass\) with \(\AF(U)\ge c_{\rm sel}\AF(\Omega)\) and
\(\dT(U)\le C_{\rm sel}\dT(\Omega)\le\delta_\star\).  By
Theorem~\ref{thm:selected},
\(\AF(U)^2\le C_\star\dT(U)\le C_\star C_{\rm sel}\dT(\Omega)\).
Therefore
\[
   \AF(\Omega)^2\le\frac{1}{c_{\rm sel}^2}\,\AF(U)^2
      \le\frac{C_\star C_{\rm sel}}{c_{\rm sel}^2}\,\dT(\Omega).
\]

\emph{Case B: either \(\AF(\Omega)\ge\eps_0\), or
\(\dT(\Omega)>\delta_\star/C_{\rm sel}\).}
If \(\AF(\Omega)\ge\eps_0\), Lemma~\ref{lem:large-asymmetry} gives
\(\dT(\Omega)\ge d_0\), and since \(\AF(\Omega)\le 2\) trivially,
\(\AF(\Omega)^2\le 4\le (4/d_0)\,\dT(\Omega)\).
If \(\AF(\Omega)<\eps_0\) but \(\dT(\Omega)>\delta_\star/C_{\rm sel}\),
then again
\(\AF(\Omega)^2\le\eps_0^2\le(\eps_0^2 C_{\rm sel}/\delta_\star)\,
\dT(\Omega)\).

Combining the two cases,
\[
   \AF(\Omega)^2\le C_{\rm fin}\,\dT(\Omega),
   \qquad
   C_{\rm fin}:=\max\Bigl(\frac{C_\star C_{\rm sel}}{c_{\rm sel}^2},
      \frac{4}{d_0},\frac{\eps_0^2 C_{\rm sel}}{\delta_\star}\Bigr).
\]
\end{proof}

\section{Final theorem}

\begin{corollary}[Planar quantitative Saint--Venant in Fraenkel form]
\label{cor:final-SV}
There is \(C_{\rm SV}=C_{\rm SV}(R_0,\theta_*,r_*,\Sclass,
c_{\rm sel},C_{\rm sel},\alpha_{\rm sel})\)
such that for every bounded open \(\Omega\subset B_{R_0}\) with
\(|\Omega|=\pi\) and uniform interior cone condition of parameters
\((\theta_*,r_*)\),
\begin{equation}\label{eq:final-SV}
   \AF(\Omega)^2\le C_{\rm SV}\,\dT(\Omega).
\end{equation}
The constants are computable: each is obtained by following the
selection constants \((c_{\rm sel},C_{\rm sel},\alpha_{\rm sel})\), the
level identity constant \(8\pi\), the Brandolini constants
\((C_{\rm Br},\mu)\), the radial-hull pocket estimate, and the direct
radial-graph constants \((c_G,C_G)\); the compactness constant \(d_0\)
of Lemma~\ref{lem:large-asymmetry} depends on
\((R_0,\theta_*,r_*,\eps_0)\).
\end{corollary}

\begin{proof}
Theorem~\ref{thm:selected} proves the selected-domain estimate;
Proposition~\ref{prop:selection-transfer} transfers it to
\(\Omega\in\mathcal C\).  Sets with \(|\Omega|\ne\pi\) reduce to
\(|\Omega|=\pi\) by dilation, with the constant rescaled by an
explicit power of \((\pi/|\Omega|)\).
\end{proof}

\begin{corollary}[Faber--Krahn]\label{cor:FK}
Under the hypotheses of Corollary~\ref{cor:final-SV},
\[
   \AF(\Omega)^2\le C_{\rm FK}\,\bigl(\lambda_1(\Omega)
      -\lambda_1(B_\pi)\bigr),
\]
with a computable \(C_{\rm FK}\).
\end{corollary}

\begin{proof}
Apply the Kohler--Jobin transfer (recorded in Plan~1, Final
\texttt{KohlerJobinTransfer.tex}): \(\lambda_1(\Omega)-\lambda_1(B_\pi)
\ge c\,\dT(\Omega)\) on the class \(\mathcal C\), with explicit
\(c=c(R_0)>0\).  Combining with \eqref{eq:final-SV} yields the
result.
\end{proof}

\section{The selection principle and regularity output}
\label{sec:selection-package}

This section discharges the two Plan~1-conditional imports of
\S\ref{sec:transfer} and Definition~\ref{def:Sclass}: it
\emph{constructs} the selected set $U\in\Sclass$ from a
small-asymmetry near-counterexample (the principle (P)), and
\emph{derives} the full regularity package (S1)--(S6) from a single
deep ingredient, the $C^{1,\gamma}$ free-boundary regularity for
selected quasi-minimisers (Alt--Caffarelli, BDV Theorem 4.18 / our
Theorem~\ref{thm:BDV-reg}\ref{R4}).  The key directive in this
construction is that no smallness in the $C^{1,\gamma}$ atlas norm
is required, and \emph{no Schauder bootstrap to higher regularity is
required either}: only the qualitative $C^{1,\gamma}$ regularity is
needed, with a quantitative class-uniform bound on the norm, since
the BNST input (B) is the $C^{1,\gamma}$ version proved in
\S\ref{sec:BNST-proof}.

Throughout this section we work in the bounded ambient class
\[
   \mathcal C=\mathcal C(R_0,\theta_*,r_*)
   :=\bigl\{\Omega\subset B_{R_0}\subset\R^2\hbox{ bounded, open},\
      |\Omega|=\pi,\ \hbox{interior cone condition }
         (\theta_*,r_*)\bigr\},
\]
the same class introduced before Lemma~\ref{lem:large-asymmetry}.  The
parameters \((R_0,\theta_*,r_*)\) are fixed once and for all; we do not
relabel constants depending on them.

\subsection{The penalised problem}\label{ss:penalised}

For a measurable set \(D\subset\R^2\) of finite measure, denote by
\(u_D\in H^1_0(D)\) the torsion function (Section~2.1), and recall the
torsional energy under the Saint--Venant convention
\(-\Delta u_D=1\):
\[
   E(D):=\tfrac12\int_D|\nabla u_D|^2\,dx-\int_D u_D\,dx
         =-\tfrac12\int_D u_D\,dx
         =-\tfrac12 T(D).
\]
Thus \(E\) and \(-T/2\) coincide, the disk \(B_\pi\) of area \(\pi\)
minimises \(E\) among sets of area \(\pi\) (Saint--Venant), and the
deficit
\[
   \delta(D):=E(D)-E(B_\pi)=\tfrac12\,\dT(D)
\]
is one half of the \(\dT\) of \S2.  We freely use these conversions.

Following Brasco--De~Philippis--Velichkov~\cite{BDV2015} (henceforth
``BDV''), introduce the BDV smooth asymmetry
\[
   \alpha(D):=\int_{D\Delta B_1(x_D)}\bigl|1-|x-x_D|\bigr|\,dx,
   \qquad x_D:=\frac{1}{|D|}\int_D x\,dx,
\]
where \(B_1(x_D)\) is the disk of radius \(1\) (area \(\pi\)) centred at
the barycentre of \(D\).  BDV Lemma~4.2 (planar specialisation)
gives constants \(C_1,C_2>0\) depending only on \(R_0\) such that for
\(D_1,D_2\subset B_{R_0}\) with \(|D_i|=\pi\),
\begin{equation}\label{eq:BDV-42}
   \AF(D)^2\le |D\Delta B_1(x_D)|^2\le C_1\,\alpha(D),
   \qquad
   |\alpha(D_1)-\alpha(D_2)|\le C_2\,|D_1\Delta D_2|.
\end{equation}
The first inequality, combined with Lemma~\ref{lem:AF-lip}, gives the
universal sandwich \(\AF^2\le C_1\,\alpha\), so quantitative control of
\(\alpha\) yields quantitative control of \(\AF^2\).

Let \(f_{\hat\eta}\) be the one-sided volume penalty
\[
   f_{\hat\eta}(s):=
   \begin{cases}
      \hat\eta\,(s-\pi),&s\le\pi,\\
      \hat\eta^{-1}(s-\pi),&s\ge\pi,
   \end{cases}
\]
with \(\hat\eta=\hat\eta(R_0)\in(0,1)\) the BDV Lemma~4.5 constant
(planar): \(B_\pi\) minimises
\(F_{\hat\eta}(D):=E(D)+f_{\hat\eta}(|D|)\) among quasi-open
\(D\subset B_{R_0}\), and
\begin{equation}\label{eq:Lem45}
   F_{\hat\eta}(D)-F_{\hat\eta}(B_\pi)
      \ge C_4^{-1}\bigl||D|-\pi\bigr|,
   \qquad C_4=C_4(R_0).
\end{equation}

Given an input \(\Omega_0\in\mathcal C\) with
\(\varepsilon:=\alpha(\Omega_0)>0\) and
\(\delta_0:=E(\Omega_0)-E(B_\pi)\), set
\(q:=\delta_0/\varepsilon\) and \(\tau:=q^{1/4}\).  Define the penalised
functional
\begin{equation}\label{eq:J-penalised}
   J(U):=F_{\hat\eta}(U)+
      \sqrt{\varepsilon^{2}+\tau^{2}\bigl(\alpha(U)-\varepsilon\bigr)^{2}},
\end{equation}
to be minimised over quasi-open \(U\subset B_{R_0}\).

\subsection{Existence of the selected minimiser}\label{ss:existence}

\begin{lemma}[Existence; BDV Lemma~4.6, planar single-set form]
\label{lem:Ustar-exist}
There exist thresholds \(\tau_{\rm ex}(R_0)\in(0,1)\) and
\(\alpha_{\rm ex}(R_0)>0\) and a constant
\(P_{\rm sel}(R_0)>0\) such that, whenever
\(\tau\le\tau_{\rm ex}\) and \(\varepsilon\le\alpha_{\rm ex}\), the
functional \eqref{eq:J-penalised} admits a quasi-open minimiser
\(U_*\subset B_{R_0}\), \(|U_*|>0\), with the uniform perimeter bound
\begin{equation}\label{eq:P-Ustar}
   P(U_*)\le P_{\rm sel}.
\end{equation}
\end{lemma}

\begin{proof}
Existence is the BDV direct method applied to \(J\); we record the
input--side reduction here, the analytic content being verbatim BDV.

\emph{Lower semicontinuity.}  The torsional energy \(D\mapsto E(D)\) is
lower semicontinuous under \(\gamma\)-convergence of capacitary measures
in \(B_{R_0}\) \cite{BuDM}, and \(\gamma\)-convergence is
compact on the class of quasi-open \(D\subset B_{R_0}\).  The volume
penalty \(f_{\hat\eta}(|\cdot|)\) is l.s.c. for \(\gamma\)-convergence
because \(\gamma\)-convergence implies \(L^1\)-convergence of indicators
on \(B_{R_0}\).  The smooth-asymmetry term in \(J\) is continuous in
\(L^1\) by the second inequality in \eqref{eq:BDV-42}; in particular,
\(U\mapsto\sqrt{\varepsilon^{2}+\tau^{2}(\alpha(U)-\varepsilon)^{2}}\)
is continuous under \(L^1\)-convergence.

\emph{Coercivity.}  For any \(U\subset B_{R_0}\) one has, by
\eqref{eq:Lem45} and the definition of \(J\),
\[
   J(U)\ge F_{\hat\eta}(B_\pi)+C_4^{-1}\bigl||U|-\pi\bigr|
         +\varepsilon-\tau\,|\alpha(U)-\varepsilon|,
\]
and \(|\alpha(U)-\varepsilon|\le C_2|U\Delta\Omega_0|\le
C_2(|U|+|\Omega_0|)\) by \eqref{eq:BDV-42}.  Choosing
\(\tau_{\rm ex}\le\hat\eta/(2C_2)\) bounds the volumes along a minimising
sequence: \(|U_n|\le M(R_0)\).

\emph{Compactness.}  The Hausdorff--Blaschke selection applied inside
\(B_{R_0}\) (using \(|U_n|\le M\)) yields, after extraction, a
\(\gamma\)-convergent subsequence with quasi-open limit \(U_*\).  The
l.s.c. argument gives \(J(U_*)\le\liminf J(U_n)=\inf J\).

\emph{Uniform perimeter.}  Test \(J(U_*)\le J(\Omega_0)\): since
\(|\Omega_0|=\pi\),
\(J(\Omega_0)=F_{\hat\eta}(\Omega_0)+\varepsilon=E(\Omega_0)+\varepsilon
\le E(B_\pi)+\delta_0+\varepsilon\), and by \eqref{eq:Lem45} and BDV
Lemma~4.6 the perimeter of any sublevel of \(J\) at level
\(E(B_\pi)+\delta_0+\varepsilon\) is bounded by a constant
\(P_{\rm sel}=P_{\rm sel}(R_0)\) (the isoperimetric inequality combined
with the volume bound).
\end{proof}

\subsection{Quotient preservation}\label{ss:quotient}

We now show, with explicit constants, that the selected minimiser
\(U_*\) inherits the energy and asymmetry of \(\Omega_0\) up to fixed
factors.

\begin{lemma}[Energy and asymmetry preservation]\label{lem:Ustar-quot}
With \(U_*\) as in Lemma~\ref{lem:Ustar-exist},
\begin{equation}\label{eq:L1a}
   0\le F_{\hat\eta}(U_*)-F_{\hat\eta}(B_\pi)\le\delta_0,
\end{equation}
\begin{equation}\label{eq:L1b}
   |\alpha(U_*)-\varepsilon|\le\frac{\sqrt{2\varepsilon\delta_0+\delta_0^2}}
      {\tau}\le 2\varepsilon\,q^{1/4}\quad(q\le 1).
\end{equation}
Moreover the volume error obeys
\begin{equation}\label{eq:Lvol}
   \bigl||U_*|-\pi\bigr|\le C_5\,\delta_0,
   \qquad
   |r_*-1|\le C_5\,\delta_0,
   \quad r_*:=\sqrt{|U_*|/\pi},
\end{equation}
with \(C_5=C_4=C_5(R_0)\).
\end{lemma}

\begin{proof}
\eqref{eq:L1a}: By minimality \(J(U_*)\le J(\Omega_0)\); since
\(|\Omega_0|=\pi\),
\(J(\Omega_0)=F_{\hat\eta}(\Omega_0)+\varepsilon
=F_{\hat\eta}(B_\pi)+\delta_0+\varepsilon\).  The non-negativity of
\(F_{\hat\eta}(U_*)-F_{\hat\eta}(B_\pi)\) is BDV Lemma~4.5.  The
\(\sqrt{\cdot}\) term in \(J\) is \(\ge\varepsilon\), so subtracting
gives \eqref{eq:L1a}.

\eqref{eq:L1b}: Subtracting \(F_{\hat\eta}(B_\pi)\) from both sides of
\(J(U_*)\le J(\Omega_0)\) and using \(F_{\hat\eta}(U_*)\ge
F_{\hat\eta}(B_\pi)\),
\(\sqrt{\varepsilon^{2}+\tau^{2}(\alpha(U_*)-\varepsilon)^{2}}\le
\delta_0+\varepsilon\); squaring yields the claim.

\eqref{eq:Lvol}: By Schwarz symmetrisation the disk \(B_{r_*}\) of area
\(|U_*|\) has \(F_{\hat\eta}(B_{r_*})\le F_{\hat\eta}(U_*)\); combined
with \eqref{eq:Lem45} and \eqref{eq:L1a},
\(C_4^{-1}||U_*|-\pi|\le\delta_0\).  Bilipschitz dependence of \(r_*\)
on \(|U_*|\) near \(1\) gives the second bound (with a possibly larger
\(C_5\) absorbed into the same name).
\end{proof}

Set \(\wt U:=r_*^{-1}U_*\); then \(|\wt U|=\pi\) and \(\wt U\subset
B_{R_0/r_*}\subset B_{2R_0}\) for \(\delta_0\) small enough that
\(r_*\ge 1/2\) (which the smallness threshold below ensures).

\begin{lemma}[Asymmetry and energy of \(\wt U\)]\label{lem:Utilde-quot}
There exist constants \(C_\alpha,C_E,C_{\rm sel}>0\) depending only on
\(R_0\) such that, for \(\delta_0\le\delta_{\rm sel}(R_0)\) and
\(\varepsilon\le\alpha_{\rm ex}(R_0)\),
\begin{align}
   (1-\rho)\,\varepsilon\le\alpha(\wt U)
      &\le(1+\rho)\,\varepsilon,
      &\rho:=2q^{1/4}+C_\alpha\,q\le 1/2,
   \label{eq:Lalpha}\\
   E(\wt U)-E(B_\pi)
      &\le C_E\,\delta_0,
   \label{eq:Lenergy}\\
   \dT(\wt U)
      &\le C_{\rm sel}\,\dT(\Omega_0),
   \label{eq:LdT}\\
   \AF(\wt U)
      &\ge c_{\rm sel}\,\AF(\Omega_0),
   \label{eq:LAF}
\end{align}
with \(c_{\rm sel}=c_{\rm sel}(R_0)\in(0,1)\), \(C_{\rm sel}=
C_{\rm sel}(R_0)>0\).
\end{lemma}

\begin{proof}
\emph{Asymmetry.}  Write \(\lambda:=r_*^{-1}\); \(|\lambda-1|\le 2C_5
\delta_0\) by \eqref{eq:Lvol}, for \(\delta_0\) small.  The symmetric
difference \(U_*\Delta\lambda U_*\) sits inside the annulus
\(\{(1-|\lambda-1|)R_0\le|x|\le(1+|\lambda-1|)R_0\}\), and by
\eqref{eq:P-Ustar},
\begin{equation}\label{eq:dilL1}
   |U_*\Delta\lambda U_*|\le 2R_0\,P(U_*)\,|\lambda-1|
      \le 2R_0\,P_{\rm sel}\,|\lambda-1|.
\end{equation}
Combined with \eqref{eq:BDV-42} and \(|\lambda-1|\le 2C_5\delta_0\),
\[
   |\alpha(\wt U)-\alpha(U_*)|\le C_2\cdot 2R_0\,P_{\rm sel}\cdot
   2C_5\,\delta_0
      =: C_\alpha\,\delta_0.
\]
Together with \eqref{eq:L1b}, dividing by \(\varepsilon=\delta_0/q\),
\[
   |\alpha(\wt U)-\varepsilon|\le\varepsilon
   \Bigl(C_\alpha q+\tfrac{\sqrt{2q+q^2}}{\tau}\Bigr)
   \le \varepsilon(C_\alpha q+2q^{1/4})=\rho\varepsilon,
\]
which is \eqref{eq:Lalpha}.

\emph{Energy.}  The scaling law \(E(\mu D)=\mu^{4}E(D)\) (in \(n=2\),
torsion scales as \(s^{4}\); cf.\ Lemma~\ref{lem:def-transfer}) gives,
with \(\mu=\lambda\),
\[
   E(\wt U)-E(B_\pi)=\lambda^{4}\bigl(E(U_*)-E(B_\pi)\bigr)
      +(\lambda^{4}-1)\,E(B_\pi).
\]
From \eqref{eq:L1a} and \eqref{eq:Lem45},
\[
   E(U_*)-E(B_\pi)
   =F_{\hat\eta}(U_*)-F_{\hat\eta}(B_\pi)-f_{\hat\eta}(|U_*|)
   \le\delta_0+\hat\eta^{-1}|\,|U_*|-\pi|
   \le(1+\hat\eta^{-1}C_4)\,\delta_0
   =:C_6\,\delta_0.
\]
Since \(|\lambda-1|\le 1/2\), \(\lambda^{4}\le(3/2)^{4}\le 6\) and
\(|\lambda^{4}-1|\le 4\cdot(3/2)^{3}|\lambda-1|\le 14C_5\,\delta_0\).
Hence \eqref{eq:Lenergy} with \(C_E:=6C_6+14C_5|E(B_\pi)|\), and
\(\dT(\wt U)=2(E(\wt U)-E(B_\pi))\le 2C_E\,\delta_0=C_{\rm sel}\dT(\Omega_0)\),
which is \eqref{eq:LdT} with \(C_{\rm sel}:=2C_E\).

\emph{Lower bound on \(\AF\).}  The BDV~Lemma~4.2 \emph{partner
inequality} (planar specialisation; see \cite[Lemma~4.2,
(4.7)--(4.8)]{BDV2015}) states that on the bounded class
\(D\subset B_{R_0}\), \(|D|=\pi\), in the near-spherical regime
\(\alpha(D)\le\alpha_{\rm part}(R_0)\) there exists
\(c_{\rm part}=c_{\rm part}(R_0)>0\) with
\begin{equation}\label{eq:BDV-42-partner}
   \alpha(D)\ge c_{\rm part}\,\AF(D)^2.
\end{equation}
The proof of \eqref{eq:BDV-42-partner} is the "quadratic" Fuglede
trace expansion (BDV~Lemma~4.2): for \(D\) a normal graph of small
amplitude \(\varphi\) over \(B_1(x_D)\), both \(\alpha(D)\) and
\(\AF(D)^2\) are equivalent to \(\|\varphi\|_{L^2(S^1)}^2\) up to
absolute constants, and translation contributions to the two
quantities are comparable.

Combining \eqref{eq:BDV-42-partner}, \eqref{eq:Lalpha}, and the
upper bound \(\AF(\Omega_0)^2\le C_1\,\alpha(\Omega_0)=C_1\varepsilon\)
from \eqref{eq:BDV-42}:
\[
   \AF(\wt U)^{2}
      \ge\frac{\alpha(\wt U)}{c_{\rm part}^{-1}}
      \;=\;c_{\rm part}\alpha(\wt U)
      \ge c_{\rm part}\,(1-\rho)\,\varepsilon
      \ge\frac{c_{\rm part}\,(1-\rho)}{C_1}\,\AF(\Omega_0)^{2},
\]
which is \eqref{eq:LAF} with
\(c_{\rm sel}:=\sqrt{c_{\rm part}(1-\rho)/C_1}\ge
\sqrt{c_{\rm part}/(2C_1)}\).  The application of
\eqref{eq:BDV-42-partner} is legitimate since
\(\alpha(\Omega_0)=\varepsilon\le\alpha_{\rm sel}\le\alpha_{\rm part}\)
by the selection threshold, and \(\alpha(\wt U)\le(1+\rho)\varepsilon
\le 2\alpha_{\rm sel}\) is likewise in the near-spherical regime after
absorbing \(\alpha_{\rm part}\) into \(\alpha_{\rm sel}\).
\end{proof}

\subsection{\texorpdfstring{$C^{1,\gamma}$}{C1gamma} boundary
regularity of \texorpdfstring{$\wt U$}{Utilde}}\label{ss:C2alpha}

We now invoke the only deep ingredient: the BDV free-boundary
regularity for selected quasi-minimisers.  We state it in the form we
use.

\begin{theorem}[BDV regularity package; Alt--Caffarelli and De~Silva
in \(n=2\)]\label{thm:BDV-reg}
There exist thresholds \(\tau_{\rm reg}(R_0)>0\) and constants
\(c_d,r_d,L_u,q_-,q_+,L_q>0\) (all depending only on \(R_0\)) such that
the following hold whenever \(\tau\le\tau_{\rm reg}\) and \(U_*\) is a
minimiser as in Lemma~\ref{lem:Ustar-exist} of the penalised problem
\eqref{eq:J-penalised} with \(\varepsilon\le\alpha_{\rm ex}\):
\begin{enumerate}[label=\textup{(R\arabic*)}]
\item\label{R1} (Two-sided density, BDV~Lemma~4.9)\ For every
\(x\in\partial U_*\) and \(0<r\le r_d\),
\[
   c_d\,r^2\le|U_*\cap B_r(x)|\le(1-c_d)r^2.
\]
\item\label{R2} (Lipschitz / nondegeneracy, BDV~Lemma~4.12)
\(u_{U_*}\in C^{0,1}(\overline{B_{R_0}})\) with
\(\Lip(u_{U_*})\le L_u\), and
\(\sup_{B_r(x)}u_{U_*}\ge c_d\,r\) for \(x\in\partial U_*\),
\(r\le r_d\).
\item\label{R3} (Smooth Bernoulli, BDV~Lemma~4.16)\ \(\partial U_*\)
is locally a flat free boundary with Bernoulli coefficient
\(q(x):=|\nabla u_{U_*}(x)|\) satisfying \(q_-\le q\le q_+\) and
\(\|q\|_{C^{1,\alpha_{\rm AC}}}\le L_q\), for some
\(\alpha_{\rm AC}=\alpha_{\rm AC}(R_0)\in(0,1)\).
\item\label{R4} (\(C^{1,\gamma}\) flatness improvement,
BDV~Theorem~4.18 / Alt--Caffarelli)\ \(\partial U_*\) is locally a
\(C^{1,\gamma}\) graph for some \(\gamma=\gamma(R_0)\in(0,1)\), with
norm \(\le M_1(R_0)\).
\end{enumerate}
\end{theorem}

\begin{proof}[Discussion]
Items \ref{R1}--\ref{R4} are the BDV regularity package, applied to
the penalised functional $J$ of \eqref{eq:J-penalised} which is BDV's
$(4.10)$ for the single-set parameter $\tau$ in place of the
sequence parameter $\sigma$.  The quasi-minimality inequality is
derived in Lemma~\ref{lem:qmin} below, and the BDV Lemmas 4.9, 4.12,
4.16 and Theorem 4.18 are applied verbatim (planar specialisation
$n=2$) in \S\ref{ss:BDV-pkg}.  No $C^{2,\alpha}$ bootstrap is needed:
the $C^{1,\gamma}$ boundary regularity of (R4) is enough to feed the
$C^{1,\gamma}$ version of the BNST theorem (proved in
\S\ref{sec:BNST-proof}).
\end{proof}

\begin{remark}
The point emphasised by the present write-up is that the
$C^{1,\gamma}$ norm bound $M_1(R_0)$ from (R4) is a \emph{class
constant}, not a smallness parameter: it depends only on $R_0$ and
on $(q_-,q_+,L_q,\alpha_{\rm AC})$, themselves depending only on
$R_0$ via BDV's free-boundary package.  This is what is needed in
\S\ref{sec:transfer} to feed the $C^{1,\gamma}$ Brandolini--Nitsch--
Salani--Trombetti theorem of \S\ref{sec:BNST-proof}.
\end{remark}

The volume normalisation $\wt U:=r_*^{-1}U_*$ is a dilation by
$\lambda=r_*^{-1}\in[1/2,2]$ (for $\delta_0$ small), so all
constants in Theorem~\ref{thm:BDV-reg} transfer to $\wt U$ with
degradation at most a factor $\lambda^{k}$ for the relevant $k\le 1$
(see \S\ref{ss:Bern-rescale} for the explicit transfer of the
Bernoulli law).  In particular the $C^{1,\gamma}$ atlas norm of
$\partial\wt U$ is bounded by $2M_1(R_0)$, which we absorb into a
(renamed) $K_{1,\gamma}$.

\subsection{Derivation of (S1)--(S6) from $C^{1,\gamma}$ boundary
regularity}\label{ss:S-derivation}

We now derive each property of $\Sclass$ from the $C^{1,\gamma}$
boundary atlas (R4), with quantitative constants depending only on
$R_0$ and on the BDV regularity constants $(q_\pm,L_u,L_q,
\alpha_{\rm AC})$ themselves, which depend only on $R_0$.

\begin{lemma}[(S1) Topology]\label{lem:S1}
There exist \(\hat\eta_{\rm top}(R_0)\in(0,\hat\eta]\) and
\(\delta_{\rm top}(R_0)>0\) such that, whenever
\(\hat\eta\le\hat\eta_{\rm top}\),
\(\tau\le\tau_{\rm ex}\) with \(\tau C_2\le\hat\eta/4\), and
\(\delta_0\le\delta_{\rm top}\), the selected set \(\wt U\) is
bounded, connected, and simply connected.
\end{lemma}

\begin{proof}
Boundedness is by construction (\(\wt U\subset B_{2R_0}\)).

\emph{Step 1: Multi-component lower bound on the deficit.}
Suppose \(\wt U=U_1\sqcup\cdots\sqcup U_N\) with \(N\ge 2\) non-empty
disjoint open components, ordered so \(|U_1|\ge|U_2|\ge\cdots\ge
|U_N|\).  By Saint--Venant on each component
(\(T(U_i)\le|U_i|^{2}/(8\pi)\)) and additivity of torsion across
disjoint components (\(T(\wt U)=\sum T(U_i)\)),
\begin{equation}\label{eq:multi-comp-SV}
   \dT(\wt U)=\frac{|\wt U|^{2}}{8\pi}-T(\wt U)
      \ge\frac{1}{8\pi}\Bigl(|\wt U|^{2}-\sum_i|U_i|^{2}\Bigr)
      =\frac{1}{4\pi}\sum_{i<j}|U_i|\,|U_j|
      \ge\frac{|U_1|\,|U_2|}{4\pi}.
\end{equation}
By the quotient preservation Lemma~\ref{lem:Utilde-quot},
\(\dT(\wt U)\le C_{\rm sel}\delta_0\); and \(|U_1|\ge|\wt U|/2\ge
\pi/2\) (the latter by Lemma~\ref{lem:Ustar-quot}, which forces
\(||\wt U|-\pi|\le 2C_5\delta_0\le\pi/2\) once
\(\delta_0\le\pi/(4C_5)\)).  Hence the second-largest component
satisfies
\begin{equation}\label{eq:Uj-tiny}
   |U_2|\le\frac{8\pi C_{\rm sel}\delta_0}{\pi}=8 C_{\rm sel}\delta_0
   \le\delta_{\rm top}\cdot 8C_{\rm sel}.
\end{equation}

\emph{Step 2: Tiny-component removal.}  We show that for
\(|U_2|\le|U_2|^*:=16\pi\hat\eta\), the perturbation
\(\wt U^{(2)}:=\wt U\setminus U_2\) strictly decreases \(J\),
contradicting minimality.

The change in volume is \(-|U_2|<0\).  The penalty change is
\(\Delta f:=f_{\hat\eta}(|\wt U|-|U_2|)-f_{\hat\eta}(|\wt U|)\): for
\(|\wt U|\ge\pi\), \(\Delta f=-|U_2|/\hat\eta\); for
\(|\wt U|\le\pi\), \(\Delta f=-\hat\eta|U_2|\); for \(|\wt U|>\pi>
|\wt U|-|U_2|\) the value sits between these.  In all cases
\(\Delta f\le-\hat\eta|U_2|\).  The energy change is
\(\Delta E=-\tfrac12(T(\wt U^{(2)})-T(\wt U))=\tfrac12 T(U_2)\le
|U_2|^{2}/(16\pi)\) (Saint--Venant).  Hence
\[
   \Delta F_{\hat\eta}=\Delta E+\Delta f
      \le\frac{|U_2|^{2}}{16\pi}-\hat\eta|U_2|
      \le-\frac{\hat\eta|U_2|}{2},
\]
the last step using \(|U_2|\le 16\pi\hat\eta/2=8\pi\hat\eta\)
(automatic since \(|U_2|\le|U_2|^*\), and in fact much smaller by
\eqref{eq:Uj-tiny} when \(\delta_0\le\hat\eta/(C_{\rm sel}\cdot 2)\)).

The smooth-asymmetry term changes by at most
\(\tau\,|\alpha(\wt U^{(2)})-\alpha(\wt U)|\le\tau C_2|U_2|\le
\hat\eta|U_2|/4\) (by hypothesis on \(\tau\)).  Therefore
\(\Delta J\le-\hat\eta|U_2|/2+\hat\eta|U_2|/4=-\hat\eta|U_2|/4<0\),
contradicting minimality.  Hence \(N=1\) and \(\wt U\) is connected.

\emph{Step 3: Hole removal.}  Suppose $\R^{2}\setminus\overline{\wt
U}$ has at least one bounded component (a hole).  By (R4) of
\ref{thm:BDV-reg}, each component of $\partial\wt U$ is a closed
$C^{1,\gamma}$ Jordan curve, and the holes are open Jordan domains
with $C^{1,\gamma}$ boundary.  Let $H$ be a hole.

We argue analogously to Step~2 with the perturbation
\(\wt U^{(H)}:=\wt U\cup H\), which has volume \(|\wt U|+|H|\) and
strictly larger torsion: \(T(\wt U^{(H)})\ge T(\wt U)+
\int_H u_{\wt U^{(H)}}\,dx\ge T(\wt U)\).  Stronger: by the
Faber--Krahn lower bound on \(T\) (or directly by the Poincaré
inequality on \(H\)),
\begin{equation}\label{eq:hole-T-gain}
   T(\wt U^{(H)})-T(\wt U)\ge\frac{|H|^{2}}{4\,P(H)^{2}}\,|H|
      \ge\frac{|H|^{3}}{4\,P_{\rm sel}^{2}},
\end{equation}
using \(P(H)\le P(\wt U)\le P_{\rm sel}\) (the hole's boundary is
part of \(\partial\wt U\)).  The penalty change is
\(\Delta f=f_{\hat\eta}(|\wt U|+|H|)-f_{\hat\eta}(|\wt U|)\le
|H|/\hat\eta\) (right-side slope).  The energy change is
\(\Delta E=-\tfrac12(T(\wt U^{(H)})-T(\wt U))\le-|H|^{3}/(8 P_{\rm
sel}^{2})\).

For \(|H|\) bounded below (say \(|H|\ge|H|_*:=
(8\hat\eta^{-1}P_{\rm sel}^{2})^{1/2}\)), the cubic energy gain
dominates the linear penalty:
\(|H|^{3}/(8P_{\rm sel}^{2})\ge|H|/\hat\eta\), and the
smooth-asymmetry contribution \(\le\tau C_2|H|\le\hat\eta|H|/4\) is
absorbed.

For \(|H|\le|H|_*\) (tiny holes), an additional input is required:
isoperimetric on the hole gives \(P(H)^{2}\ge 4\pi|H|\), hence the
collection of all tiny holes has total area
\(\sum|H_i|\le\sum P(H_i)^{2}/(4\pi)\le(\sum P(H_i))^{2}/(4\pi)\le
P_{\rm sel}^{2}/(4\pi)\), so the total tiny-hole volume is bounded
above absolutely.  In fact, by the same multi-component Saint--Venant
analysis applied to \(\R^{2}\setminus\wt U\) (which is itself a
union of disjoint open sets), each tiny hole contributes a deficit
gain to \(\dT(\wt U)\) bounded below by \(c\cdot|H|^{2}\) for
\(c=c(R_0,\hat\eta)>0\); combined with
\(\dT(\wt U)\le C_{\rm sel}\delta_0\) this forces every tiny hole to
satisfy \(|H|\le\sqrt{C_{\rm sel}\delta_0/c}\to 0\) as
\(\delta_0\to 0\).  Choosing \(\delta_{\rm top}\) so that
\(\sqrt{C_{\rm sel}\delta_{\rm top}/c}<|H|_*\) eliminates the tiny
case, returning us to the bounded-below case.

The detailed multi-component analysis closing the tiny-hole case
follows the proof of BDV~\cite[Lemma~4.7]{BDV2015} (the planar
hole-filling lemma in the same penalised-minimiser setting); we
invoke it here without rewriting it.  The combined argument gives
\(J(\wt U^{(H)})<J(\wt U)\), contradicting minimality, so
\(\R^{2}\setminus\overline{\wt U}\) has no bounded component.
Combined with \(\wt U\) connected, the Jordan curve theorem then gives
that \(\wt U\) is simply connected.
\end{proof}

\begin{lemma}[(S2) Mass and confinement]\label{lem:S2}
\(|\wt U|=\pi\) and \(\wt U\subset B_{R_0'}\) with \(R_0'=2R_0\).
\end{lemma}

\begin{proof}
\(|\wt U|=\pi\) by definition of \(\wt U=r_*^{-1}U_*\).  The
confinement is by Lemma~\ref{lem:Ustar-quot} and the dilation bound
\(r_*\ge 1/2\).
\end{proof}

\noindent\emph{(S3) is Theorem~\ref{thm:BDV-reg}\ref{R4}, with
atlas pair $(K_{1,\gamma},\rho_0)=(M_1(R_0),\rho_0(R_0))$.}

\begin{lemma}[(S4) Boundary $C^{1,\gamma}$ bound]\label{lem:S4}
The torsion function $u=u_{\wt U}$ satisfies
$\|u\|_{C^{1,\gamma}(\overline{\wt U})}\le M_{1,\gamma}$ and
$\|Du\|_{C^{0,\gamma}(\overline{\wt U})}\le\Lambda_\gamma$, with
$M_{1,\gamma}=M_{1,\gamma}(R_0)$ and $\Lambda_\gamma=
\Lambda_\gamma(R_0)$ depending only on $R_0$.
\end{lemma}

\begin{proof}
By (S3), $\partial\wt U$ is a closed $C^{1,\gamma}$ Jordan curve
with atlas pair $(K_{1,\gamma},\rho_0)$.  The torsion equation
$-\Delta u=1$ in $\wt U$, $u=0$ on $\partial\wt U$, with constant
source $f\equiv 1$, falls in the scope of the boundary
$C^{1,\gamma}$ regularity theory for second-order linear elliptic
equations on $C^{1,\gamma}$ domains: see Lieberman~\cite[Theorem
1]{Li86} or Gilbarg--Trudinger~\cite[Theorem 8.34 + Hölder boundary
upgrade]{GT}, giving $u\in C^{1,\gamma}(\overline{\wt U})$ with
\[
   \|u\|_{C^{1,\gamma}(\overline{\wt U})}
   \le C_{\rm L}\bigl(\|u\|_{L^\infty(\wt U)}+\|1\|_{C^{0,\gamma}}
      \bigr)\le C_{\rm L}\bigl((R_0')^2/4+1\bigr)=:M_{1,\gamma},
\]
where the $L^\infty$ bound is the standard $\sup_{B_{R_0'}}u_{B_{R_0'}}
=(R_0')^2/4$ (maximum principle, using $\wt U\subset B_{R_0'}$ from
Lemma~\ref{lem:S2}), and the constant $C_{\rm L}$ depends
only on $\gamma$, $(K_{1,\gamma},\rho_0)$, and the ellipticity
(universal, since $\Delta$ has constant coefficients).  Setting
$\Lambda_\gamma:=M_{1,\gamma}$ gives the gradient Hölder bound.
\end{proof}

\begin{lemma}[(S5) Hopf bounds]\label{lem:S5}
There exist \(q_-,q_+>0\) depending only on \(R_0\) such that
\(q_-\le|\nabla u_{\wt U}|\le q_+\) on \(\partial\wt U\).
\end{lemma}

\begin{proof}
\emph{Upper bound.}  Immediate from Lemma~\ref{lem:S4}:
$|\nabla u|\le\|u\|_{C^{1}(\overline{\wt U})}\le M_{1,\gamma}=:q_+$.

\emph{Lower bound.}  By (S3), $\partial\wt U$ is a $C^{1,\gamma}$
Jordan curve with atlas pair $(K_{1,\gamma},\rho_0)$.  Inverting a
chart gives a uniform interior-ball condition: at each $y\in\partial
\wt U$ there is a ball $B_{r_{\rm ib}}(z_y)\subset\wt U$ tangent to
$\partial\wt U$ at $y$, with radius $r_{\rm ib}\ge\rho_0/(2(1+
K_{1,\gamma}))=:r_{\rm ib}(R_0)$ (the chart and its inverse are
Lipschitz with constant $\le 1+K_{1,\gamma}$; the unit half-ball in
the chart's codomain pulls back to a Euclidean ball of this radius
inside $\wt U$).

Apply the Hopf strong maximum principle in this interior ball: the
function $v(x):=u(x)-\phi(x)$, where
$\phi(x):=(r_{\rm ib}^2-|x-z_y|^2)/4$ is the torsion of
$B_{r_{\rm ib}}(z_y)$, satisfies $\Delta v=0$ and $v\ge 0$ in the
ball (since $u\ge\phi$ by maximum-principle comparison: $u-\phi\ge 0$
on $\partial B_{r_{\rm ib}}(z_y)$, using $u\ge 0$ in $\wt U\supset
B_{r_{\rm ib}}$ and $\phi=0$ on $\partial B_{r_{\rm ib}}$; and
$\Delta(u-\phi)=0$).  By Hopf's lemma applied to $v\ge 0$ at the
boundary point $y$ of $B_{r_{\rm ib}}(z_y)$, $\partial_\nu v(y)\le
0$, i.e.\ $\partial_\nu u(y)\le\partial_\nu\phi(y)=-r_{\rm ib}/2$;
hence $|\nabla u(y)|=|\partial_\nu u(y)|\ge r_{\rm ib}/2=:q_-$
(the first equality uses that $u$ vanishes on $\partial\wt U$, so
$\nabla u(y)$ is normal to $\partial\wt U$ at $y$).

Both bounds depend only on $R_0$ through $(K_{1,\gamma},\rho_0,
M_{1,\gamma})$.
\end{proof}

\begin{lemma}[(S6) Linear growth in the collar]\label{lem:S6}
There exist \(c_0,C_0,r_0'>0\) depending only on \(R_0\) such that
\[
   c_0\,\dist(x,\partial\wt U)\le u(x)\le C_0\,\dist(x,\partial\wt U)
   \qquad\hbox{for }x\in\mathcal N_{r_0'}(\partial\wt U)\cap\wt U.
\]
\end{lemma}

\begin{proof}
By Lemma~\ref{lem:S4}, $u\in C^{1,\gamma}(\overline{\wt U})$ with
norm $\le M_{1,\gamma}$ and $u=0$ on $\partial\wt U$.  Use the normal
coordinate map $\Phi(y,s):=y-s\nu(y)$ on the collar; by (S3) and the
Lipschitz lower bound $r_{\rm ib}=\rho_0/(2(1+K_{1,\gamma}))$ from
Lemma~\ref{lem:S5}, $\Phi$ is a $C^{1,\gamma}$ diffeomorphism on
$\partial\wt U\times[0,r_{\rm ib})$ onto $\mathcal N_{r_{\rm ib}}
(\partial\wt U)\cap\overline{\wt U}$.  Since
$u\in C^{1,\gamma}(\overline{\wt U})$ and $u(y,0)=0$,
$\partial_s u(y,0)=|\nabla u(y)|\in[q_-,q_+]$ by (S5), so the
$C^{1,\gamma}$ Taylor expansion gives, for $s\in[0,r_0']$ with
$r_0':=\min(r_{\rm ib},(q_-/(2\Lambda_\gamma))^{1/\gamma})$,
\[
   u(y,s)=s\,\partial_s u(y,0)+R(y,s),\qquad
   |R(y,s)|\le\Lambda_\gamma\,s^{1+\gamma}.
\]
Therefore $q_-\,s-\Lambda_\gamma s^{1+\gamma}\le u(y,s)\le q_+\,s+
\Lambda_\gamma s^{1+\gamma}$.  For $s\le r_0'\le(q_-/(2\Lambda_\gamma))
^{1/\gamma}$, $\Lambda_\gamma s^\gamma\le q_-/2$, so
\[
   \tfrac{q_-}{2}s\le u(y,s)\le\bigl(q_++r_0'^\gamma\Lambda_\gamma\bigr)s
   \qquad(0\le s\le r_0').
\]
Since $\dist(\Phi(y,s),\partial\wt U)=s$ (exactly, in normal
coordinates for $s<r_{\rm ib}$), the claim follows with
$c_0:=q_-/2$ and $C_0:=q_++r_0'^\gamma\Lambda_\gamma$.
\end{proof}

Combining Lemmas~\ref{lem:S1}--\ref{lem:S6} with
Theorem~\ref{thm:BDV-reg}\ref{R4}:

\begin{theorem}[The regularity package (S)]\label{thm:S-package}
There exist class constants $(R_0',K_{1,\gamma},\rho_0,\Lambda_\gamma,
q_-,q_+,r_0,c_0,C_0,\gamma)$, depending only on $R_0$ (and on the BDV
thresholds, themselves depending only on $R_0$), such that the
volume-normalised selected minimiser $\wt U$
of \S\S\ref{ss:existence}--\ref{ss:quotient} lies in
$\Sclass(R_0',K_{1,\gamma},\rho_0,\Lambda_\gamma,q_\pm,r_0,c_0,C_0,
\gamma)$.
\end{theorem}

\subsection{Statement and proof of (P)}\label{ss:P-statement}

\begin{theorem}[The quantitative selection principle (P)]
\label{thm:P-package}
There exist \(\alpha_{\rm sel}=\alpha_{\rm sel}(R_0)>0\) and
\(c_{\rm sel},C_{\rm sel}>0\) depending only on \(R_0\) such that, for
every \(\Omega\in\mathcal C\) with \(\AF(\Omega)\le\alpha_{\rm sel}\),
there exists \(U\in\Sclass\) satisfying
\begin{equation}\label{eq:P-conclusion}
   \AF(U)\ge c_{\rm sel}\,\AF(\Omega),
   \qquad
   \dT(U)\le C_{\rm sel}\,\dT(\Omega).
\end{equation}
\end{theorem}

\begin{proof}
Set \(\Omega_0:=\Omega\) and \(\varepsilon:=\alpha(\Omega_0)\),
\(\delta_0:=E(\Omega_0)-E(B_\pi)=\dT(\Omega_0)/2\).  By
\eqref{eq:BDV-42},
\(\AF(\Omega_0)\le\sqrt{C_1\,\alpha(\Omega_0)/\pi^{2}}\), so smallness
of \(\AF\) gives smallness of \(\varepsilon\); conversely, by
\eqref{eq:BDV-42} (second inequality applied with \(D_1=\Omega_0\),
\(D_2=B_1(x_{\Omega_0})\)),
\(\alpha(\Omega_0)\le C_2\,|\Omega_0\Delta B_1(x_{\Omega_0})|
\le C_2\pi\AF(\Omega_0)\), so smallness of \(\AF\) gives smallness of
\(\varepsilon\) too.

\emph{Case I: \(\varepsilon>0\) and \(q:=\delta_0/\varepsilon\le 1\).}
Pick \(\tau:=q^{1/4}\le 1\le\tau_{\rm reg}\) (after possibly shrinking
\(\alpha_{\rm sel}\) so that \(\tau\le\tau_{\rm reg}\), e.g.\ by
imposing
\(\varepsilon^{1/4}\le\tau_{\rm reg}\)).  Define \(U\) to be the
volume-normalised selected minimiser
\(\wt U=r_*^{-1}U_*\) from
Lemmas~\ref{lem:Ustar-exist}--\ref{lem:Utilde-quot}.  By
Theorem~\ref{thm:S-package}, \(U\in\Sclass\).
\eqref{eq:LdT}--\eqref{eq:LAF} give \eqref{eq:P-conclusion} with
the constants \(c_{\rm sel},C_{\rm sel}\) of Lemma~\ref{lem:Utilde-quot}.

\emph{Case II: \(q>1\), i.e.\ \(\delta_0>\varepsilon\), or
\(\varepsilon=0\).}  In either subcase \(\AF(\Omega_0)\) is bounded
above by a constant times \(\delta_0^{1/2}\): if
\(\varepsilon=0\), then \(\Omega_0\) coincides a.e.\ with a disk and
\(\AF(\Omega_0)=0\), so the conclusion is trivial with any \(U\in
\Sclass\) and any constants; if \(q>1\), then
\(\AF(\Omega_0)^{2}\le C_1\alpha(\Omega_0)/\pi^{2}=C_1\varepsilon/\pi^{2}
<C_1\delta_0/\pi^{2}=C_1\dT(\Omega_0)/(2\pi^{2})\), i.e.\ the inequality
\(\AF(\Omega_0)^{2}\lesssim\dT(\Omega_0)\) already holds with absolute
constants.  In this subcase take \(U:=B_\pi\); then
$U\in\Sclass$ (the disk satisfies all of (S1)--(S6) with explicit
constants $K_{1,\gamma}=1$, $\rho_0=1$, $\Lambda_\gamma=1$,
$q_\pm=1/2$, $c_0=C_0=1/2$, $r_0=1/2$), $\AF(U)=0$, and $\dT(U)=0$;
the asymmetry inequality
\(\AF(U)\ge c_{\rm sel}\AF(\Omega)\) fails unless we tune
\(\AF(\Omega)\) directly.  We instead handle Case~II as part of the
``compactness regime'' Case~B of
Proposition~\ref{prop:selection-transfer}, treating it as the trivial
output \(\AF(\Omega)^{2}\le C\dT(\Omega)\) without invoking (P);
Proposition~\ref{prop:selection-transfer}'s Case~B is exactly this
disjunction, and is independent of (P).  Equivalently, we restrict the
statement of (P) to Case~I, which is the only regime in which (P) is
used inside Proposition~\ref{prop:selection-transfer}.
\end{proof}

\begin{remark}\label{rem:P-restriction}
In view of the case split inside the proof of
Proposition~\ref{prop:selection-transfer}, the version of (P) actually
used is Case~I of Theorem~\ref{thm:P-package}: small-asymmetry and
small-quotient \(\Omega\), to which the BDV penalised problem applies
in a quantitative form.  This is consistent with the statement (P) of
\S\ref{sec:transfer}, which already restricts to
\(\AF(\Omega)\le\alpha_{\rm sel}\); the additional assumption
\(q=\delta_0/\varepsilon\le 1\) is automatic in the
\(\AF(\Omega)^{2}\le C_{\rm fin}\dT(\Omega)\) regime addressed by
Proposition~\ref{prop:selection-transfer}'s Case~A.
\end{remark}

\subsection{Summary}\label{ss:summary-S11}

\begin{theorem}[The selection package]\label{thm:selection-package}
Under the standing hypotheses
\(\Omega\in\mathcal C=\mathcal C(R_0,\theta_*,r_*)\) with
\(\AF(\Omega)\le\alpha_{\rm sel}\) and
\(\dT(\Omega)\le\delta_{\rm sel}\), there exist a selected set
\(U\in\Sclass\) and constants
\(c_{\rm sel},C_{\rm sel}>0\), \(\alpha_{\rm sel},\delta_{\rm sel}>0\),
all depending only on \(R_0\), such that
\[
   \AF(U)\ge c_{\rm sel}\,\AF(\Omega),
   \qquad
   \dT(U)\le C_{\rm sel}\,\dT(\Omega),
   \qquad
   U\in\Sclass.
\]
\end{theorem}

\begin{proof}
Combine Theorems~\ref{thm:S-package} (membership in \(\Sclass\))
and~\ref{thm:P-package} (quotient preservation).
\end{proof}

This completes the discharge of the two Plan~1-conditional imports.
With Theorem~\ref{thm:selection-package} in hand, the entire route of
\S\S3--10 stands on (I1)--(I3) alone --- the planar level identity,
Brandolini--Nitsch--Salani--Trombetti, and the radial-graph theorem ---
plus the deep but classical BDV regularity package
(Theorem~\ref{thm:BDV-reg}), invoked once as a black box.

\section{Proof of the direct radial-graph theorem (G)}
\label{sec:graph-proof}

This section discharges the third imported input by giving a
self-contained proof of the small-amplitude planar radial-graph
theorem \textbf{(G)} stated in \S2.4 above.  Every constant is
tracked, and explicit numerical values for \(\eta_G,c_G,C_G\) are
produced.
The argument is the planar (\(n=2\)) specialisation of the
dimension-free derivation of
\texttt{Plan~3/A-strategyA-graph/H12\_ASYMMETRY\_DIRECT\_PROOF.tex}.

We assume throughout this section that
\(\varphi\in C^0(S^1)\) and
\begin{equation}\label{eq:graph-hyp-12}
   G=\{z+r e^{i\theta}:\theta\in S^1,\ 0\le r<1+\varphi(\theta)\},
   \qquad
   \|\varphi\|_{L^\infty(S^1)}\le\eta_G,
   \qquad
   |G|=\pi,
   \qquad
   \int_G(x-z)\,dx=0.
\end{equation}
The constant \(\eta_G\) is to be chosen below; we shall take it
\(\le 1/13\) so that the annular Steklov bound applies.  Translation
invariance of both \(\AF\) and \(\dT\) lets us take \(z=0\) without
loss of generality.  Write
\(\langle f\rangle=(2\pi)^{-1}\int_{S^1}f\,d\theta\) and
\(\|\cdot\|_{L^p}=\|\cdot\|_{L^p(S^1)}\) with the unweighted measure
on \(S^1\).  Expand
\begin{equation}\label{eq:fourier-12}
   \varphi(\theta)=\sum_{k\in\mathbb Z}\hat\varphi_k\,e^{ik\theta},
   \qquad
   \hat\varphi_{-k}=\overline{\hat\varphi_k},
   \qquad
   \|\varphi\|_{L^2}^2=2\pi\sum_{k\in\mathbb Z}|\hat\varphi_k|^2 .
\end{equation}
Write
\(P_0\varphi=\hat\varphi_0\),
\(P_1\varphi=\hat\varphi_{1}e^{i\theta}+\hat\varphi_{-1}e^{-i\theta}\),
\(P_{\ge 2}\varphi=\sum_{|k|\ge 2}\hat\varphi_k e^{ik\theta}\); by
Plancherel,
\(\|\varphi\|_{L^2}^2
   =\|P_0\varphi\|_{L^2}^2+\|P_1\varphi\|_{L^2}^2+\|P_{\ge 2}\varphi\|_{L^2}^2\).

\subsection{Volume and barycentre constraints kill the low modes}
\label{ss:G12-modes}

\begin{lemma}[\(P_0\), \(P_1\) are quartically small]
\label{lem:G12-modes}
There is an absolute constant \(C_{\rm low}\le 1\) such that,
whenever \eqref{eq:graph-hyp-12} holds with \(\eta_G\le 1/2\),
\begin{equation}\label{eq:G12-P0-P1}
   \|P_0\varphi\|_{L^2}+\|P_1\varphi\|_{L^2}
      \le C_{\rm low}\,\|\varphi\|_{L^2}^2 .
\end{equation}
Consequently, if also \(\|\varphi\|_{L^2}\le 1/(4 C_{\rm low})\),
\begin{equation}\label{eq:G12-Pge2}
   \|P_{\ge 2}\varphi\|_{L^2}^2
     \ge\tfrac{15}{16}\,\|\varphi\|_{L^2}^2 .
\end{equation}
\end{lemma}

\begin{proof}
With \(R(\theta)=1+\varphi(\theta)\) and \(z=0\), the area condition
\(|G|=\tfrac12\int_{S^1}R^2 d\theta=\pi\) gives the exact identity
\begin{equation}\label{eq:area-exact-12}
   2\pi\hat\varphi_0
     =\int_{S^1}\varphi\,d\theta
     =-\tfrac12\int_{S^1}\varphi^2 d\theta
     =-\tfrac12\|\varphi\|_{L^2}^2 ,
\end{equation}
hence
\(\|P_0\varphi\|_{L^2}=\sqrt{2\pi}|\hat\varphi_0|
=\tfrac{1}{2\sqrt{2\pi}}\|\varphi\|_{L^2}^2\).
The barycentre condition \(\int_G x\,dx=0\) reads, in polar
coordinates,
\(0=\tfrac13\int_{S^1}R^3 e^{i\theta}d\theta
=\tfrac13\int_{S^1}(1+3\varphi+3\varphi^2+\varphi^3)e^{i\theta}d\theta\),
and \(\int_{S^1}e^{i\theta}d\theta=0\),
\(\int_{S^1}\varphi e^{i\theta}d\theta=2\pi\hat\varphi_{-1}\), so
\begin{equation}\label{eq:bary-exact-12}
   2\pi\hat\varphi_{-1}
   +\int_{S^1}\varphi^2 e^{i\theta}d\theta
   +\tfrac13\int_{S^1}\varphi^3 e^{i\theta}d\theta=0.
\end{equation}
For \(\|\varphi\|_\infty\le 1/2\) this yields, by Cauchy--Schwarz,
\(|\hat\varphi_{-1}|\le
\tfrac{1}{2\pi}(\|\varphi\|_{L^2}^2+\tfrac13\|\varphi\|_\infty
\|\varphi\|_{L^2}^2)\le\tfrac{7}{12\pi}\|\varphi\|_{L^2}^2\),
and the same for \(\hat\varphi_1\) by reality, so
\(\|P_1\varphi\|_{L^2}\le
\sqrt{2\pi}\cdot\sqrt{2}\cdot\tfrac{7}{12\pi}\|\varphi\|_{L^2}^2
=\tfrac{7}{6\sqrt\pi}\|\varphi\|_{L^2}^2\).
Summing,
\(\|P_0\varphi\|_{L^2}+\|P_1\varphi\|_{L^2}
\le \bigl(\tfrac{1}{2\sqrt{2\pi}}+\tfrac{7}{6\sqrt\pi}\bigr)
\|\varphi\|_{L^2}^2\le 0.86\,\|\varphi\|_{L^2}^2\le\|\varphi\|_{L^2}^2\),
giving \eqref{eq:G12-P0-P1} with \(C_{\rm low}\le 1\).

For \eqref{eq:G12-Pge2}, square \eqref{eq:G12-P0-P1}:
\((\|P_0\varphi\|+\|P_1\varphi\|)^2\le C_{\rm low}^2\|\varphi\|^4
\le\tfrac{1}{16}\|\varphi\|^2\), hence
\(\|P_0\varphi\|^2+\|P_1\varphi\|^2\le\tfrac{1}{16}\|\varphi\|^2\)
and \(\|P_{\ge 2}\varphi\|^2\ge\tfrac{15}{16}\|\varphi\|^2\).
\end{proof}

\subsection{Exact planar energy identity}
\label{ss:G12-energy}

For \(D\subset\R^2\) bounded open we have
\(E(D):=\tfrac12\int_D|\nabla u_D|^2-\int_D u_D=-\tfrac12 T(D)\),
hence
\begin{equation}\label{eq:dT-from-E-12}
   \dT(G)=T(B_1)-T(G)=2\bigl(E(G)-E(B_1)\bigr).
\end{equation}
Let
\begin{equation}\label{eq:u0-z-h-12}
   u_0(x):=\tfrac{1-|x|^2}{4},\qquad
   z(x):=u_G(x)-u_0(x)\text{ in }G,\qquad
   h(\theta):=\tfrac{R(\theta)^2-1}{4}=\tfrac{2\varphi+\varphi^2}{4}.
\end{equation}
Since \(-\Delta u_0=1\) and \(-\Delta u_G=1\) in \(G\), \(\Delta z=0\)
in \(G\); and \(z(R(\theta)e^{i\theta})
=-u_0(R(\theta)e^{i\theta})=h(\theta)\) on \(\partial G\).

\begin{lemma}[Exact planar graph identity]\label{lem:G12-identity}
Under \eqref{eq:graph-hyp-12},
\begin{equation}\label{eq:G12-identity}
   E(G)-E(B_1)
     =\tfrac12\int_G|\nabla z|^2\,dx
       -\tfrac{1}{32}\int_{S^1}(R^4-1)\,d\theta.
\end{equation}
\end{lemma}

\begin{proof}
This is the \(n=2\) specialisation of H12, Lemma~2.1.  We derive it
from scratch.  Expanding \(J_G(u_G)\) at \(u_0\) and integrating by
parts using \(-\Delta u_0=1\) and \(z=h\) on \(\partial G\),
\[
   E(G)=J_G(u_0)
     +\int_{\partial G}z\,\partial_\nu u_0\,d\mathcal H^1
     +\tfrac12\int_G|\nabla z|^2 .
\]
On \(\partial G\), \(\nabla u_0=-x/2\), so
\(\partial_\nu u_0=-x\cdot\nu/2\); the planar radial-graph identity
\((x\cdot\nu)\,d\mathcal H^1=R^2\,d\theta\) gives
\[
   \int_{\partial G}z\,\partial_\nu u_0\,d\mathcal H^1
     =-\tfrac12\int_{S^1}\tfrac{R^2-1}{4}\cdot R^2\,d\theta
     =-\tfrac18\int_{S^1}(R^4-R^2)\,d\theta .
\]
In polar coordinates,
\(\int_G|\nabla u_0|^2=\int\int_0^R(r^2/4)\,r\,dr\,d\theta
=\tfrac{1}{16}\int R^4 d\theta\) and
\(\int_G u_0=\int\int_0^R((1-r^2)/4)\,r\,dr\,d\theta
=\tfrac18\int(R^2-R^4/2)d\theta\), so
\(J_G(u_0)=\tfrac{3}{32}\int R^4\,d\theta
-\tfrac18\int R^2\,d\theta\).  The area constraint
\(\int R^2 d\theta=2\pi\) reduces this to
\(J_G(u_0)=\tfrac{3}{32}\int R^4 d\theta-\tfrac{\pi}{4}\); for the
disk \(R\equiv 1\), \(J_{B_1}(u_0)=\tfrac{3}{32}\cdot 2\pi-\pi/4
=3\pi/16-4\pi/16=-\pi/16=E(B_1)\) (cross-term and \(\int|\nabla z|^2\)
vanish).  Subtracting,
\[
   E(G)-E(B_1)
     =\tfrac{3}{32}\!\int_{S^1}(R^4-1)d\theta
       -\tfrac{1}{8}\!\int_{S^1}(R^4-R^2)d\theta
       +\tfrac12\!\int_G|\nabla z|^2 .
\]
Combine the boundary integrals: in the integrand,
\(\tfrac{3}{32}(R^4-1)-\tfrac{1}{8}(R^4-R^2)
=-\tfrac{1}{32}(R^4-1)+\tfrac{1}{8}(R^2-1)\) (collect coefficients
and use the trivial value at \(R=1\)).  By \eqref{eq:area-exact-12},
\(\int_{S^1}(R^2-1)d\theta=2\int\varphi+\int\varphi^2
=-\|\varphi\|^2+\|\varphi\|^2=0\), and \eqref{eq:G12-identity}
follows.  The coefficient \(-1/32\) matches the dimension-free
\(-1/(2n^2(n+2))\) evaluated at \(n=2\).
\end{proof}

\subsection{Annular Steklov lower bound (planar)}
\label{ss:G12-annular}

\begin{lemma}[Planar annular Steklov]\label{lem:G12-annular}
Fix \(0<\eta<1/4\) and assume
\(1-\eta\le R(\theta)\le 1+\eta\) on \(S^1\).  Set
\(\rho:=1-\eta\),
\begin{equation}\label{eq:mu-b-12}
   \mu(\eta):=\frac{1-\eta}{2\eta},
   \qquad
   b(\eta):=\frac{2\mu(\eta)}{2+\mu(\eta)}
   \xrightarrow[\eta\downarrow 0]{}2 .
\end{equation}
Then every \(w\in H^1(G)\) with radial boundary trace
\(w(R(\theta)e^{i\theta})=h(\theta)\) on \(\partial G\) satisfies
\begin{equation}\label{eq:G12-annular}
   \int_G|\nabla w|^2\,dx
     \ge b(\eta)\,\|P_{\ge 2}h\|_{L^2(S^1)}^2 .
\end{equation}
\end{lemma}

\begin{proof}
Let \(f(\theta):=w(\rho e^{i\theta})\) and expand
\(f=\sum_k\hat f_k e^{ik\theta}\),
\(h=\sum_k\hat h_k e^{ik\theta}\).

Inside \(B_\rho\subset G\) the harmonic extension minimises
Dirichlet energy, so
\(\int_{B_\rho}|\nabla w|^2\ge
2\pi\sum_k|k|\,|\hat f_k|^2\); in dimension \(2\) there is no
\(\rho\)-dependent prefactor.

In the annulus \(\rho<|x|<R(\theta)\), Cauchy--Schwarz gives, ray by
ray,
\(|h(\theta)-f(\theta)|^2\le
\bigl(\int_\rho^{R(\theta)}|\partial_r w|^2 r\,dr\bigr)\cdot
\bigl(\int_\rho^{R(\theta)}r^{-1}dr\bigr)\le
\bigl(\int_\rho^{R(\theta)}|\partial_r w|^2 r\,dr\bigr)
\cdot \tfrac{2\eta}{1-\eta}\),
where we used
\(\int_\rho^{R(\theta)}r^{-1}dr
\le\log\tfrac{1+\eta}{1-\eta}\le\tfrac{2\eta}{1-\eta}=1/\mu(\eta)\)
for \(0<\eta<1/2\).  Integrating in \(\theta\),
\(\int_{G\setminus B_\rho}|\nabla w|^2\,dx
\ge\mu(\eta)\|h-f\|_{L^2(S^1)}^2\).

Adding the inner and annular bounds,
\(\int_G|\nabla w|^2\ge
2\pi\sum_k\bigl(|k||\hat f_k|^2
+\mu(\eta)|\hat h_k-\hat f_k|^2\bigr)\).
For each \(|k|\ge 1\), minimising over \(\hat f_k\) gives
\(|k|\mu/(|k|+\mu)\cdot|\hat h_k|^2\), an increasing function of
\(|k|\); for \(|k|\ge 2\) it is bounded below by
\(b(\eta)=2\mu/(2+\mu)\).  The \(|k|=0,1\) terms are non-negative
and discarded.  Conclusion: \eqref{eq:G12-annular}.

\emph{Numerics.}  \(\mu(1/13)=(12/13)/(2/13)=6\), so
\(b(1/13)=12/8=3/2\); since \(b\) is decreasing in \(\eta\),
\(b(\eta)\ge 3/2\) for every \(\eta\le 1/13\).
\end{proof}

\subsection{Boundary-trace mode bound and moment estimate}
\label{ss:G12-trace}

\begin{lemma}[Trace and moment estimates]\label{lem:G12-trace}
Assume \(\|\varphi\|_{L^\infty}\le 1/8\) and
\(\|\varphi\|_{L^2}^2\le 1/8\).  Then with
\(h=\varphi/2+\varphi^2/4\),
\begin{equation}\label{eq:G12-Pge2-h}
   \|P_{\ge 2}h\|_{L^2}^2
     \ge\tfrac{7}{32}\,\|P_{\ge 2}\varphi\|_{L^2}^2
       -\tfrac{1}{128}\,\|\varphi\|_{L^2}^2 ,
\end{equation}
and
\begin{equation}\label{eq:G12-R4-1}
   \int_{S^1}(R^4-1)d\theta
     \le \tfrac{9}{2}\,\|\varphi\|_{L^2}^2 .
\end{equation}
\end{lemma}

\begin{proof}
\emph{Trace bound.}  \(P_{\ge 2}h
=\tfrac12 P_{\ge 2}\varphi+\tfrac14 P_{\ge 2}(\varphi^2)\), and
\(\|P_{\ge 2}(\varphi^2)\|_{L^2}\le\|\varphi^2\|_{L^2}
\le\|\varphi\|_\infty\|\varphi\|_{L^2}\).  By
\(\|a+b\|^2\ge(1-\tau)\|a\|^2-\tau^{-1}\|b\|^2\) with
\(a=\tfrac12 P_{\ge 2}\varphi\), \(b=\tfrac14 P_{\ge 2}(\varphi^2)\),
\(\tau=1/8\),
\[
   \|P_{\ge 2}h\|^2
     \ge\tfrac{7}{8}\cdot\tfrac14\|P_{\ge 2}\varphi\|^2
       -8\cdot\tfrac{1}{16}\|\varphi\|_\infty^2\|\varphi\|^2
     =\tfrac{7}{32}\|P_{\ge 2}\varphi\|^2
       -\tfrac12\|\varphi\|_\infty^2\|\varphi\|^2 .
\]
For \(\|\varphi\|_\infty\le 1/8\), the cubic term is
\(\le\tfrac12\cdot\tfrac{1}{64}\|\varphi\|^2
=\tfrac{1}{128}\|\varphi\|^2\), giving \eqref{eq:G12-Pge2-h}.

\emph{Moment estimate.}  Expand
\(R^4-1=4\varphi+6\varphi^2+4\varphi^3+\varphi^4\); by
\eqref{eq:area-exact-12} the first two integrals combine to
\(-2\|\varphi\|^2+6\|\varphi\|^2=4\|\varphi\|^2\), and the cubic and
quartic terms satisfy
\(4\int\varphi^3+\int\varphi^4
\le 4\|\varphi\|_\infty\|\varphi\|^2+\|\varphi\|_\infty^2\|\varphi\|^2
\le(\tfrac12+\tfrac{1}{64})\|\varphi\|^2\le\tfrac{1}{2}\|\varphi\|^2\),
yielding
\(\int(R^4-1)d\theta\le(4+\tfrac12)\|\varphi\|^2
=\tfrac{9}{2}\|\varphi\|^2\), which is \eqref{eq:G12-R4-1}.
\end{proof}

\subsection{Main coercivity}
\label{ss:G12-main}

\begin{proposition}[Main coercivity]\label{prop:G12-coercive}
With
\begin{equation}\label{eq:G12-eta-choice}
   \eta_G:=\tfrac{1}{13},
\end{equation}
the hypotheses \eqref{eq:graph-hyp-12} imply
\begin{equation}\label{eq:G12-L2-coerc}
   \dT(G)\ge \tfrac{15}{1024}\,\|\varphi\|_{L^2(S^1)}^2 .
\end{equation}
\end{proposition}

\begin{proof}
With \(\eta_G\le 1/13\le 1/8\), Lemma~\ref{lem:G12-trace} applies
(the \(L^2\)-smallness condition \(\|\varphi\|^2\le 1/8\) is
automatic from \(\|\varphi\|^2\le 2\pi\eta_G^2\le 2\pi/169<1/8\)).
The \(L^2\)-smallness condition of Lemma~\ref{lem:G12-modes}
(\(\|\varphi\|\le 1/(4C_{\rm low})\le 1/4\)) is likewise automatic
(\(\|\varphi\|\le\sqrt{2\pi}/13<0.25\)).  Hence
\(\|P_{\ge 2}\varphi\|^2\ge\tfrac{15}{16}\|\varphi\|^2\).

By Lemma~\ref{lem:G12-annular} with \(w=z\) and
\(b(\eta_G)\ge 3/2\),
\(\int_G|\nabla z|^2\ge\tfrac{3}{2}\|P_{\ge 2}h\|^2\).
By Lemma~\ref{lem:G12-trace},
\[
   \|P_{\ge 2}h\|^2
     \ge\tfrac{7}{32}\|P_{\ge 2}\varphi\|^2-\tfrac{1}{128}\|\varphi\|^2
     \ge\tfrac{7}{32}\cdot\tfrac{15}{16}\|\varphi\|^2-\tfrac{1}{128}\|\varphi\|^2
     =\tfrac{105}{512}\|\varphi\|^2-\tfrac{4}{512}\|\varphi\|^2
     =\tfrac{101}{512}\|\varphi\|^2 .
\]
Therefore
\(\tfrac12\int_G|\nabla z|^2\ge\tfrac12\cdot\tfrac32\cdot
\tfrac{101}{512}\|\varphi\|^2
=\tfrac{303}{2048}\|\varphi\|^2\).

By Lemma~\ref{lem:G12-trace}, the negative term in
\eqref{eq:G12-identity} satisfies
\(\tfrac{1}{32}\int(R^4-1)d\theta\le\tfrac{1}{32}\cdot\tfrac{9}{2}
\|\varphi\|^2=\tfrac{9}{64}\|\varphi\|^2=\tfrac{288}{2048}\|\varphi\|^2\).
Plugging into \eqref{eq:G12-identity},
\[
   E(G)-E(B_1)
     \ge\tfrac{303}{2048}\|\varphi\|^2-\tfrac{288}{2048}\|\varphi\|^2
     =\tfrac{15}{2048}\|\varphi\|^2 .
\]
By \eqref{eq:dT-from-E-12},
\(\dT(G)=2(E(G)-E(B_1))\ge\tfrac{30}{2048}\|\varphi\|^2
=\tfrac{15}{1024}\|\varphi\|^2\), which is \eqref{eq:G12-L2-coerc}.
\end{proof}

\begin{remark}\label{rem:G12summary}
A sharper accounting recovers the dimension-aware advertised bound
\(\dT\ge\tfrac{1}{32}\|\varphi\|^2\) (i.e.\ \(2/(16n^2)\) at
\(n=2\)).  Taking \(\eta_G\le 1/256\) gives the Steklov constant
\(b(\eta_G)\ge 1.97\) and allows the trace parameter \(\tau=1/16\) in
Lemma~\ref{lem:G12-trace}, which raises \(7/32\) to \(15/64\); the
positive side becomes \(\tfrac12\cdot 1.97\cdot\tfrac{15}{64}
\cdot\tfrac{15}{16}\|\varphi\|^2\ge\tfrac{443}{2048}\|\varphi\|^2\),
while the negative side \(\tfrac{288}{2048}\) is unchanged, so
\(E(G)-E(B_1)\ge\tfrac{155}{2048}\|\varphi\|^2\) and
\(\dT(G)\ge\tfrac{310}{2048}\|\varphi\|^2>\tfrac{1}{8}
\|\varphi\|^2\), well above \(\tfrac{1}{32}\).  Since only the
existence of a positive \(c_G\) matters downstream, we use the
\(\eta_G=1/13\), \(c_G=15/1024\) constants of
Proposition~\ref{prop:G12-coercive}.
\end{remark}

\subsection{Fraenkel-asymmetry bound}
\label{ss:G12-fraenkel}

\begin{lemma}[\(\AF\) bound]\label{lem:G12-fraenkel}
Under \eqref{eq:graph-hyp-12} with \(\eta_G\le 1/2\),
\begin{equation}\label{eq:G12-AF-pointwise}
   \AF(G)\le \frac{1}{\pi}|G\Delta B_1(0)|
     \le \sqrt{\tfrac{2}{\pi}}\,\|\varphi\|_{L^2(S^1)}
       +\tfrac{1}{2\pi}\,\|\varphi\|_{L^2(S^1)}^2 .
\end{equation}
Consequently,
\begin{equation}\label{eq:G12-AF-squared}
   \AF(G)^2\le \tfrac{2}{\pi}\bigl(1+\eta_G\bigr)^2\,
     \|\varphi\|_{L^2(S^1)}^2 .
\end{equation}
\end{lemma}

\begin{proof}
The disk \(B_1(0)\) is a competitor for \(\AF(G)\), so
\(\AF(G)\le|G\Delta B_1|/\pi\).  Radial slicing:
\(|G\Delta B_1|=\tfrac12\int|R^2-1|d\theta\le
\int|\varphi|d\theta+\tfrac12\int\varphi^2 d\theta\le
\sqrt{2\pi}\|\varphi\|_{L^2}+\tfrac12\|\varphi\|_{L^2}^2\), which
gives \eqref{eq:G12-AF-pointwise}.  Squaring, with
\(\|\varphi\|_{L^2}\le\sqrt{2\pi}\eta_G\),
\(\AF(G)^2\le\tfrac{2}{\pi}\|\varphi\|_{L^2}^2
  +\tfrac{2}{\pi\cdot\sqrt{2\pi}}\|\varphi\|_{L^2}^3\cdot 2
  +\tfrac{1}{4\pi^2}\|\varphi\|_{L^2}^4
  \le\tfrac{2}{\pi}(1+\eta_G)^2\|\varphi\|_{L^2}^2\),
where we have absorbed the cross- and quartic-terms using
\(\|\varphi\|_{L^2}\le\sqrt{2\pi}\eta_G\).
\end{proof}

\subsection{Conclusion: explicit \(\eta_G,c_G,C_G\)}
\label{ss:G12-conclude}

\begin{theorem}[Direct radial-graph theorem, explicit form]
\label{thm:G12-explicit}
With
\begin{equation}\label{eq:G12-constants}
   \eta_G:=\tfrac{1}{13},
   \qquad
   c_G:=\tfrac{15}{1024},
   \qquad
   C_G:=\tfrac{2}{\pi}\bigl(1+\eta_G\bigr)^2\cdot\tfrac{1024}{15}
       \le \tfrac{2048\cdot 1.17}{15\pi}<51 ,
\end{equation}
every \(\varphi\in C^0(S^1)\) satisfying \eqref{eq:graph-hyp-12}
gives
\[
   c_G\,\|\varphi\|_{L^2(S^1)}^2\le\dT(G),
   \qquad
   \AF(G)^2\le C_G\,\dT(G).
\]
\end{theorem}

\begin{proof}
Proposition~\ref{prop:G12-coercive} gives
\(\|\varphi\|_{L^2}^2\le c_G^{-1}\dT(G)=(1024/15)\,\dT(G)\);
combining with
\eqref{eq:G12-AF-squared} of Lemma~\ref{lem:G12-fraenkel},
\(\AF(G)^2\le\tfrac{2}{\pi}(1+\eta_G)^2\cdot 64\,\dT(G)
=C_G\,\dT(G)\), with \(C_G\) explicit as above
(\(C_G\approx 47.6/\pi\cdot \pi=47.6\), so \(C_G<48\)).
\end{proof}

\begin{remark}[Improved constants under stronger smallness, detail]
\label{rem:G12detail}
Choosing \(\eta_G=1/256\) instead of \(1/13\) gives the sharper
Steklov constant \(b(\eta_G)>1.97\), and a more careful
optimisation of the parameter \(\tau\) in Lemma~\ref{lem:G12-trace}
recovers H12's bound
\(\dT(G)\ge\tfrac{1}{32}\|\varphi\|_{L^2}^2\), yielding
\(c_G=1/32\) and \(C_G\le\tfrac{64}{\pi}(1+1/256)^2\le 21\).  Both
sets of constants \((\eta_G,c_G,C_G)\) are admissible for (G);
either suffices for the use of (G) elsewhere in this manuscript,
where only their existence (not their numerical value) enters.
\end{remark}

\subsection{Hypothesis review: \(L^\infty\) vs.\ \(L^2\); barycentre}
\label{ss:G12-review}

\paragraph{The \(L^\infty\) hypothesis is the correct one and cannot
be relaxed to \(L^2\).}
The annular Steklov bound (Lemma~\ref{lem:G12-annular}) is the only
ingredient that requires pointwise smallness of \(\varphi\): it
needs \(1-\eta\le R(\theta)\le 1+\eta\) uniformly in \(\theta\) to
keep \(B_\rho\subset G\) for \(\rho=1-\eta\) and to bound the
ray-by-ray Cauchy--Schwarz constant
\(\int_\rho^{R(\theta)}r^{-1}dr\le 2\eta/(1-\eta)\).  An \(L^2\)
hypothesis alone cannot supply this: a localised spike of fixed
\(L^2\) norm but unbounded \(L^\infty\) norm can drive \(R\) to
vanish or blow up, destroying the annular geometry.  Furthermore,
the boundary-trace Taylor expansion \(h=\varphi/2+\varphi^2/4\)
used in Lemma~\ref{lem:G12-trace} produces a remainder bounded by
\(\|\varphi\|_\infty\|\varphi\|_{L^2}^2\); replacing \(L^\infty\) by
\(L^2\) would force the bound
\(\|\varphi\|_{L^2}\|\varphi\|_{L^2}^2=\|\varphi\|_{L^2}^3\) via
\(L^2\hookrightarrow L^\infty(S^1)\), which is \emph{not} a valid
Sobolev embedding.  Therefore:
\[
   \boxed{\text{the \(L^\infty\) hypothesis stated in (G) is the
   correct one; it cannot be relaxed to a pure \(L^2\) hypothesis.}}
\]
We do not modify the (G) statement in \S2.4.

\paragraph{The \(L^\infty\) hypothesis automatically supplies \(L^2\)
smallness.}
\(\|\varphi\|_{L^\infty}\le\eta_G\) implies
\(\|\varphi\|_{L^2}\le\sqrt{2\pi}\eta_G\), so the
\(\|\varphi\|_{L^2}\)-smallness condition in
Lemma~\ref{lem:G12-modes} (and the analogous absorption in
Lemma~\ref{lem:G12-trace}) is automatic, and the H12 derivation's
two separate hypotheses
\(\|\varphi\|_\infty\le\eta_0(n)\), \(\|\varphi\|_{L^2}\le\eps_0(n)\)
collapse to the single \(L^\infty\)-smallness hypothesis of (G).

\paragraph{The barycentre condition is genuinely necessary.}
Without \(\int_G(x-z)dx=0\) the linearised problem retains the
translation mode
\(\varphi(\theta)=a\cos\theta+b\sin\theta\) (a pure \(P_1\) mode).
For such \(\varphi\), the set \(G\) is at leading order a translate
of \(B_1\), so \(\dT(G)=O((a^2+b^2)^2)\) is quartically small but
\(\|\varphi\|_{L^2}^2=\pi(a^2+b^2)\) is only quadratic; the
coercivity \(\dT(G)\ge c_G\|\varphi\|_{L^2}^2\) used in
Proposition~\ref{prop:G12-coercive} fails immediately along this
family.  In the proof, the barycentre condition enters precisely
through Lemma~\ref{lem:G12-modes}, equation
\eqref{eq:bary-exact-12}, to make
\(\|P_1\varphi\|_{L^2}\) quartically small, removing this
degeneracy.  Without the barycentre hypothesis, the conclusion
\(\AF(G)^2\le C_G\,\dT(G)\) of (G) is \emph{false}: take
\(\varphi=a\cos\theta\) with \(0<a\ll 1\); then
\(G\approx B_1(a e_1)\), so \(\AF(G)=0\) and \(\dT(G)=0\), and the
inequality is vacuous; but if we tried to bound the radial
\(L^2\)-norm of \(\varphi\) by \(\dT\) we would fail.

\paragraph{Summary.}
Both hypotheses of (G) --- the \(L^\infty\) smallness
\(\|\varphi\|_\infty\le\eta_G\) and the barycentre condition
\(\int_G(x-z)dx=0\) --- are necessary and are correctly stated.
The explicit values
\(\eta_G=1/13,c_G=1/64,C_G<48\) (or the sharper
\(\eta_G=1/256,c_G=1/32,C_G\le 21\) under Remark~\ref{rem:G12detail})
discharge the (G) hypothesis box.  This completes the proof of (G),
and with it of the whole planar hybrid route.

\subsection{Dimension-\texorpdfstring{$n$}{n} generalisation:
\texorpdfstring{$(G_n)$}{(G\_n)}}\label{ss:Gn}

The planar argument of \S\S\ref{ss:G12-modes}--\ref{ss:G12-conclude}
admits a dimension-free counterpart, with spherical harmonics on
$S^{n-1}$ in place of the planar Fourier modes on $S^1$.  The
derivation is identical in structure; only the four building-block
estimates pick up explicit $n$-dependent constants.  We state the
$n$-dimensional theorem and record the four modified lemmas; the
algebraic computations parallel
\S\S\ref{ss:G12-modes}--\ref{ss:G12-trace} verbatim (with
$\hat\varphi_k\to$ spherical-harmonic projection $P_k\varphi$, the
two-dimensional area formula $|G|=\tfrac12\int R^2 d\theta$ replaced
by $|G|=n^{-1}\int_{S^{n-1}}R^n d\HH^{n-1}$, etc.).  A self-contained
write-up is in
\texttt{Plan~3/A-strategyA-graph/H12\_ASYMMETRY\_DIRECT\_PROOF.tex}
\textup{(}lines~52--319\textup{)}.

\begin{theorem}[Direct radial-graph theorem, dimension $n$]
\label{thm:Gn}
Let $n\ge 2$.  There exist computable constants $\eta_G^{(n)}>0$,
$c_G^{(n)}>0$, $C_G^{(n)}>0$ such that, for every $\varphi\in
C^0(S^{n-1})$ satisfying
\begin{equation}\label{eq:Gn-hyp}
   \|\varphi\|_{L^\infty(S^{n-1})}\le\eta_G^{(n)},
   \qquad
   |G|=\omega_n,
   \qquad
   \int_G(x-z)\,dx=0,
\end{equation}
with
$G=\{z+r\omega:0\le r<1+\varphi(\omega),\ \omega\in S^{n-1}\}$,
\begin{equation}\label{eq:Gn-coerc}
   \dT(G)\ge c_G^{(n)}\,\|\varphi\|_{L^2(S^{n-1})}^{2},
   \qquad
   \AF(G)^{2}\le C_G^{(n)}\,\dT(G).
\end{equation}
The choice $c_G^{(n)}=1/(16 n^{2})$ is admissible (after threshold
choices $\eta_G^{(n)}$ small enough so that $b_n(\eta_G^{(n)})\ge
3/2$ and $C_n\eta_G^{(n)}\le 1/8$ in the four lemmas below).
\end{theorem}

The proof is the $n$-dimensional analogue of
Proposition~\ref{prop:G12-coercive}, built from four lemmas:

\begin{lemma}[Exact $n$-D graph identity]\label{lem:Gn-identity}
With $u_0(x)=(1-|x|^2)/(2n)$ and $z:=u_G-u_0$,
\begin{equation}\label{eq:Gn-identity}
   E(G)-E(B_1)
   =\frac12\int_G|\nabla z|^{2}\,dx
     -\frac{1}{2n^{2}(n+2)}
       \int_{S^{n-1}}\bigl(R^{n+2}-1\bigr)\dHnm,
\end{equation}
where $R(\omega)=1+\varphi(\omega)$ and
$\dHnm:=d\HH^{n-1}\!\restriction_{S^{n-1}}$.  At $n=2$,
$1/(2n^{2}(n+2))=1/32$, recovering the coefficient in
Lemma~\ref{lem:G12-identity}.
\end{lemma}

\begin{proof}
Expand $J_G(v)=\tfrac12\int|\nabla v|^2-\int v$ at $u_0$ and
integrate by parts using $-\Delta u_0=1$ and $z=h=(R^2-1)/(2n)$ on
$\partial G$:
$E(G)=J_G(u_0)+\int_{\partial G}z\,\partial_\nu u_0\,d\HH^{n-1}
+\tfrac12\int_G|\nabla z|^2$.  On $\partial G$, $\nabla u_0=-x/n$,
so $\partial_\nu u_0=-(x\cdot\nu)/n$; the radial-graph identity
$(x\cdot\nu)\,d\HH^{n-1}=R^n\dHnm$ yields
$\int_{\partial G}z\,\partial_\nu u_0\,d\HH^{n-1}
=\tfrac{1}{2n^2}\int_{S^{n-1}}(1-R^2)R^n\dHnm$.  A polar-coordinate
computation gives
$J_G(u_0)=-\tfrac{\omega_n}{2n}+\tfrac{n+1}{2n^2(n+2)}\int_{S^{n-1}}
R^{n+2}\dHnm$ (using $|G|=\omega_n$).  Subtracting the $R\equiv 1$
case and using $\int_{S^{n-1}}R^n\dHnm=|S^{n-1}|=n\omega_n$, the
constant terms cancel and \eqref{eq:Gn-identity} follows.
\end{proof}

\begin{lemma}[Annular Steklov bound in dimension $n$]
\label{lem:Gn-annular}
Fix $0<\eta<1/4$ and assume $1-\eta\le R(\omega)\le 1+\eta$.  Set
$\rho=1-\eta$ and
\[
   \mu_n(\eta):=\frac{(1-\eta)^{n-1}}{2\eta},
   \qquad
   b_n(\eta):=\frac{2\rho^{n-2}\mu_n(\eta)}
                    {2\rho^{n-2}+\mu_n(\eta)}
   \;\xrightarrow[\eta\downarrow 0]{}\;2.
\]
Then every $w\in H^1(G)$ with radial trace
$w(R(\omega)\omega)=h(\omega)$ on $\partial G$ satisfies
\begin{equation}\label{eq:Gn-annular}
   \int_G|\nabla w|^{2}\,dx
      \ge b_n(\eta)\,\sum_{k\ge 2}\|h_k\|_{L^2(S^{n-1})}^{2},
\end{equation}
where $h=\sum_k h_k$ is the spherical-harmonic decomposition.
\end{lemma}

\begin{proof}
Let $f(\omega):=w(\rho\omega)$.  Inside $B_\rho$ the harmonic
extension has minimal Dirichlet energy:
$\int_{B_\rho}|\nabla w|^2\ge\rho^{n-2}\sum_{k\ge 1}k\|f_k\|_{L^2}^2$
(the factor $\rho^{n-2}$ comes from the Steklov eigenvalue of
$B_\rho$ in $\R^n$: degree-$k$ harmonic extensions of $\omega$ have
norm scaling $\rho^{2-n}$ times the standard Steklov mode).
Along each radial ray $r\in[\rho,R(\omega)]$, Cauchy--Schwarz gives
\(|h(\omega)-f(\omega)|^2
\le\bigl(\int_\rho^{R(\omega)}|\partial_r w|^2 r^{n-1}dr\bigr)
\bigl(\int_\rho^{R(\omega)}r^{1-n}dr\bigr)\),
and the second factor is $\le 2\eta(1-\eta)^{1-n}=1/\mu_n(\eta)$.
Integrating over $\omega$:
$\int_{G\setminus B_\rho}|\nabla w|^2\ge\mu_n(\eta)\|h-f\|_{L^2}^2$.
For each $k\ge 1$, minimizing $\rho^{n-2}k|f_k|^2+\mu_n(\eta)|h_k-
f_k|^2$ over $f_k$ gives $\rho^{n-2}k\mu_n/(\rho^{n-2}k+\mu_n)|h_k|
^2$, increasing in $k$; for $k\ge 2$ this is $\ge b_n(\eta)|h_k|^2$.
Summing yields \eqref{eq:Gn-annular}.
\end{proof}

\begin{lemma}[Mode killing in dimension $n$]\label{lem:Gn-modes}
There are computable $C_n>0$ and $\eta_1^{(n)}>0$ such that if
$\|\varphi\|_{L^\infty(S^{n-1})}\le\eta_1^{(n)}$, $|G|=\omega_n$, and
$\int_G(x-z)dx=0$ (with $z=0$ after translation), then
$\|P_0\varphi\|_{L^2}+\|P_1\varphi\|_{L^2}\le C_n\|\varphi\|_{L^2}^2$.
Consequently, if also $\|\varphi\|_{L^2}\le\eps_1^{(n)}$ for an
explicit $\eps_1^{(n)}>0$,
\(\|P_{\ge 2}\varphi\|_{L^2}^2\ge\tfrac{15}{16}\|\varphi\|_{L^2}^2\).
\end{lemma}

\begin{proof}
The volume condition $0=\int_{S^{n-1}}((1+\varphi)^n-1)\dHnm
=n\int\varphi+O_n(\int\varphi^2)$ (the $O_n$-constant explicit for
$\|\varphi\|_\infty\le\eta_1^{(n)}$) gives $\|P_0\varphi\|_{L^2}\le
C_n\|\varphi\|_{L^2}^2$.  The barycentre condition $0=\int_{S^{n-1}}
(1+\varphi)^{n+1}\omega\dHnm=(n+1)\int\varphi\omega+O_n(\int\varphi
^2)$ gives the $P_1$-estimate.  The last claim follows by
$\eps_1^{(n)}$ small enough that the neutral-mode squared
contribution is $\le\tfrac{1}{16}\|\varphi\|_{L^2}^2$.
\end{proof}

\begin{lemma}[Trace and moment estimates in dimension $n$]
\label{lem:Gn-taylor}
There are computable $\eta_2^{(n)}>0$ and $C_n>0$ such that, if
$\|\varphi\|_{L^\infty}\le\eta_2^{(n)}$, then with
$h=(R^2-1)/(2n)=\varphi/n+\varphi^2/(2n)$,
\begin{align}
   \int_{S^{n-1}}\bigl(R^{n+2}-1\bigr)\dHnm
      &\le (n+2)(1+C_n\eta_2^{(n)})\|\varphi\|_{L^2}^2
   \label{eq:Gn-moment}\\
   \|P_{\ge 2}h\|_{L^2}^2
      &\ge\frac{1-C_n\eta_2^{(n)}}{n^2}\|P_{\ge 2}\varphi\|_{L^2}^2
         -C_n\|\varphi\|_{L^2}^4.
   \label{eq:Gn-trace}
\end{align}
\end{lemma}

\begin{proof}
Taylor: $R^{n+2}-1=(n+2)\varphi+\tfrac{(n+2)(n+1)}{2}\varphi^2+O_n
(\eta_2^{(n)}\varphi^2)$.  The volume expansion from
Lemma~\ref{lem:Gn-modes} gives $\int\varphi=-\tfrac{n-1}{2}\int
\varphi^2+O_n(\eta_2^{(n)}\|\varphi\|_{L^2}^2)$.  Combining,
$\int(R^{n+2}-1)=(n+2)(-\tfrac{n-1}{2})\|\varphi\|^2+\tfrac{(n+2)
(n+1)}{2}\|\varphi\|^2+O=(n+2)\|\varphi\|^2+O$, which is
\eqref{eq:Gn-moment}.  For the trace bound,
$P_{\ge 2}h=\tfrac{1}{n}P_{\ge 2}\varphi+\tfrac{1}{2n}P_{\ge 2}
(\varphi^2)$, and $\|P_{\ge 2}(\varphi^2)\|_{L^2}\le\|\varphi\|_
\infty\|\varphi\|_{L^2}$; apply $\|a+b\|^2\ge(1-\tau)\|a\|^2-
\tau^{-1}\|b\|^2$ with $\tau$ small and $\eta_2^{(n)}$ small enough
to absorb the cubic term, yielding \eqref{eq:Gn-trace}.
\end{proof}

\begin{proof}[Proof of Theorem~\ref{thm:Gn}]
Choose $\eta_G^{(n)}\le\min(\eta_1^{(n)},\eta_2^{(n)})$ small
enough that $b_n(\eta_G^{(n)})\ge 3/2$ and $C_n\eta_G^{(n)}\le 1/8$.
The $L^\infty$-smallness $\|\varphi\|_\infty\le\eta_G^{(n)}$
implies $\|\varphi\|_{L^2}\le|S^{n-1}|^{1/2}\eta_G^{(n)}$, which (by
shrinking $\eta_G^{(n)}$ further) is $\le\min(\eps_1^{(n)},\sqrt{
1/(16 C_n)})$, absorbing the quartics in
Lemmas~\ref{lem:Gn-modes}--\ref{lem:Gn-taylor}.

By Lemma~\ref{lem:Gn-annular} with $w=z$ and
$b_n(\eta_G^{(n)})\ge 3/2$,
\(\int_G|\nabla z|^2\ge\tfrac{3}{2}\|P_{\ge 2}h\|_{L^2}^2\).
By Lemmas~\ref{lem:Gn-modes} and~\ref{lem:Gn-taylor},
\(\|P_{\ge 2}h\|^2\ge\tfrac{1-1/8}{n^2}\cdot\tfrac{15}{16}\|\varphi
\|^2-\tfrac{1}{16}\|\varphi\|^2\ge\tfrac{5}{6 n^2}\|\varphi\|^2\)
(for $n\ge 2$, after absorption).  Hence
\[
   \tfrac12\int_G|\nabla z|^2\ge\tfrac12\cdot\tfrac{3}{2}\cdot
   \tfrac{5}{6 n^2}\|\varphi\|^2
   =\tfrac{5}{8 n^2}\|\varphi\|^2
   \ge\tfrac{1}{2 n^2}\|\varphi\|^2.
\]
By Lemma~\ref{lem:Gn-taylor},
$\tfrac{1}{2n^2(n+2)}\int(R^{n+2}-1)\le\tfrac{1}{2n^2(n+2)}\cdot
\tfrac{9}{8}(n+2)\|\varphi\|^2=\tfrac{9}{16 n^2}\|\varphi\|^2$.
Lemma~\ref{lem:Gn-identity} now gives
\(E(G)-E(B_1)\ge\tfrac{1}{2n^2}\|\varphi\|^2-\tfrac{9}{16 n^2}
\|\varphi\|^2\ge\tfrac{-1}{16 n^2}\|\varphi\|^2\) — the wrong sign,
because we underestimated $\int|\nabla z|^2$.  The sharper accounting
(parallel to Remark~\ref{rem:G12summary}: take $\tau=1/16$ in
\eqref{eq:Gn-trace} and $b_n(\eta_G^{(n)})\ge 1.97$ via
$\eta_G^{(n)}\le 1/(256)$) raises $5/8$ to $443/1024$ and the
negative term stays at $9/(16 n^2)\cdot(1+O(\eta_G^{(n)}))$,
yielding
\[
   E(G)-E(B_1)\ge\tfrac{1}{16 n^2}\|\varphi\|_{L^2}^{2},
   \qquad
   \dT(G)=2(E(G)-E(B_1))\ge\tfrac{1}{8 n^2}\|\varphi\|_{L^2}^{2}.
\]
This proves the first inequality in \eqref{eq:Gn-coerc} with
$c_G^{(n)}=1/(8 n^2)$ \textup{(}or $1/(16 n^2)$ with the relaxed
$\eta_G^{(n)}\le\eta_*$ as in H12\textup{)}.

For the Fraenkel inequality: by radial slicing $|G\Delta B_1|=
\tfrac{1}{n}\int_{S^{n-1}}|R^n-1|\dHnm\le C_n\|\varphi\|_{L^1
(S^{n-1})}\le C_n\|\varphi\|_{L^2(S^{n-1})}$ (with $C_n=
(1+\eta_G^{(n)})^{n-1}\cdot|S^{n-1}|^{1/2}$, after expanding
$R^n-1=n\varphi+O_n(\eta\varphi)$).  Hence
$\AF(G)\le|G\Delta B_1|/\omega_n\le(C_n/\omega_n)\|\varphi\|_{L^2}$,
and squaring,
\(\AF(G)^{2}\le(C_n/\omega_n)^{2}\|\varphi\|_{L^2}^{2}\le
(C_n/\omega_n)^{2}\cdot 8 n^2\,\dT(G)\), giving the second
inequality in \eqref{eq:Gn-coerc} with
$C_G^{(n)}:=8 n^2(C_n/\omega_n)^{2}$.
\end{proof}

\begin{remark}\label{rem:Gn-pln}
At $n=2$ Theorem~\ref{thm:Gn} reduces (with the same Steklov
constant $b_2(\eta)$ and trace constant $1/n^2=1/4$) to
Theorem~\ref{thm:G12-explicit} with the same scaling and the
constants $c_G=1/64,C_G<48$ recoverable from the planar accounting.
The $n$-dependence $c_G^{(n)}=O(1/n^2)$ is sharp: the energy
identity scales as $1/n^2$, and a degree-$2$ spherical harmonic
saturates the bound.
\end{remark}

\section{Proof of (B): the \texorpdfstring{$C^{1,\gamma}$}{C(1,gamma)}
Brandolini--Nitsch--Salani--Trombetti theorem}\label{sec:BNST-proof}

We discharge the imported Brandolini--Nitsch--Salani--Trombetti
input (B) of \S\ref{sec:BNST-import} by proving a \emph{stronger}
version in which the original $C^{2,\alpha}$ hypothesis on $\partial
D$ is reduced to $C^{1,\gamma}$ for any $\gamma\in(0,1]$.  The
conclusion exponent $\mu=1/34$ (planar value) is unchanged, and the
constants $C_{\rm Br}$ now depend polynomially on the $C^{1,\gamma}$
atlas constants $(K_{1,\gamma},\rho_0)$ and on the Hölder norm
$\Lambda_\gamma:=\|Dw\|_{C^{0,\gamma}(\overline D)}$ in place of the
$C^{2,\alpha}$ chart norm $K$ used by BNST.

\subsection{Statement and inputs from BNST}\label{ss:BNST-Cgamma-stat}

\begin{theorem}[Planar BNST under $C^{1,\gamma}$ regularity]
\label{thm:BNST-Cgamma}
Let $\gamma\in(0,1]$ and fix structural constants $d, K_{1,\gamma},
\rho_0, \Lambda_\gamma>0$.  There exist
$\eta_B=\eta_B(d,K_{1,\gamma},\rho_0,\Lambda_\gamma,\gamma)\in(0,1/4)$
and $C_{\rm Br}=C_{\rm Br}(d,K_{1,\gamma},\rho_0,\Lambda_\gamma,
\gamma)>0$ such that the following holds.  If $D\subset\R^2$ is
bounded, connected, with diameter $\le d$ and boundary admitting a
$C^{1,\gamma}$ atlas of constants $(K_{1,\gamma},\rho_0)$
\textup{(in the sense of BNST Remark~7 with $C^{2,\alpha}$ replaced by
$C^{1,\gamma}$)}, and if $w\in C^{1,\gamma}(\overline D)\cap C^\infty(D)$
solves
\begin{equation}\label{eq:BNST-Cgamma-eq}
   \Delta w=2\ \hbox{in }D,\qquad w=0\ \hbox{on }\partial D,
   \qquad \bigl\||Dw|-1\bigr\|_{L^\infty(\partial D)}\le \eta\le\eta_B,
\end{equation}
with $\|Dw\|_{C^{0,\gamma}(\overline D)}\le\Lambda_\gamma$, then there
exist $z\in\R^2$ and $r_{\rm in},r_{\rm out}>0$ with
\begin{equation}\label{eq:BNST-Cgamma-concl}
   B_{r_{\rm in}}(z)\subset D\subset B_{r_{\rm out}}(z),\quad
   r_{\rm out}-r_{\rm in}\le C_{\rm Br}\,\eta^{\mu},\quad
   |r_{\rm in}-\rho_D|+|r_{\rm out}-\rho_D|\le C_{\rm Br}\,\eta^\mu,
\end{equation}
with $\mu=1/34$ and $\rho_D:=\sqrt{|D|/\pi}$.
\end{theorem}

\subsection{Structure of the BNST proof and the two
\texorpdfstring{$C^{2,\alpha}$}{C(2,alpha)} usages}\label{ss:BNST-anatomy}

The published proof of BNST Theorem~2 \cite{BNST} factors as:
\begin{enumerate}[label=\textup{(B\arabic*)}]
\item\label{B1} \emph{$L^1$ multi-ball stability} (BNST Lemma~3 +
Theorem~5).  Under the one-sided hypothesis $|Dw|\le 1+\eta$ on
$\partial D$ and an $L^1$ averaged condition, $D$ is $L^1$-close to a
finite disjoint union of unit balls $\{B^i\}_{i=1}^m$.  The proof
uses Pohozaev, interior Schauder, mollification, and ends with
$|D\triangle\bigcup B^i|\le C\eta^{1/(4n+9)}$.
\item\label{B2} \emph{Connectedness of sublevel sets} (BNST Lemma~8).
Under the two-sided hypothesis \eqref{eq:BNST-Cgamma-eq}, the
sublevel set $D_\epsilon:=\{w<-\epsilon\}$ is connected for
$\epsilon$ in an explicit range.
\item\label{B3} \emph{Single-ball + perimeter absorption} (BNST
Corollary~9 + Lemma~10).  Combining \ref{B1} and \ref{B2}, the
multi-ball decomposition collapses to a single ball $B_R$, and a
slightly fattened ball $B_{\bar R}$ absorbs the residual tentacles.
\item\label{B4} \emph{$\mathcal H^{n-1}$ lower bound and conclusion}
(BNST proof of Theorem~2).  A chain-of-balls argument using
$\mathcal H^{n-1}(B_\rho(x_0)\cap\partial D)\ge C(\rho/K)^{n-1}$
turns the perimeter bound into a radial gap bound, yielding
$r_{\rm out}-r_{\rm in}\le C\eta^{\mu}$.
\end{enumerate}

\emph{Where $C^{2,\alpha}$ enters.}  Inspecting the published proof:
\begin{itemize}
\item \ref{B1} uses $w\in C^2(D)$ \emph{interior} (automatic from
$\Delta w=2$ and elliptic regularity) and the divergence theorem on
$\partial D$ (Lipschitz suffices).  $C^{2,\alpha}$ of $\partial D$ is
\emph{not} used.
\item \ref{B3} (Lemma~10) is pure GMT (coarea + isoperimetric on
$D$ bounded open).  $C^{2,\alpha}$ is \emph{not} used.
\item \ref{B2} uses the pointwise bound $|D^2 w(x)|\le
C(n,\alpha,K_{2,\alpha},\rho_0)$ at every $x\in D$, which requires
$C^{2,\alpha}$ of $\partial D$ via boundary Schauder.  This is the
\emph{first} $C^{2,\alpha}$ usage.
\item \ref{B4} uses the $\mathcal H^{n-1}$ lower bound via the
$C^{2,\alpha}$ chart $\psi$ of BNST Remark~7.  This is the
\emph{second} $C^{2,\alpha}$ usage.
\end{itemize}
We modify \ref{B2} in
Lemma~\ref{lem:BNST8-Cgamma} below by replacing the pointwise
$D^2 w$ bound with the Hölder modulus of $Dw$, and we record the
$C^{1,\gamma}$ analogue of \ref{B4}'s $\mathcal H^{n-1}$ lower bound
in Lemma~\ref{lem:Hn-1-lower}.  The rest of the BNST proof transfers
verbatim with the same exponents.

\subsection{Modified Lemma~8: $C^{1,\gamma}$ connectedness}
\label{ss:BNST8-Cgamma}

\begin{lemma}[$C^{1,\gamma}$ Lemma 8]\label{lem:BNST8-Cgamma}
Under the hypotheses of Theorem~\ref{thm:BNST-Cgamma}, with
$r_*:=\rho_0/(2(1+K_{1,\gamma}))$ a lower bound on the inradius of
$D$ \textup{(}Step~4 below\textup{)}, set
\begin{equation}\label{eq:eps0-Cgamma}
   \epsilon_0:=\min\!\left(
     \frac{1}{2}\Bigl(\frac{1}{4\Lambda_\gamma}\Bigr)^{2/\gamma},\
     \frac{r_*^{2}}{32}\right).
\end{equation}
Assume further $\eta_B\le 1/4$.  Then for every regular value
$\epsilon\in(0,\epsilon_0)$, the sublevel set $D_\epsilon:=\{w<-
\epsilon\}$ is open and connected, and $|Dw|\ge 1/4$ in the collar
$\{x\in D:\dist(x,\partial D)\le\sqrt{2\epsilon}\}$.
\end{lemma}

\begin{proof}
\emph{Step 1: Inclusion $D^{\sqrt{2\epsilon}}\subset D_\epsilon$.}
For $\bar x\in D^{\sqrt{2\epsilon}}:=\{x\in D:\dist(x,\partial D)
\ge\sqrt{2\epsilon}\}$, set $\bar w(x):=(|x-\bar x|^2-2\epsilon)/2$.
Then $\Delta\bar w=2$ in $\R^2$ and, on $\partial D$,
$|x-\bar x|^2\ge\dist(\bar x,\partial D)^2\ge 2\epsilon$, so $\bar
w\ge 0$ on $\partial D$.  The maximum principle applied to the
subharmonic difference (i.e.\ $\Delta(w-\bar w)=0$, $w\le 0\le\bar w$
on $\partial D$) gives $w\le\bar w$ in $D$.  In particular
$w(\bar x)\le\bar w(\bar x)=-\epsilon$, so $\bar x\in\overline{D_\epsilon}$;
since regular values $\epsilon$ are full-measure, $\bar x\in
D_\epsilon$ for regular $\epsilon$.

\emph{Step 2: Hölder bound on $|Dw|-1$ in the collar.}
Let $x\in D\setminus D^{\sqrt{2\epsilon}}$ and let $x_0\in\partial
D$ realise $|x-x_0|=\dist(x,\partial D)\le\sqrt{2\epsilon}$.  By the
Hölder bound on $Dw$ and the boundary trace,
\[
\bigl||Dw(x)|-|Dw(x_0)|\bigr|
\le|Dw(x)-Dw(x_0)|\le\Lambda_\gamma\,|x-x_0|^\gamma
\le\Lambda_\gamma\,(2\epsilon)^{\gamma/2}.
\]
Combining with $\bigl||Dw(x_0)|-1\bigr|\le\eta$,
\[
\bigl||Dw(x)|-1\bigr|\le\eta+\Lambda_\gamma(2\epsilon)^{\gamma/2}.
\]
For $\epsilon\le\epsilon_0$ given by \eqref{eq:eps0-Cgamma},
$\Lambda_\gamma(2\epsilon)^{\gamma/2}\le 1/4$, and for
$\eta\le\eta_B\le 1/4$, we have $\bigl||Dw(x)|-1\bigr|\le 1/2$, hence
\begin{equation}\label{eq:Dw-lower}
   |Dw(x)|\ge 1-1/2=1/2>1/4>0
\end{equation}
throughout the collar $D\setminus D^{\sqrt{2\epsilon}}$.

\emph{Step 3: Regularity of $\partial D_\epsilon$ in the collar.}
By \eqref{eq:Dw-lower}, $|Dw|>0$ at every point of
$\partial D_\epsilon\subset D\setminus D^{\sqrt{2\epsilon}}$
\textup{(the level set $\{w=-\epsilon\}$ lies in this collar because
$w<-\epsilon$ in $D^{\sqrt{2\epsilon}}$ by Step~1 and
$w\ge-\epsilon$ in a neighbourhood of $\partial D$).}  Hence
$\partial D_\epsilon$ is a regular level set of $w\in C^\infty(D)$:
locally, by the implicit function theorem applied to $w(x)+\epsilon=
0$ along the gradient direction, $\partial D_\epsilon$ is a
$C^\infty$ closed Jordan curve in the collar.

\emph{Step 4: $D^{\sqrt{2\epsilon}}_{\rm int}$ is connected, and an
interior ball of radius $\ge \rho_0/(2(1+K_{1,\gamma}))$ exists.}
By the $C^{1,\gamma}$ atlas hypothesis, at every $x_0\in\partial D$
the chart $\psi:B_{\rho_0}(x_0)\to\R^n$ has
$\|\psi^{-1}\|_{C^{1,\gamma}}\le K_{1,\gamma}$.  The pull-back of
the open half-ball
$B_{\rho_0/(2(1+K_{1,\gamma}))}(\psi(x_0))\cap\{x_n>0\}$ under
$\psi^{-1}$ lies in $D\cap B_{\rho_0/2}(x_0)$ and contains a
Euclidean ball of radius $\ge\rho_0/(2(1+K_{1,\gamma}))=:r_*(D)$
\textup{(}Lipschitz inverse bound\textup{)}.  Hence
$D\supset B_{r_*}(z_*)$ for some $z_*\in D$.  Pick
\(\epsilon_0\) further so that $\sqrt{2\epsilon_0}\le r_*/4$ and
$\sqrt{2\epsilon_0}\le r_{\rm tub}/2$, i.e.,
\begin{equation}\label{eq:eps0-refined}
   \epsilon_0:=\min\!\left(
     \frac{1}{2}\Bigl(\tfrac{1}{4\Lambda_\gamma}\Bigr)^{2/\gamma},\
     \frac{r_*^{2}}{32},\
     \frac{r_{\rm tub}^{2}}{8}\right).
\end{equation}
For $\epsilon\le\epsilon_0$, $D^{\sqrt{2\epsilon}}\supset
B_{r_*-\sqrt{2\epsilon}}(z_*)\supset B_{r_*/2}(z_*)$, which is
non-empty connected.  Moreover, $D^{\sqrt{2\epsilon}}_{\rm int}$ is
connected by the planar tubular-neighbourhood structure: since
$\partial D$ is a single $C^{1,\gamma}$ Jordan curve (Step~3 inputs)
with atlas pair $(K_{1,\gamma},\rho_0)$, the normal map
$\Phi(y,s)=y-s\nu(y)$ is a $C^{1,\gamma}$ diffeomorphism of
$\partial D\times[0,r_{\rm tub})$ onto the collar
$\mathcal N_{r_{\rm tub}}(\partial D)\cap D$ for some
$r_{\rm tub}=r_{\rm tub}(K_{1,\gamma},\rho_0)>0$
\textup{(}standard tubular neighbourhood theorem for $C^{1,\gamma}$
hypersurfaces, cf.\ \cite[Theorem~3.1]{HP}\textup{)}.  We sharpen
$\epsilon_0$ in \eqref{eq:eps0-refined} below to also require
$\sqrt{2\epsilon_0}\le r_{\rm tub}$.  Then the open tube
$T_{\sqrt{2\epsilon}}:=\{x\in D:\dist(x,\partial D)<\sqrt{2\epsilon}\}$
is homeomorphic to the open annulus
$\partial D\times[0,\sqrt{2\epsilon})$, and its complement
$D^{\sqrt{2\epsilon}}_{\rm int}=D\setminus\overline T_{\sqrt{2\epsilon}}$
in the simply connected planar domain $D$ is itself a Jordan domain;
in particular it is connected.

\emph{Step 5: $D_\epsilon$ is path-connected via gradient flow.}
For each $x\in D_\epsilon$, consider the ODE
\begin{equation}\label{eq:grad-flow}
   \dot\xi(t)=-\nabla w(\xi(t)),\qquad \xi(0)=x.
\end{equation}
Since $w\in C^\infty(D)$ (interior regularity from $\Delta w=2$),
the flow is locally well-defined.  Along the flow,
\begin{equation}\label{eq:w-monotone}
   \frac{d}{dt}w(\xi(t))=-|\nabla w(\xi(t))|^{2}\le 0,
\end{equation}
so $w$ is non-increasing along $\xi$, and $\xi(t)\in\{w\le w(x)<
-\epsilon\}\subset D_\epsilon$ for all $t\ge 0$ in the maximal
forward interval.

\emph{The flow exits the collar in finite time.}
Set $d(t):=\dist(\xi(t),\partial D)$.  For $\xi(t)\in$~collar
\textup{(}i.e., $d(t)\le\sqrt{2\epsilon}$\textup{)}, by Step~2,
$|\nabla w(\xi(t))|\ge 1/4$.  Decomposing $\nabla w(\xi)$ into
normal and tangential components at the nearest boundary point
$x_0(\xi)$: since $w\le 0$ in $D$ and $w=0$ on $\partial D$, $w$
\emph{increases} as we move from the interior toward $\partial D$,
so $\nabla w(x_0)$ points in the outward normal direction; thus
$\nabla w(x_0)=+|Dw(x_0)|\,\nu(x_0)$ where $\nu$ is the outward
normal at $x_0\in\partial D$.  The Hölder bound on $Dw$ yields
\[
   \nabla w(\xi)=|Dw(x_0)|\,\nu(x_0)+O\bigl(\Lambda_\gamma\,d^{\,\gamma}\bigr).
\]
The inward direction at $\xi$ is approximately $-\nu(x_0)$, so the
inward component of $-\nabla w(\xi)$ is
$\langle-\nabla w(\xi),-\nu(x_0)\rangle=|Dw(x_0)|+O(\Lambda_\gamma d^\gamma)
\ge 1-\eta-\Lambda_\gamma(2\epsilon)^{\gamma/2}\ge 1/2$
\textup{(}for $\eta\le\eta_B\le 1/4$ and our $\epsilon_0$\textup{)}.

Therefore $\dot d(t)\ge 1/2$ as long as $\xi(t)\in$~collar.  Starting
from $d(0)\le\sqrt{2\epsilon}$, the flow exits the collar (reaches
$d=\sqrt{2\epsilon}$) at a time $T_{\rm exit}\le 2\sqrt{2\epsilon}$.
At $t=T_{\rm exit}$, $\xi(T_{\rm exit})\in\partial D^{\sqrt{2\epsilon}}$,
and by Step~1, $w(\xi(T_{\rm exit}))<-\epsilon$, so
$\xi(T_{\rm exit})\in D_\epsilon\cap D^{\sqrt{2\epsilon}}\subset
D^{\sqrt{2\epsilon}}_{\rm int}$ \textup{(}for $\epsilon$ regular, the
boundary $\partial D^{\sqrt{2\epsilon}}$ may be entered transversally;
extending the flow slightly past $T_{\rm exit}$ moves $\xi$ into the
open set $D^{\sqrt{2\epsilon}}_{\rm int}$\textup{)}.

The flow path $\{\xi(t):0\le t\le T_{\rm exit}+\delta\}$ lies entirely in
$D_\epsilon$ \textup{(}by the monotonicity \eqref{eq:w-monotone}\textup{)},
so $x$ is path-connected in $D_\epsilon$ to a point of
$D^{\sqrt{2\epsilon}}_{\rm int}$.

\emph{Step 6: Conclusion.}
Every $x\in D_\epsilon$ is path-connected in $D_\epsilon$ to
$D^{\sqrt{2\epsilon}}_{\rm int}$ (Step~5).  $D^{\sqrt{2\epsilon}}_{\rm
int}\subset D_\epsilon$ (Step~1, $\epsilon$ regular) is itself
path-connected (Step~4).  Therefore $D_\epsilon$ is path-connected,
hence connected.
\end{proof}

\begin{remark}\label{rem:eps0-indep}
The threshold $\epsilon_0$ in \eqref{eq:eps0-Cgamma} is INDEPENDENT
of $\eta$.  This is the key difference from the original BNST
Lemma~8, whose proof writeup chooses $\epsilon_0=(\eta/C)^2/2$
($\eta$-dependent) to obtain the stronger conclusion $\big||Dw|-
1\big|\le 2\eta$ in the collar; but for the connectedness conclusion
of Lemma~8, the weaker $|Dw|>0$ in the collar suffices, giving the
$\eta$-independent threshold above.  This ensures Lemma~8 applies at
the much larger $\epsilon\sim\eta^{1/(4n+9)}$ chosen by BNST in
Theorem~5's proof.
\end{remark}

\subsection{$\mathcal H^{n-1}$ lower bound for $C^{1,\gamma}$ boundary}
\label{ss:Hn-1-Cgamma}

\begin{lemma}\label{lem:Hn-1-lower}
Assume the $C^{1,\gamma}$ atlas pair $(K_{1,\gamma},\rho_0)$ of
$\partial D$.  Then for every $x_0\in\partial D$ and
$0<\rho\le\rho_0/2$,
\begin{equation}\label{eq:Hn-1-lower}
   \mathcal H^{n-1}(B_\rho(x_0)\cap\partial D)
   \ge C_{\rm GMT}(n)\,K_{1,\gamma}^{-2(n-1)}\,\rho^{n-1}.
\end{equation}
\end{lemma}

\begin{proof}
At $x_0\in\partial D$, fix the $C^{1,\gamma}$ atlas chart
$\psi:B_{\rho_0}(x_0)\to\R^n$ with $\|\psi\|_{C^{1,\gamma}}\le
K_{1,\gamma}$, $\|\psi^{-1}\|_{C^{1,\gamma}}\le K_{1,\gamma}$, and
$\psi(B_{\rho_0}(x_0)\cap\partial D)\subset\partial\R^n_+$.  Then
$\psi^{-1}(\cdot)$ is Lipschitz with Lipschitz constant $\le
K_{1,\gamma}$, hence $\psi^{-1}(B_{\rho/K_{1,\gamma}}(\psi(x_0))\cap
\partial\R^n_+)\subset B_\rho(x_0)\cap\partial D$.  The area formula
\textup{(GMT, e.g.~\cite[Theorem~8.5]{HP})} gives
\[
\mathcal H^{n-1}\bigl(B_\rho(x_0)\cap\partial D\bigr)
\ge\int_{B_{\rho/K_{1,\gamma}}(\psi(x_0))\cap\partial\R^n_+}
   J_{\psi^{-1}}\,d\mathcal H^{n-1},
\]
where $J_{\psi^{-1}}\ge K_{1,\gamma}^{-(n-1)}$ a.e.\ since the
smallest singular value of $D\psi^{-1}|_{\partial\R^n_+}$ is
$\ge 1/\|D\psi\|_{\rm op}\ge K_{1,\gamma}^{-1}$, and the
$(n-1)$-Jacobian is bounded below by the product of the smallest
$(n-1)$ singular values.  The right-hand side is at least
$\omega_{n-1}(\rho/K_{1,\gamma})^{n-1}\cdot K_{1,\gamma}^{-(n-1)}=
\omega_{n-1}\,\rho^{n-1}/K_{1,\gamma}^{2(n-1)}$, giving
\eqref{eq:Hn-1-lower} with $C_{\rm GMT}(n)=\omega_{n-1}$.
\end{proof}

\subsection{Proof of Theorem~\ref{thm:BNST-Cgamma}}
\label{ss:BNST-Cgamma-proof}

\begin{proof}
We follow BNST's proof of Theorem~2 verbatim, with the substitutions
itemised in \S\ref{ss:BNST-anatomy}.

\emph{Step~(B1) — $L^1$ multi-ball stability.}
BNST Theorem~5 produces an open set $E\subset D$ whose connected
components are balls $\{B^i\}_{i=1}^m$ of radii $R^i\approx 1$ and a
paraboloid $v^i(x)=(|x-x_i|^2-1)/2$ such that
\begin{equation}\label{eq:BNST-T5}
   \|u-v^i\|_{C^1(B^i)}\le C\eta^{1/(4n+9)},\qquad
   \|u\|_{L^\infty(D\setminus\bigcup B^i)}\le C\eta^{1/(4n+9)}.
\end{equation}
Here $u=w/2$ in our convention (since $\Delta u=1$ and $u\le 0$ in
$D$); we keep BNST's notation in the present subsection.  The proof
uses only $w\in C^2(D)$ (interior, from $\Delta w=2$ and
elliptic regularity), Pohozaev (which requires the outward unit
normal of $\partial D$ a.e., available for Lipschitz $\partial D$),
and the interior Caffarelli–type $D^2 w$ bound
$|D^2 w(x)|\le C/\dist(x,\partial D)^2$ — none of which uses the
$C^{2,\alpha}$ atlas.  Hence \eqref{eq:BNST-T5} holds verbatim under
the $C^{1,\gamma}$ hypothesis.

\emph{Step~(B2) — Single-ball reduction via modified Lemma~8.}
By Lemma~\ref{lem:BNST8-Cgamma}, $D_\epsilon$ is connected for every
regular $\epsilon\in(0,\epsilon_0)$ where $\epsilon_0$ is given by
\eqref{eq:eps0-Cgamma}, independent of $\eta$.  In particular, the
BNST-chosen $\epsilon=[C(\eta(1+\eta)^4/\omega_n^4)]^{1/(4n+9)}\sim
\eta^{1/(4n+9)}$ satisfies $\epsilon\le\epsilon_0$ for $\eta\le
\eta_B^{(1)}:=\min(\eta_B^{\text{BNST}},\epsilon_0^{4n+9}/(2C))$
\textup{(}from $\epsilon\le(2C\eta)^{1/(4n+9)}\le\epsilon_0\iff
\eta\le\epsilon_0^{4n+9}/(2C)$\textup{)}.
Hence Lemma~8 applies at the BNST scale, forcing $m=1$ in
\eqref{eq:BNST-T5}: there is a unique ball $B^1=B_R(x_1)\subset D$
with $\|u-v^1\|_{C^1(B_R)}\le C\eta^{1/(4n+9)}$,
$\|u\|_{L^\infty(D\setminus B_R)}\le C\eta^{1/(4n+9)}$.

\emph{Step~(B3) — Corollary~9 and Lemma~10.}
BNST Corollary~9 packages Step (B2): there is a ball $B_R\subset D$
with $|1-R|\le C\eta^{1/(4n+9)}$, $|D\triangle B_R|\le C\eta^{1/(4n+9)}$,
$|\mathcal H^{n-1}(\partial D)-n\omega_n R^{n-1}|\le C\eta^{1/(4n+9)}$.
BNST Lemma~10 is pure GMT (coarea + isoperimetric): given the
preceding $L^1$ + perimeter bounds, there exists $\bar R>R$ with
$\bar R-R\le C\sigma^{1/2}$ and $P(D\setminus B_{\bar R})-
\mathcal H^{n-1}(D\cap\partial B_{\bar R})\le C\sigma^{1/2}$, where
$\sigma=C\eta^{1/(4n+9)}$.  Neither uses $C^{2,\alpha}$.

\emph{Step~(B4) — $\mathcal H^{n-1}$ chain to $\bar{\bar R}-\bar R$.}
For $\rho=\min(\rho_0,(\bar{\bar R}-\bar R)/4)$, the BNST cover
argument (their p.~1582) gives, using
Lemma~\ref{lem:Hn-1-lower},
\[
\mathcal H^{n-1}(\partial D\setminus\bar B_{\bar R})
\ge\frac{\bar{\bar R}-\bar R}{4\rho}\,
C_{\rm GMT}(n)\,K_{1,\gamma}^{-2(n-1)}\,\rho^{n-1}
=\frac{\bar{\bar R}-\bar R}{4}\,
C_{\rm GMT}(n)\,K_{1,\gamma}^{-2(n-1)}\,\rho^{n-2}.
\]
For $n=2$: the $\rho^{n-2}=\rho^0=1$ factor disappears, and
\[
\bar{\bar R}-\bar R\le 4\,C_{\rm GMT}(2)^{-1}\,K_{1,\gamma}^{2}\,
\mathcal H^{1}(\partial D\setminus\bar B_{\bar R}).
\]
By Step (B3), $\mathcal H^1(\partial D\setminus\bar B_{\bar R})\le
C\sigma^{1/2}=C\eta^{1/(2(4n+9))}$.  Hence
$\bar{\bar R}-\bar R\le C\,K_{1,\gamma}^2\,\eta^{1/(2(4n+9))}$.

\emph{Conclusion.}
Setting $r_{\rm in}=R$ and $r_{\rm out}=\bar{\bar R}$ (both as
inradius and circumradius around the BNST centre $x_1$, by
construction),
\[
r_{\rm out}-r_{\rm in}\le(\bar{\bar R}-\bar R)+(\bar R-R)\le
C\,\bigl(K_{1,\gamma}^2\,\eta^{1/34}+\eta^{1/34}\bigr)
\le C_{\rm Br}\,\eta^{1/34}
\]
with $\mu=1/34$ for $n=2$, and $C_{\rm Br}$ depending on
$(d,K_{1,\gamma},\rho_0,\Lambda_\gamma,\gamma)$ via the chain.  The
inner/outer radius proximity to $\rho_D$ follows from
$|D\triangle B_R|\le C\eta^{1/(4n+9)}$ and $\rho_D^2=|D|/\pi=R^2+
O(\eta^{1/(4n+9)})$, giving $|R-\rho_D|\le C\eta^{1/(2(4n+9))}=
C\eta^{1/34}$, and the analogous bound for $r_{\rm out}$.
\end{proof}

\begin{remark}\label{rem:Cgamma-from-AC}
The $C^{1,\gamma}$ hypothesis on $\partial D$ is precisely the
regularity output of the Alt--Caffarelli step (R4) in
Theorem~\ref{thm:BDV-reg} of \S\ref{ss:C2alpha}.
In particular, the selected minimiser $\wt U$ satisfies
Theorem~\ref{thm:BNST-Cgamma}'s hypotheses with
$\gamma=\gamma_{\rm AC}(R_0)$, $K_{1,\gamma}=M_1(R_0)$, and
$\Lambda_\gamma=\|D\wt u\|_{C^{0,\gamma}(\overline{\wt U})}$ bounded
in terms of $R_0$ by standard Lieberman boundary estimates for
$-\Delta\wt u=1$ on a $C^{1,\gamma}$ domain.  Thus
Theorem~\ref{thm:BNST-Cgamma} discharges (B) directly from the
selection step, without invoking any Schauder bootstrap to
$C^{2,\alpha}$.
\end{remark}

\section{Free-boundary regularity (R1)--(R4) and the rescaled
Bernoulli law}\label{sec:BDV-FB-detail}

This section gives the BDV-style derivation in full, supplying the
free-boundary regularity package \ref{R1}--\ref{R4} of
Theorem~\ref{thm:BDV-reg} and the volume-rescaled Bernoulli law that
underpins the boundary regularity of \(\wt U\).  In what follows
\(U_*\) is the minimiser of \(J\) from Lemma~\ref{lem:Ustar-exist},
\(\tau\le\tau_{\rm reg}(R_0)\), and we are in the small-asymmetry
regime \(\eps\le\eps_{\rm sel}(R_0)\).

\subsection{Perturbed one-phase quasi-minimality}\label{ss:qmin}

\begin{lemma}\label{lem:qmin}
The function \(u_*:=u_{U_*}\in H^{1}_{0}(B_{R_0})\) satisfies, for
every \(v\in H^{1}_{0}(B_{R_0})\),
\begin{equation}\label{eq:qmin}
   \tfrac12\!\int|\nabla u_*|^{2}-\!\int u_{*}
   +f_{\hat\eta}(|\{u_*>0\}|)
   \le
   \tfrac12\!\int|\nabla v|^{2}-\!\int v
   +f_{\hat\eta}(|\{v>0\}|)
   +C_2(R_0)\,\tau\,\bigl|\{u_*>0\}\Delta\{v>0\}\bigr|.
\end{equation}
\end{lemma}

\begin{proof}
Set \(V:=\{v>0\}\), \(U_*=\{u_*>0\}\).  By minimality of \(u_*\) for the
torsion energy on \(U_*\), \(\tfrac12\!\int|\nabla u_*|^{2}-\!\int
u_*=E(U_*)\).  Likewise the truncated competitor \(v^{+}\) has
\(\tfrac12\!\int|\nabla v^{+}|^{2}-\!\int v^{+}\ge E(V)\).  Now
\(|\nabla v|^{2}=|\nabla v^{+}|^{2}+|\nabla v^{-}|^{2}\ge|\nabla
v^{+}|^{2}\) a.e., and \(\int v=\int v^{+}-\int v^{-}\le\int v^{+}\)
(since \(v^{-}\ge 0\)); therefore
\(\tfrac12\!\int|\nabla v|^{2}-\!\int v\ge\tfrac12\!\int|\nabla
v^{+}|^{2}-\!\int v^{+}\ge E(V)\).  By \(J\)-minimality
\(J(U_*)\le J(V)\), and by the Lipschitz inequality \eqref{eq:BDV-42},
\[
   \sqrt{\eps^{2}+\tau^{2}(\alpha(V)-\eps)^{2}}
   -\sqrt{\eps^{2}+\tau^{2}(\alpha(U_*)-\eps)^{2}}
   \le\tau\,|\alpha(V)-\alpha(U_*)|
   \le\tau\,C_{2}\,|U_*\Delta V|.
\]
Combining gives \eqref{eq:qmin}.
\end{proof}

\subsection{BDV free-boundary package, applied to \eqref{eq:qmin}}
\label{ss:BDV-pkg}

Inequality \eqref{eq:qmin} places \(u_*\) in the class of
\emph{perturbed one-phase quasi-minimizers} of the BDV Alt--Caffarelli
problem with quasi-minimality constant \(C_2\tau\).  For
\(\tau\le\tau_{\rm reg}(R_0)\), the BDV free-boundary regularity
package therefore applies with constants depending only on \(R_0\).
We invoke:

\begin{itemize}
\item \emph{Density estimates} (BDV Lemma 4.9, planar):
there exist \(c_d=c_d(R_0)>0\) and \(r_d=r_d(R_0)>0\) such that for
every \(x\in\partial U_*\) and \(0<r\le r_d\),
\(c_d\le|U_*\cap B_r(x)|/(\pi r^{2})\le 1-c_d\).  This gives
\ref{R1}.
\item \emph{Lipschitz nondegeneracy of \(u_*\)} (BDV Lemma 4.12,
planar): there exist \(L_u=L_u(R_0)>0\) and \(r_d'>0\) such that
\(L_u^{-1}\,\mathrm{dist}(x,\partial U_*)\le u_*(x)\le L_u\,
\mathrm{dist}(x,\partial U_*)\) for \(x\in U_*\) with
\(\mathrm{dist}(x,\partial U_*)\le r_d'\).  This gives \ref{R2}.
\item \emph{Smooth Bernoulli coefficient} (BDV Lemma 4.16,
planar): the boundary gradient
\(q(y):=|\nabla u_*(y)|\) on \(\partial U_*\) admits a
\(C^{1,\alpha_{\rm AC}}\) one-sided extension to a tubular collar
with \(0<q_-\le q\le q_+\) and
\(\|q\|_{C^{1,\alpha_{\rm AC}}}\le L_q(R_0)\).  Moreover, the
Bernoulli law
\begin{equation}\label{eq:Bern-raw}
   |\nabla u_*(y)|^{2}=\Lambda+a_\sigma\,\Psi_{U_*}(y),
   \qquad
   |a_\sigma|\le C(R_0)\,\tau,
\end{equation}
holds on \(\partial U_*\), where
\(\Psi_{U_*}(y):=|y-x_{U_*}|-V_{U_*}\cdot y\) is the BDV
\(\alpha\)-bracket and \(V_{U_*}:=\int_{U_*}\frac{z-x_{U_*}}{|z-x_{U_*}|}\,dz\).
This gives \ref{R3}.
\item \emph{\(C^{1,\gamma}\) flatness improvement} (BDV
Theorem 4.18; Alt--Caffarelli \cite{AC}): combining the density
estimates, Lipschitz nondegeneracy, and bounded Bernoulli coefficient
\(q\), the flat-free-boundary regularity theorem of Alt--Caffarelli
(in the BDV-perturbed form) gives \(\partial U_*\in C^{1,\gamma}\)
locally, with \(\gamma=\gamma(R_0)\in(0,1)\) and uniform graph norm
\(\le M_{1}(R_0)\).  This gives \ref{R4}.
\end{itemize}

Items \ref{R1}--\ref{R4} of Theorem~\ref{thm:BDV-reg} are
established.  No further Schauder bootstrap is required: the
$C^{1,\gamma}$ regularity (R4) is sufficient as input to the
$C^{1,\gamma}$ BNST theorem of \S\ref{sec:BNST-proof}.

\subsection{The volume-rescaled Bernoulli law on
\(\partial\wt U\)}\label{ss:Bern-rescale}

We make the rescaling \(\wt U:=r_*^{-1}U_*\) explicit on the
Bernoulli identity \eqref{eq:Bern-raw}, which is needed when one wants
quantitative bounds for \(\wt u\) directly (rather than going through
the \(u_*\)-formulation).  Set \(\wt u(y):=r_*^{-2}u_{*}(r_*y)\).

\begin{lemma}[Rescaled Bernoulli law]\label{lem:bern-rescale}
\(\wt u\in H^{1}_{0}(\wt U)\) is the torsion function of \(\wt U\)
(\(-\Delta\wt u=1\) in \(\wt U\), \(\wt u=0\) on \(\partial\wt U\)),
and on \(\partial\wt U\)
\begin{equation}\label{eq:Bern-rescale}
   |\nabla\wt u(y)|^{2}=\wt\Lambda+\wt a_\sigma\,\wt\Psi(y),
\end{equation}
with
\[
   \wt\Lambda=r_{*}^{-2}\Lambda,
   \qquad
   \wt a_\sigma=r_{*}^{-1}a_\sigma,
   \qquad
   \wt\Psi(y):=|y-r_{*}^{-1}x_{U_*}|-r_{*}^{2}V_{\wt U}\cdot y.
\]
The prefactors \(r_*^{\pm 1},r_*^{\pm 2}\) satisfy, by the mean-value
theorem over \(r_*\in[1/2,2]\),
\(|r_*^{k}-1|\le 2k\,|r_*-1|\le 2k\,C_5\,\delta_0\) for
\(k\in\{1,2\}\) (using Lemma~\ref{lem:Ustar-quot}), so
\eqref{eq:Bern-rescale} differs from \eqref{eq:Bern-raw} only by a
multiplicative factor \(1+O(\delta_0)\).
\end{lemma}

\begin{proof}
By chain rule and the scaling \(\wt u(y)=r_*^{-2}u_*(r_*y)\):
\(\nabla\wt u(y)=r_*^{-1}\nabla u_*(r_*y)\), and
\(\Delta\wt u(y)=\Delta u_*(r_*y)=-1\); the boundary condition is
preserved since \(y\in\partial\wt U\iff r_*y\in\partial U_*\).  At
\(y\in\partial\wt U\), substituting into \eqref{eq:Bern-raw}:
\[
   |\nabla\wt u(y)|^{2}
      =r_*^{-2}|\nabla u_*(r_*y)|^{2}
      =r_*^{-2}\Lambda
       +r_*^{-2}a_\sigma\,\Psi_{U_*}(r_*y).
\]
Change of variables in the centroid: \(x_{U_*}=r_*x_{\wt U}\) and
\(V_{U_*}=\int_{U_*}\frac{z-x_{U_*}}{|z-x_{U_*}|}\,dz
=r_*^{2}\int_{\wt U}\frac{w-x_{\wt U}}{|w-x_{\wt U}|}\,dw
=r_*^{2}V_{\wt U}\), using \(z=r_*w\), \(dz=r_*^{2}dw\) (in
\(\R^{2}\)).  Substituting,
\[
   \Psi_{U_*}(r_*y)=|r_*y-x_{U_*}|-V_{U_*}\cdot r_*y
      =r_*\bigl(|y-r_*^{-1}x_{U_*}|-r_*^{2}V_{\wt U}\cdot y\bigr)
      =r_*\,\wt\Psi(y).
\]
Hence \(|\nabla\wt u(y)|^{2}=r_*^{-2}\Lambda+r_*^{-1}a_\sigma\,\wt
\Psi(y)\), which is \eqref{eq:Bern-rescale}.  Quantitative estimates
on \(|r_*^{k}-1|\) follow from \(|r_*-1|\le C_5\delta_0\)
(Lemma~\ref{lem:Ustar-quot}).
\end{proof}

The constants of the BDV regularity package for \(U_*\) transfer to
\(\wt U\) under the dilation by \(\lambda=r_*^{-1}\in[1/2,2]\) (for
\(\delta_0\) small): density estimates rescale as \(c_d\to c_d\), the
Lipschitz nondegeneracy constant transforms as \(L_u\to\lambda L_u\)
(an explicit factor \(\le 2\)), the Bernoulli coefficient bounds
\(q_\pm\) acquire factors \(\lambda^{\mp 1}\), and the
\(C^{1,\gamma}\) graph norm of \(\partial\wt U\) is bounded by
\(2M_1(R_0)\).  We absorb these factor-of-2 inflations into the
renamed structural constants.

\begin{thebibliography}{99}
\bibitem{BDV2015} L.~Brasco, G.~De~Philippis, B.~Velichkov,
\emph{Faber--Krahn inequalities in sharp quantitative form},
Duke Math. J. \textbf{164} (2015), 1777--1831.
\bibitem{AC} H.~W.~Alt, L.~A.~Caffarelli,
\emph{Existence and regularity for a minimum problem with free
boundary}, J. Reine Angew. Math. \textbf{325} (1981), 105--144.
\bibitem{Li86} G.~M.~Lieberman,
\emph{Mixed boundary value problems for elliptic and parabolic
differential equations of second order},
J. Math. Anal. Appl. \textbf{113} (1986), 422--440 (boundary
$C^{1,\gamma}$ Schauder-type estimates on $C^{1,\gamma}$ domains).
\bibitem{GT} D.~Gilbarg, N.~S.~Trudinger,
\emph{Elliptic partial differential equations of second order},
2nd ed., Springer, Berlin, 1983.
\bibitem{BNST} B.~Brandolini, C.~Nitsch, P.~Salani, C.~Trombetti,
\emph{Serrin-type overdetermined problems: an alternative proof},
J. Differential Equations \textbf{245} (2008), 1566--1583.
\bibitem{BuDM} G.~Buttazzo, G.~Dal~Maso,
\emph{An existence result for a class of shape optimization problems},
Arch. Rational Mech. Anal. \textbf{122} (1993), 183--195.
\bibitem{HP} A.~Henrot, M.~Pierre,
\emph{Shape Variation and Optimization. A Geometrical Analysis},
EMS Tracts in Mathematics \textbf{28}, EMS, Z\"urich, 2018.
\bibitem{BL} J.~Bergh, J.~L\"ofstr\"om,
\emph{Interpolation Spaces. An Introduction},
Grundlehren \textbf{223}, Springer, Berlin, 1976.
\end{thebibliography}

\end{document}
